\documentclass[10pt,leqno,twoside]{article}

% packages
\usepackage[alphabetic]{amsrefs}
\usepackage{physics}
% margin spacing
\usepackage[top=1in, bottom=1in, left=1in, right=1in]{geometry}
\usepackage{amsfonts, amsmath, amssymb, amsthm}
\usepackage{fancyhdr}
\usepackage{graphicx}
\graphicspath{{./images/}}
\usepackage{enumitem}
\usepackage{quiver}
% indexing
\usepackage{imakeidx}
    \AtBeginDocument{\renewcommand\indexname{Symbols and definitions}}
    \makeindex
\indexsetup{level=\section*,toclevel=section,noclearpage,firstpagestyle=frontmatter}
\usepackage{hyperref}
% \hypersetup{colorlinks=true,linkcolor=blue}
\usepackage[noabbrev, capitalise]{cleveref}
\crefformat{equation}{equation~#2#1#3}
\crefformat{lemma}{\textrm{Lemma}~#2#1#3}

% theorems
\theoremstyle{plain}
\newtheorem{lem}{Lemma}
\newtheorem{lemma}[lem]{Lemma}
\newtheorem{thm}[lem]{Theorem}
\newtheorem{theorem}[lem]{Theorem}
\newtheorem{prop}[lem]{Proposition}
\newtheorem{proposition}[lem]{Proposition}
\newtheorem{cor}[lem]{Corollary}
\newtheorem{corollary}[lem]{Corollary}
\newtheorem{conj}[lem]{Conjecture}
\newtheorem{fact}[lem]{Fact}
\newtheorem{form}[lem]{Formula}

\theoremstyle{definition}
\newtheorem{defn}[lem]{Definition}
\newtheorem{definition/}[lem]{Definition}
\newenvironment{definition}
  {\renewcommand{\qedsymbol}{\textdagger}%
   \pushQED{\qed}\begin{definition/}}
  {\popQED\end{definition/}}
\newtheorem{example}[lem]{Example}
\newtheorem{remark}[lem]{Remark}
\newtheorem{exercise}[lem]{Exercise}
\newtheorem{notation}[lem]{Notation}

\numberwithin{equation}{section}
\numberwithin{lem}{section}

% % colors
% \usepackage{xcolor}
% \definecolor{p}{HTML}{FFDDDD}
% \definecolor{g}{HTML}{D9FFDF}
% \definecolor{y}{HTML}{FFFFCF}
% \definecolor{b}{HTML}{D9FFFF}
% \definecolor{o}{HTML}{FADECB}
% %\definecolor{}{HTML}{}

% % \highlight[<color>]{<stuff>}
% \newcommand{\highlight}[2][p]{\mathchoice%
%   {\colorbox{#1}{$\displaystyle#2$}}%
%   {\colorbox{#1}{$\textstyle#2$}}%
%   {\colorbox{#1}{$\scriptstyle#2$}}%
%   {\colorbox{#1}{$\scriptscriptstyle#2$}}}%

% header/footer formatting
\fancypagestyle{frontmatter}{
    \fancyhf{}
    \pagestyle{fancy}
    \fancyhead[LE,RO]{\thepage}
}
\fancypagestyle{body}{
    \pagestyle{fancy}
    \fancyhead[LO,RE]{\nouppercase{\rightmark}}
    \fancyhead[LE,RO]{\thepage}
}
\renewcommand{\headrulewidth}{1pt}

% paragraph indentation/spacing
\setlength{\parindent}{0pt}
\setlength{\parskip}{6pt}
\renewcommand{\baselinestretch}{1.00}

% extra commands defined here (rarely used)
\newcommand{\br}[1]{\left(#1\right)}
\newcommand{\sbr}[1]{\left[#1\right]}
\newcommand{\cbr}[1]{\left\{#1\right\}}

% bracket notation for inner product
\usepackage{mathtools}

\DeclarePairedDelimiterX{\abr}[1]{\langle}{\rangle}{#1}

% new commands
\newcommand{\boldindex}[1]{\textbf{\textit{#1\index{#1}}}} % adds to index
\newcommand{\textib}[1]{\textbf{\textit{#1}}}
\DeclareMathOperator{\Mat}{M}
\DeclareMathOperator{\GL}{GL}
\DeclareMathOperator{\SL}{SL}
\newcommand{\smod}[1]{\;(\bmod\; #1)}
\newcommand{\smallabcd}{\big(\!\begin{smallmatrix}
    a & b \\ c & d
\end{smallmatrix}\!\big)}
\newcommand{\abcd}{\begin{pmatrix}
    a & b \\ c & d
\end{pmatrix}}
\newcommand{\slz}{\SL_2(\mathbb{Z})}
\newcommand{\glr}{\GL_2(\mathbb{R})}
\newcommand{\glrp}{\GL_2^+(\mathbb{R})}
\newcommand{\glqp}{\GL_2^+(\mathbb{Q})}
\newcommand{\old}{\text{old}}
\newcommand{\new}{\text{new}}
\DeclareMathOperator{\sgn}{sgn}
% \DeclareMathOperator{\Span}{span}
% \DeclareMathOperator{\nullity}{nullity}
% \DeclareMathOperator\Aut{Aut}
% \DeclareMathOperator\Inn{Inn}
% \DeclareMathOperator{\Orb}{Orb}
\DeclareMathOperator{\lcm}{lcm}
% \DeclareMathOperator{\Hol}{Hol}
% \DeclareMathOperator{\Jac}{Jac}
% \DeclareMathOperator{\rad}{rad}
% \DeclareMathOperator{\Tor}{Tor}
% \DeclareMathOperator{\End}{End}
% \DeclareMathOperator{\Gal}{Gal}
% \DeclareMathOperator{\Nat}{Nat}
% \DeclareMathOperator{\Frac}{Frac}
\DeclareMathOperator{\id}{id}
% \DeclareMathOperator{\im}{im}
% \DeclareMathOperator{\Hom}{Hom}
% \DeclareMathOperator{\Ext}{Ext}
% \DeclareMathOperator{\aug}{aug}

% indices
\setcounter{page}{1}
\setcounter{section}{-1}

% array column and row separation
\arraycolsep = 2pt
\renewcommand{\arraystretch}{.8}

% temporaries for editing
\usepackage{color}
\newcommand{\sai}[1]{\textcolor{red}{#1}}

\begin{document}
\begin{titlepage}
    \begin{center}
        \vspace*{4em}
        {\Large\textbf{Strong multiplicity one for classical modular forms}}

        \vspace{12em}
        \includegraphics[scale=0.14]{uf.png}

        \vspace{14em}
        Saisudharshan Sivakumar

        Undergraduate Honors Thesis -- Spring 2024\\
        Department of Mathematics, University of Florida\\
        v1.0
    \end{center}
\end{titlepage}
\begin{center}
    \textbf{Abstract}
\end{center}
Classical modular forms are analytic functions of the upper half plane invariant under the action of congruence subgroups of $\slz$, up to a factor of automorphy. In this setting, the Hecke operators on the space of modular forms are normal operators on the subspace of cusp forms, the modular forms that vanish at cusps. Cusp forms satisfy the strong multiplicity one property, which roughly says that normalized eigenfunctions of the Hecke operators are uniquely determined by their eigenvalues. Strong multiplicity one, and other multiplicity one results, are often presented in the language of automorphic forms and representation theory. We introduce the theory of classical modular forms and provide a self-contained proof of the strong multiplicity one property for classical modular forms, following the proof in Modular Forms by Toshitsune Miyake. This approach avoids using the language of automorphic forms and representation theory by studying $L$-functions associated to modular forms, and their Euler products. An application of the strong multiplicity one property yields a basis for the space of cusp forms, which is nice since the space of modular forms decomposes into the direct sum of the space of cusp forms and the space of Eisenstein series.\newpage
\pagestyle{frontmatter}
\section*{Acknowledgements}
I would like to thank my advisor, Dr.\@ Jeremy Booher, for his guidance as I learned about modular forms and the difficulties of writing mathematics. Prior to his advising, I had not even had the foggiest idea of how to write a thesis, much less of how select a suitable topic. I would like to thank him specifically for reading and commenting on drafts of my thesis. He also helped me greatly in the application process for mathematics graduate school, by writing letters of recommendation, giving advice on my personal statements, and encouraging me to apply to good programs that would fit my interests. He also taught me algebra in a kind manner that reaffirmed my interest in the subject, and provided friendship in difficult times. 

I also wish to thank Dr.\@ Sergei Shabanov for his guidance and support throughout my entire time as an undergraduate. I began as a freshman in his Honors Calculus 3 course, and took his courses all the way through his Mathematical Methods in Physics course, all of which were taught at a higher standard that encouraged me to do my best. He had also written many letters of recommendation for REUs and for graduate school. I feel that my attitude towards mathematics was shaped in great part by his, specifically valuing the ability and endurance to do (lengthy) computations and the ability to understand abstract concepts in concrete ways.

Another professor I would like to thank is Dr.\@ Lei Zhang for teaching me Fourier analysis. It was his enthusiastic teaching that made it quite clear to me that it is the most beautiful part of mathematics, and inspired me to do my own reading. The Fourier transform and its avatars have appeared several times during my undergraduate years, in different areas of mathematics (and this thesis!). I have the unshakeable feeling that I will need the rest of my life to meditate on these ideas.

I would like to thank my friends from the math department as well as from the physics department for their company and support. I hope to see them many times in the future. My highschool math teachers, Mrs.\@ Erika Sheehan and Mr.\@ Martin Doherty, and my highschool physics teacher, Dr.\@ Rajasekhar Narisetty, also deserve thanks due to initiating and encouraging my interest in mathematics.

Lastly, I would like to thank my parents for their lifelong support and love for me, and for having faith in me as I progress into the next stages of life.

\newpage\tableofcontents

\newpage\section*{Introduction}
\addcontentsline{toc}{section}{Introduction}
Modular forms are functions defined on the complex upper half plane that satisfy certain holomorphy conditions, and are invariant under a certain action of so-called \textit{congruence subgroups} of $\slz$, up to some factor. A modular form $f$ admits a Fourier series expansion 
\[f(z) = \sum_{n=0}^\infty a_ne^{2\pi i n z}, \] where $z$ is in the upper half plane and the $a_n$ are complex numbers. Often, the symbol $q$ is used in place of $e^{2\pi i z}$, so that the Fourier series expansion becomes the $q$-series expansion $f(z) = \sum_{n=0}^\infty a_nq^n$. Each modular form has an associated integer \textit{weight} $k$ and an integer \textit{level} $N$ which characterizes how it transforms under the action of a particular congruence subgroup. There are interesting connections between the theory of modular forms and other areas of mathematics, as we will see in the following examples.

A famous example of a modular form is the weight $12$, level $1$ \textit{modular discriminant} function $\Delta$ given by
\[\Delta(q) = q\prod_{n=1}^\infty (1-q^n)^{24},\quad q= e^{2\pi i z}.\]
(See \hyperlink{https://www.lmfdb.org/ModularForm/GL2/Q/holomorphic/1/12/a/a/}{LMFDB newform orbit 1.12.a.a}.) The \textit{Ramanujan tau function} $\tau$ denotes the Fourier coefficients of $\Delta$; that is, $\Delta(z) = \sum_{n=1}^\infty \tau(n)q^n$. The tau function satisfies several curious congruence relations, such as:
\begin{align*}
    \tau(n) &\equiv \sigma_{11}(n) \smod{2^{11}} \quad \text{for } n\equiv 1\smod{8}\\
    \tau(n) &\equiv 1217\sigma_{11}(n) \smod{2^{13}} \quad \text{for } n\equiv 3\smod{8}\\
    \tau(n) &\equiv  1537\sigma_{11}(n) \smod{2^{12}} \quad \text{for } n\equiv 5\smod{8},
\end{align*}
where $\sigma_k$ is the divisor function given by $\sigma_k(n) = \sum_{d\mid n}d^k$. In past decades, it was shown that the Fourier coefficients of other modular forms also satisfy similar kinds of congruences. As a result, the theory of congruences of modular forms has garnered significant interest in contemporary number theory.

Another interesting property of the tau function that Ramanujan observed in the early 1900s was that it was a multiplicative function: $\tau(mn) = \tau(m)\tau(n)$ for coprime integers $m,n$. Mordell gave a proof of this property in 1917, and in the late 1930s, Hecke provided a general theory involving Hecke operators that generalized Mordell's approach to spaces of modular forms and obtained similar results for modular forms. 

For another example of a modular form, consider the weight $2$, level $11$ modular form $f$ given by 
\[f(q) = q\prod_{n=1}^\infty (1-q^n)^2(1-q^{11n})^2,\quad q= e^{2\pi i z}.\]

(See \hyperlink{https://www.lmfdb.org/ModularForm/GL2/Q/holomorphic/11/2/a/a/}{LMFDB newform orbit 11.2.a.a}.) Formally expand this infinite product into the $q$-series $\sum_{n=1}^\infty a_nq^n$. The first few terms of this sum are $q - 2q^2-q^3+2q^4+q^5+2q^6-2q^7 -2q^9+\cdots$. Some interesting facts are that the coefficients $a_p$ for primes $p\neq 11$ satisfy $a_p\equiv p+1\smod 5$, and that the coefficients $a_n$ are multiplicative.

Now consider the \textit{elliptic curve} defined by the equation $y^2+y = x^3-x$. For each $p$, let $b_p$ be $p$ minus the number of solutions in $(\mathbb Z/p\mathbb Z)^2$ to the equation $y^2+y = x^3-x$. The first few values of $b_p$ for the primes $2,3,5,7$ are $-2,-1,1,-2$. But if we compare these values to $a_p$ for the same primes $2,3,5,7$, we find that $a_p = b_p$. Somewhat miraculously, it is true this equality holds for all primes $p$ except for $p =11$. 

The connection between modular forms and elliptic curves is quite profound. The book \cite{diamond} focuses on a result called the \textit{Modularity theorem}, which states that all rational elliptic curves may be obtained from modular forms. This was formerly called the Taniyama-Shimura conjecture, which was posed around 1955. A full proof of the modularity theorem appeared at the turn of the century, but partial progress by Wiles and Taylor on this result led to the proof of Fermat's Last Theorem.

In this thesis, we focus on the Hecke operators and prove the \textib{strong multiplicity one} theorem for modular forms. This result is usually stated using the language of \textit{automorphic} forms and representation theory. Automorphic forms are, broadly speaking, a generalization of modular forms in that they are well-behaved functions on the complex upper half plane with respect to the action of some action of a discrete topological (usually matrix) group. We focus on classical modular forms because they are a natural entry point into the theory of automorphic forms, and the strong multiplicity one theorem is essential to the development of the theory of modular forms. (``Classical'' modular forms are so named in order to distinguish them from other specialized modular forms appearing in the zoo of almost-invariant functions on the upper half plane with respect to the action of other matrix groups; for example, the Maass forms or the Bianchi modular forms.) 

We continue setting the stage for the strong multiplicity one theorem. The collection of all modular forms of a fixed weight and level forms a vector space. This vector space may be decomposed into the direct sum of the space of \textit{cusp forms} and the space of \textit{Eisenstein series}. The space of cusp forms is characterized by the property that the modular forms in this space vanish ``at infinity'' in a particular way. This property manifests in cusp forms as having constant term zero in their $q$-series expansions. For example, both the $\Delta$ function and $f$ above are cusp forms. An important feature of the space of cusp forms is that we may endow it with the Petersson inner product. The Eisenstein series are well studied modular forms with a number of other properties which we do not discuss.

The \textit{Hecke operators} $\abr{n}$ and $T_n$ are certain kinds of ``averaging'' operators that decompose the spaces of modular forms and cusp forms into the direct sum of what are called $\chi$\textit{-eigenspaces}, which interact nicely with the Hecke operators. The spectral properties of the Hecke operators are of great importance. For starters, the Hecke operators are normal operators on the space of cusp forms, so there exists an orthogonal basis of simultaneous eigenfunctions for the space of cusp forms. More interestingly, we have the following: If the modular form $g$ has Fourier series expansion $\sum_{n=0}^\infty t_ne^{2\pi i n z}$ with $t_1 = 1$, and is a simultaneous eigenfunction of the Hecke operators $T_n$, then 
\[T_ng = t_ng\] for all $n\geq 1$. Hecke eigenvalues are also multiplicative; that is, $t_{mn} = t_mt_n$ for $m,n$ coprime. Thus the Fourier coefficients of $f$, and the Ramanujan tau function $\tau(n)$, are multiplicative. 

The strong multiplicity one property of cusp forms is as follows: Let $g$ be a cusp form at level $M$ that is normalized; that is, $g$ has first Fourier coefficient equal to $1$, and is a common eigenfunction of all of the Hecke operators $T_n$. If there exists a cusp form $f$ that is a normalized eigenfunction $f$ of the Hecke operators $T_n$ at level $N$ (which need not be equal to $M$) whose Fourier coefficients, that is, its eigenvalues, agree with the Fourier coefficients of $g$ at all indices $n$ coprime to some integer $L$, we show that $g$ is equal to the cusp form $f$ at level $N$. In other words, the eigenvalues of $g$ determine the cusp form uniquely.

One elementary application of this result in the theory of modular forms is to determine a basis for the space of cusp forms. In \cite[Theorem 5.8.3]{diamond}, a particular set of cusp forms, belonging to an even finer collection of cusp forms called \textit{newforms}, is shown to span the space of cusp forms. To obtain the linear independence of these newforms, suppose that there is a nontrivial linear combination of elements in the spanning set that equals zero, $\sum_{i}c_if_i = 0$. One can show that this linear combination is of newforms that are eigenfunctions of the Hecke operators, with equal Fourier coefficients away from some fixed integer $L$. So by strong multiplicity one, each of the $f_i$ are equal, from which linear dependence follows.

We develop the necessary theory of classical modular forms and Hecke operators, mostly following \cite{diamond}, and then transition into the proof of strong multiplicity one in \cite[Theorem 4.6.19]{miyake}. This approach avoids the language of automorphic forms and representation theory, and instead obtains the result by studying the $L$-functions associated to modular forms and their associated Euler products. Almost all of of the results coming from \cite{miyake} have been reformulated to follow the more modern perspective found in \cite{diamond}, especially in the development of the Hecke operators.

\newpage\section{Preliminaries} \pagestyle{body} In this section, we collect the definitions and results needed to define modular forms. We describe the action of $\slz$ and its congruence subgroups on the upper half plane of the complex plane and introduce the congruence subgroups $\varGamma_0(N)$ and $\varGamma_1(N)$, which will appear in every section thereafter. We discuss modular curves and compactified modular curves briefly.
\subsection{The modular group and congruence subgroups}
Let $R$ be a unital ring. The \textib{general linear group} $\GL_n(R)$ is the group of $n\times n$ invertible matrices with entries from $R$; that is, matrices with unit determinant, under matrix multiplication. The \textib{special linear group} $\SL_n(R)$ is the (normal) subgroup of $\GL_n(R)$ whose matrices have determinant $1_R$.

\begin{definition}
    The \textib{modular group} \boldindex{$\slz$} is the group of $2\times 2$ integer-valued matrices with determinant $1$, and is a subgroup of $\glr$.
\end{definition} 
% Equivalently, the elements of $\slz$ are the integer-valued matrices $\big(\!\begin{smallmatrix}
%     a & b \\ c & d
% \end{smallmatrix}\!\big)$ with $\gcd(c,d) = 1$, since $ad-bc = 1$ is equivalent to $c$ and $d$ being coprime. 
It is well known (e.g., \cite[Chapter 7, Section 1.2, Theorem 2]{serre}) that 
\begin{equation}\label{eqn: generators modular group}
    \slz = \left\langle\begin{pmatrix}
        1 & 1 \\ 0 & 1
    \end{pmatrix}, \begin{pmatrix}
        0 & -1 \\ 1 & 0
    \end{pmatrix}\right\rangle, 
\end{equation}
and that $\slz$ acts on the Riemann sphere $\widehat{\mathbb{C}} = \mathbb{C}\cup \{\infty\}$ by fractional linear transformations
\[\abcd(z) = \frac{az + b}{cz + d}.\] For $c\neq 0$, $-d/c$ is sent to $\infty$ and $\infty$ is sent to $a/c$; if $c = 0$ then $\infty$ is sent to $\infty$. We check that this action defines a group action: It is clear that the identity matrix sends $z$ to itself, and that 
\[\begin{pmatrix}
    q & r \\ s & t
\end{pmatrix}\abcd(z) = \begin{pmatrix}
    q & r \\ s & t
\end{pmatrix}\Big(\frac{az+b}{cz+d}\Big) = \frac{q\frac{az+b}{cz+d}+r}{s\frac{az+b}{cz+d}+t} = \frac{(qa+rc)z+qb+rd}{(sa+tb)z+sb+td} = \begin{pmatrix}
    qa+rc & qb+rd \\ sa+tb & sb+td
\end{pmatrix}(z)\] for $\big(\!\begin{smallmatrix}
    q & r \\ s & t
\end{smallmatrix}\!\big),\big(\!\begin{smallmatrix}
    a & b \\ c & d
\end{smallmatrix}\!\big)\in \slz$. Similarly, $\glr$ acts on $\widehat{\mathbb C}$ by fractional linear transformations.
It follows that these transformations are (bi)holomorphic and are automorphisms.
% ; indeed, \[\abcd^{-1}\abcd(z) = \begin{pmatrix}
%     d & -b \\ -c & a
% \end{pmatrix}\Big(\frac{az+b}{cz+d}\Big) = \frac{\frac{d(az+b)}{cz+d}- b}{\frac{-c(az+b)}{cz+d}+a} = \frac{(ad-bc)z}{ad-bc} = z.\]
Furthermore, both $I$ and $-I$ act as the identity on $\widehat{\mathbb{C}}$, so for any $A\in \slz$, the actions of $A$ and $-A$ agree. The generators in \eqref{eqn: generators modular group} correspond to the maps
\[z\mapsto z+1\quad\text{and}\quad z\mapsto -1/z.\]
We consider particular subgroups of the modular group. 
\begin{definition}
    For $N$ a positive integer, the \textib{principal congruence subgroup of level $N$}, \boldindex{$\varGamma(N)$}, is given by the subgroup
\[\varGamma(N) = \cbr{\abcd\in \slz: \abcd\equiv \begin{pmatrix}
    1 & 0 \\ 0 & 1
\end{pmatrix}\smod N}\]
with the matrix congruence interepreted as congruence modulo $N$ entrywise.
\end{definition}
Alternatively, $\varGamma(N) = (I+N\Mat_2(\mathbb Z))\cap \slz$. Furthermore, observe that if $N\mid M$, then $\varGamma(M)\subset \varGamma(N)$, since $x \equiv y\smod M$ implies $x\equiv y\smod N$ whenever $N\mid M$.
\begin{lemma}\label{lem: varGamma(N) kernel of projection}
    The principal congruence subgroup $\varGamma(N)$ is the kernel of the natural homomorphism $\slz\to \SL_2(\mathbb{Z}/N\mathbb{Z})$, so $\varGamma(N)$ is normal in $\slz$; furthermore, the natural homomorphism is surjective.
\end{lemma}
\begin{proof}
    It is clear that $\varGamma(N)$ is the kernel of the natural homomorphism $\slz\to \SL_2(\mathbb{Z}/N\mathbb{Z})$, so we prove that this map is surjective. When $N = 1$, the natural homomorphism is the zero map (so $\varGamma(1) = \slz$). Let $N>1$, and consider an element \[\begin{pmatrix}
        \overline a & \overline b \\ \overline c & \overline d
    \end{pmatrix}\in \SL_2(\mathbb{Z}/N\mathbb{Z}),\] so that $\overline{ad-bc} = \overline 1$ (here $\overline{\,\cdot\,}$ denotes reduction modulo $N$). Then for some integer $k$, $ad-bc = 1+kN$. It follows that $1+kN$ is a multiple of $g = \gcd(c,d)$, so that $\gcd(g,N)=1$ and $\gcd(c,d,N)=1$. We show that there exist integers $i,j$ such that $c+iN$ and $d+jN$ are coprime.
    
    % Suppose there is no choice of $i,j$ such that $c+iN$ and $d+jN$ are coprime, so that for fixed $i$, $\gcd(c+iN,d+jN)>1$ for every $j$. By Dirichlet's theorem on arithmetic progressions, the sequence $(d+mN)$ contains infinitely many primes, which yields a contradiction.
    
    If $c\neq 0$, consider a solution $j$ to the system of congruences 
    \[\begin{cases}
        j\equiv 1\smod{p} & p\mid g\\
        j\equiv 0\smod{p} & p\nmid g, p\mid c,
    \end{cases}\] which may be obtained via the Chinese remainder theorem. Then $\gcd(c,d+jN) = 1$ since any prime $p$ dividing $c$ will not divide $d+jN$ with $j$ chosen as above (of primes $p$ dividing $c$, when $p\mid g$, we have $p\nmid N$, and when $p\nmid g$, we have $p\nmid d$). If $c = 0$, then $d\neq 0$ and repeat this argument with $d,i$ in place of $c,j$ respectively.

    With integers $c+iN$ and $d+jN$ coprime, there exist integers $s,t$ with $s(c+iN) + t(d+jN) = 1$. Then \begin{multline*}
        \det\begin{pmatrix}
            a-(k+aj-bi)tN & b+(k+aj-bi)sN \\ c+iN & d+jN
        \end{pmatrix} = ad+ajN-(k+aj-bi)tN(d+jN)\\-bc-biN-(k+aj-bi)sN(c+iN) = 1+kN+ajN-biN-(k+aj-bi)N = 1
    \end{multline*} and 
    \[\begin{pmatrix}
        \overline{a-(k+aj-bi)tN} & \overline{b+(k+aj-bi)sN} \\ \overline{c+iN} & \overline{d+jN}
    \end{pmatrix} = \begin{pmatrix}
        \overline a & \overline b \\ \overline c & \overline d
    \end{pmatrix}\] as needed.
\end{proof}
Therefore, by the first isomorphism theorem, $\slz/\varGamma(N)\cong \SL_2(\mathbb{Z}/N\mathbb{Z})$. 
\begin{lemma}
    The index of $\varGamma(N)$ in $\slz$ is $[\slz: \varGamma(N)] = N^3\prod_{p\mid N}(1-1/p^2)$.
\end{lemma}
\begin{proof}
    We first show that $\abs{\SL_2(\mathbb{Z}/p^e\mathbb{Z})} = p^{3e}(1-1/p^2)$ for a prime $p$ by induction on $e$.

    The determinant is a surjective homomorphism from $\GL_2(\mathbb{Z}/p\mathbb{Z})$ to $(\mathbb{Z}/p\mathbb{Z})^\times\cong \mathbb{Z}/(p-1)\mathbb{Z}$ with kernel $\SL_2(\mathbb{Z}/p\mathbb{Z})$. Observe that $\abs{\GL_2(\mathbb{Z}/p\mathbb{Z})}$ is the number of ordered bases of $(\mathbb{Z}/p\mathbb{Z})^2$, given by $(p^2-1)(p^2-p)$. (The first factor counts the number of admissible first basis vectors, and the second factor counts the number of admissible second basis vectors after choosing a first basis vector.) Then by the first isomorphism theorem, $\abs{\SL_2(\mathbb{Z}/p\mathbb{Z})} = (p^2-1)(p^2-p)/(p-1) = p(p^2-1) = p^3(1-1/p^2)$.
    
    By \cref{lem: varGamma(N) kernel of projection}, the natural homomorphism $\slz\to \SL_2(\mathbb{Z}/p^e\mathbb{Z})$ is surjective. The surjection $\slz\to \SL_2(\mathbb{Z}/p^e\mathbb{Z})$ is equal to the composition of the natural homomorphisms $\slz\to \SL_2(\mathbb{Z}/p^{e+1}\mathbb{Z})$ and $\SL_2(\mathbb{Z}/p^{e+1}\mathbb{Z})\to \SL_2(\mathbb{Z}/p^e\mathbb{Z})$, from which it follows that $\SL_2(\mathbb{Z}/p^{e+1}\mathbb{Z})\to \SL_2(\mathbb{Z}/p^e\mathbb{Z})$ is surjective also.
    % \[\begin{tikzcd}
    %     {\slz} & {\SL_2(\mathbb{Z}/p^e\mathbb{Z})} \\
    %     {\SL_2(\mathbb{Z}/p^{e+1}\mathbb{Z})}
    %     \arrow[from=2-1, to=1-2]
    %     \arrow[two heads, from=1-1, to=2-1]
    %     \arrow[two heads, from=1-1, to=1-2]
    % \end{tikzcd}\]
    % from which we find that the map $\SL_2(\mathbb{Z}/p^{e+1}\mathbb{Z})\to \SL_2(\mathbb{Z}/p^e\mathbb{Z})$ is surjective.

    Any element $\gamma$ of $\ker(\SL_2(\mathbb{Z}/p^{e+1}\mathbb{Z})\to \SL_2(\mathbb{Z}/p^e\mathbb{Z}))$ is of the form
    \[\begin{pmatrix}
        \overline{1+ip^e} & \overline{rp^e} \\ \overline{sp^e} & \overline{1+jp^e}
    \end{pmatrix},\] for $i,j,r,s\in \{0,\dots,p-1\}$ and $i = p-j$ (so that $\det \gamma = 1$). Hence $\abs{\ker(\SL_2(\mathbb{Z}/p^{e+1}\mathbb{Z})\to \SL_2(\mathbb{Z}/p^e\mathbb{Z}))} = p^3$. Suppose that $\abs{\SL_2(\mathbb{Z}/p^e\mathbb{Z})} = p^{3e}(1-1/p^2)$. Then by the first isomorphism theorem, $\abs{\SL_2(\mathbb{Z}/p^{e+1}\mathbb{Z})} = p^3\abs{\SL_2(\mathbb{Z}/p^e\mathbb{Z})} = p^3p^{3e}(1-1/p^2) = p^{3(e+1)}(1-1/p^2)$ as needed. By induction, $\abs{\SL_2(\mathbb{Z}/p^e\mathbb{Z})} = p^{3e}(1-1/p^2)$ for any $e$.

    There is an isomorphism of matrix groups $\Mat_2(\prod_{i=1}^nR_i)\cong \prod_{i=1}^n\Mat_2(R_i)$ for rings $R_i$ which restricts to the isomorphisms $\GL_2(\prod_{i=1}^nR_i)\cong \prod_{i=1}^n\GL_2(R_i)$ and $\SL_2(\prod_{i=1}^nR_i)\cong \prod_{i=1}^n\SL_2(R_i)$ since $(\prod_{i=1}^nR_i)^\times\cong \prod_{i=1}^nR_i^\times$ and $1_{\prod_{i=1}^nR_i} = \prod_{i=1}^n1_{R_i}$.

    Let $N$ have prime factorization $N = \prod_{p\mid N} p^{e}$. By the Chinese remainder theorem, $\mathbb{Z}/N\mathbb{Z}\cong \prod_{p\mid N}\mathbb{Z}/p^{e}\mathbb{Z}$. It follows that $\SL_2(\mathbb{Z}/N\mathbb{Z})\cong \prod_{p\mid N}\SL_2(\mathbb{Z}/p^{e}\mathbb{Z})$, from which we have \[\abs{\SL_2(\mathbb{Z}/N\mathbb{Z})}= \prod_{p\mid N}\abs{\SL_2(\mathbb{Z}/p^{e}\mathbb{Z})} = \prod_{p\mid N}p^{3e}(1-1/p^2) = N^3\prod_{p\mid N}(1-1/p^2)\] as desired.
\end{proof}

\begin{definition}\label{def: congruence subgroup}
    A subgroup $\varGamma$ of $\slz$ is a \boldindex{congruence subgroup} \textib{\textup{(}of level $N$\textup{)}} if for some positive integer $N$, $\varGamma(N)\subset \varGamma$.
\end{definition}
Note that congruence subgroups need not be normal in $\slz$. % terminology for congruence subgps, see keith conrad page

Let $\varGamma$ be a level $N$ congruence subgroup. We have $[\slz : \varGamma(N)] = [\slz:\varGamma][\varGamma:\varGamma(N)]$, and since $[\slz : \varGamma(N)]$ is finite, $[\slz:\varGamma]$ must also be finite; that is, congruence subgroups have finite index in $\slz$.

\begin{definition}
    Frequently used congruence subgroups are 
    \begin{align*}
        \varGamma_0(N) &= \cbr{\abcd\in\slz: \abcd\equiv\begin{pmatrix}
            \ast & \ast \\ 0 & \ast
        \end{pmatrix}\smod N}\text{ and}\\
        \varGamma_1(N) &= \cbr{\abcd\in \slz: \begin{pmatrix}
            a & b \\ c & d 
        \end{pmatrix}\equiv \begin{pmatrix}
            1 & \ast \\ 0 & 1
        \end{pmatrix}\smod N},
    \end{align*}
    where $\ast$ denotes unspecified quantities and the matrix congruences are to be taken entrywise.    
\end{definition}
For any positive integer $N$, the inclusions $\varGamma(N)\subset\varGamma_1(N)\subset\varGamma_0(N)\subset\slz$ hold.

\begin{lemma}\label{lem: principal normal in gamma 1}
    The map $\varGamma_1(N)\to \mathbb{Z}/N\mathbb{Z}$ given by $\smallabcd\mapsto \overline b$ is surjective with kernel $\varGamma(N)$. 
\end{lemma}
\begin{proof}
    For $\overline b\in\mathbb{Z}/N\mathbb{Z}$, $\big(\!\begin{smallmatrix}
        1 & b \\ 0 & 1
    \end{smallmatrix}\!\big)$ has unit determinant and maps to $\overline b$ under the above map. It is clear that the kernel of the above map is $\varGamma(N)$.
    % Let $\overline b\in\mathbb{Z}/N\mathbb{Z}$ be given. Let $a,d\equiv 1\smod N$ and $c\equiv 0\smod N$, and write $a= iN+1$, $d = jN+1$, and $c = kN$ for some integers $i,j,k$. If $1= ad-bc = ijN^2 + (i+j)N + 1 - bkN$, then $(ijN+ i + j)/k = b$. Take $k=1$ and reduce modulo $N$ to obtain the equation $i+j\equiv b \smod N$, which has many solutions, for example $i = b$ and $j =0$. These choices for $i,j,k$ produce a preimage of $b\smod N$ under the above map. Indeed, \[\begin{pmatrix}
    %     a & b \\ c & d
    % \end{pmatrix} = \begin{pmatrix}
    %     bN+1 & b \\ N & 1
    % \end{pmatrix}\mapsto b\quad\text{and}\quad \det\begin{pmatrix}
    %     bN+1 & b \\ N & 1
    % \end{pmatrix} = 1.\]
    % If $A\in \varGamma_1(N)$ and \[A = \begin{pmatrix}
    %     a & b \\ c & d
    % \end{pmatrix}\mapsto \overline 0,\] then $b \equiv 0\smod N$ for some integer $k$. Since $a,d\equiv 1\smod N$ and $c\equiv 0\smod N$, it follows that $A\in \varGamma(N)$. Certainly any element of $\varGamma(N)$ is contained in the kernel of the above map.
\end{proof} 
It follows that $\varGamma_1(N)/\varGamma(N)\cong \mathbb{Z}/N\mathbb{Z}$, so that the index of $\varGamma(N)$ in $\varGamma_1(N)$ is $N$. What follows is a similar result.

\begin{lemma}\label{lem: gamma 1 normal in gamma 0}
    The map $\varGamma_0(N)\to (\mathbb{Z}/N\mathbb{Z})^\times$ given by $\smallabcd\mapsto \overline d$ 
    is surjective with kernel $\varGamma_1(N)$.
\end{lemma}
\begin{proof}
    % Exercise 1.2.3. (d) in \cite{diamond}.
    % DO NOT DELETE SOLUTION, but recommended to omit solution
    Let $\overline d \in (\mathbb{Z}/N\mathbb{Z})^\times$ be given, and note that $\gcd(d,N)=1$. Let $c \equiv 0\smod N$ so that $c=kN$ for some integer $k$. Choose $k$ coprime to $d$ so that $\gcd(d,kN) = \gcd(d,c) = 1$. Then there exist integers $a,b$ with $ad-bc = 1$. Then \[\begin{pmatrix}
        a & b \\ c & d
    \end{pmatrix}\mapsto d\quad\text{and}\quad \det\begin{pmatrix}
        a & b \\ c & d
    \end{pmatrix} = 1.\] 
    If $A\in \varGamma_0(N)$ and \[A = \begin{pmatrix}
        a & b \\ c & d
    \end{pmatrix}\mapsto \overline 1,\] so $d = jN+1$ for some integer $j$. Since $c\equiv 0\smod N$, $c= kN$ and $\gcd(d,c) = 1$. There exist integers $a,b$ with $1 = ad-bc = a(jN+1) + bkN = a+ (aj + bk)N$. Reducing modulo $N$ gives $a\equiv 1\smod N$. It follows that $A\in \varGamma_1(N)$. Conversely, if $A\in \varGamma_1(N)$, then $A\mapsto \overline 1$.
\end{proof}
Hence $\varGamma_0(N)/\varGamma_1(N)\cong (\mathbb{Z}/N\mathbb{Z})^\times$, so that the index of $\varGamma_1(N)$ in $\varGamma_0(N)$ is $\phi(N) = \abs{(\mathbb{Z}/N\mathbb{Z})^\times} = N\prod_{p\mid N}(1-1/p)$, where $\phi$ denotes the Euler totient function. 
\begin{corollary}
    We obtain \[[\slz : \varGamma_0(N)] = \frac{[\slz : \varGamma(N)]}{[\varGamma_0(N) : \varGamma_1(N)][\varGamma_1(N) : \varGamma(N)]} = \frac{N^3\prod_{p\mid N}(1-1/p^2)}{N\cdot N\prod_{p\mid N}(1-1/p)} = N\prod_{p\mid N}(1+1/p).\]
\end{corollary} 

\subsection{Modular curves}
Recall that $\slz$, and more broadly $\glr$, acts on the Riemann sphere $\widehat{\mathbb{C}}$ by fractional linear transformations \[\abcd(z) = \frac{az+b}{cz+d}.\] 
If $\gamma\in \glr$ with $\gamma = \smallabcd$ and $z\in \mathbb C$, the computation $(az+b)(c\overline z + d) = ac\abs{z}^2+bd + adz + bc\overline z = ac\abs{z}^2+bd + adz + (ad-\det\gamma)\overline z = ac\abs{z}^2+bd + ad(z+\overline z) -\det\gamma\overline z$ shows that
\begin{equation}\label{eqn: imaginary parts}
    \Im(\gamma(z)) = \Im\Big(\frac{az+b}{cz+d}\Big) = \frac{\Im((az+b)(c\overline z + d))}{\abs{cz+d}^2} = \frac{\Im(-\det\gamma\overline z)}{\abs{cz+d}^2} = \frac{\det\gamma\Im(z)}{\abs{cz+d}^2}.
\end{equation}
From the above calculation it follows that if $\det \gamma >0$ and $z$ belongs to the \textib{upper half plane} \boldindex{$\mathcal H$} $= \cbr{z\in\mathbb{C} : \Im(z)>0}$, then $\gamma(z)\in\mathcal H$. In particular, it follows that $\glrp = \{\gamma\in \glr : \det \gamma > 0\}$, $\slz$, and any congruence subgroup $\varGamma$ of $\slz$ act on the upper half plane $\mathcal H$ by fractional linear transformations.

\begin{definition}
    Let $\varGamma$ be a congruence subgroup of $\slz$. The \textib{modular curve} \boldindex{$Y(\varGamma)$} is the orbit space $\varGamma\backslash \mathcal{H} = \cbr{\varGamma z : z\in\mathcal H}$.
\end{definition}
The modular curves for the congruence subgroups $\varGamma_0(N), \varGamma_1(N)$, and $\varGamma(N)$ are denoted by $Y_0(N), Y_1(N)$, and $Y(N)$, respectively.

We topologize modular curves and briefly show they are compact Riemann surfaces. Let $\varGamma$ be a congruence subgroup. The upper half plane $\mathcal H$ is given the Euclidean topology (as a subspace of $\mathbb{C}$ or of $\mathbb{R}^2$), and the natural surjection \[\pi \colon \mathcal H \to Y(\varGamma)\quad \text{defined by}\quad  z\mapsto \varGamma z\] induces the quotient topology on $Y(\varGamma)$ (i.e., the open sets in $Y(\varGamma)$ are those with open preimages under $\pi$).

\begin{lemma}\label{lem: slz acts as open maps}
    The quotient map $\pi \colon \mathcal H \to Y(\varGamma)$ is an open map.
\end{lemma}
\begin{proof}
    It suffices to show that the projection $\pi$ takes an open disk $B$ in $\mathcal{H}$ to an open set in $Y(\varGamma)$; that is, to show that \[\bigcup_{\gamma\in \varGamma}\gamma(B)\] is open. Since each $\gamma^{-1}$ is continuous, each $\gamma(B)$ is open, so $\pi$ is an open map.
\end{proof}

\begin{proposition}\label{prop: properly discontinuous}
    The action of the modular group $\slz$ on $\mathcal H$ is properly discontinuous; that is, for $z_1,z_2\in \mathcal H$, (including the case $z_1 = z_2$) there exist neighborhoods $U_1$ of $z_1$ and $U_2$ of $z_2$ in $\mathcal H$ such that for any $\gamma\in \slz$, if $\gamma(U_1)\cap U_2\neq \emptyset$, then $\gamma(z_1) = z_2$. Equivalently, for any $\gamma\in\slz$, if $\gamma(z_1)\neq z_2$, then $\gamma(U_1)\cap U_2= \emptyset$.
\end{proposition}
\begin{proof}
    %\sai{i kind of don't think I need this proof, could just cite it}
    Let $V_1$ and $V_2$ be neighborhoods of $z_1$ and $z_2$, respectively, with compact closure (e.g. open disks or bounded open sets). Let $y_1 = \sup_{z\in V_1}\cbr{\Im(z)}$, $Y_1 = \sup_{z\in V_1}\cbr{\Im(z)}$, and $y_2 = \inf_{z\in V_2}\cbr{\Im(z)}$. Then from \eqref{eqn: imaginary parts}, we have for $\gamma = \big(\!\begin{smallmatrix}
        a & b \\ c & d
    \end{smallmatrix}\!\big)\in \slz$ not equal to $\pm I$ (i.e., $c\neq 0$) and $z\in V_1$ that \[\Im(\gamma(z)) = \frac{\Im(z)}{\abs{cz+d}^2}\leq \frac{\Im(z)}{c^2\abs{z}^2}\leq \frac{1}{c^2\Im(z)}\leq \frac{1}{c^2y_1}\] and \[\Im(\gamma(z)) = \frac{\Im(z)}{\abs{cz+d}^2}\leq \frac{Y_1}{\Re(cz+d)^2} = \frac{Y_1}{(c\Re(z)+d)^2},\] so that $\Im(\gamma(z))\leq\min\cbr{1/c^2y_1,Y_1/(c\Re(z)+d)^2}$. All but finitely many integers $c$ may be chosen so that $1/c^2y_1< y_2$, and of the finitely many $c$ where this inequality does not hold, we may choose all but finitely many $d$ such that $Y_1/(c\Re(z)+d)^2< y_2$ uniformly in $z$ (since $V_1$ has compact closure). Thus for all but finitely many integers $c,d$, if $\gamma  = \big(\!\begin{smallmatrix}
        a & b \\ c & d
    \end{smallmatrix}\!\big)\in \slz$, then \begin{equation}\label{eqn: sup inf}
        \sup_{z\in V_1}\cbr{\Im(\gamma(z))}< \inf_{z\in V_2}\cbr{\Im(z)}.
    \end{equation} Thus for $\gamma = \big(\!\begin{smallmatrix}
        a & b \\ c & d
    \end{smallmatrix}\!\big)\in \slz$ satisfying \eqref{eqn: sup inf}, $\gamma(V_1)\cap V_2=\emptyset$. We show that there are only finitely many $\gamma\in \slz$ for which $\gamma(V_1)\cap V_2\neq\emptyset$.

    We determine the matrices in $\slz$ which have fixed bottom row $(c,d)$. This is equivalent to finding all matrices $\eta = \big(\!\begin{smallmatrix}
        A & B \\ C & D
    \end{smallmatrix}\!\big)\in \slz$ such that $\eta\gamma$ has bottom row $(c,d)$, where $\gamma\in \slz$ is a fixed matrix with  bottom row $(c,d)$. Evidently $C=0$ and $D = 1$, and since $1 = AD - BC = A$, we have that $\eta$ must have the form $\big(\!\begin{smallmatrix}
        1 & B \\ 0  & 1
    \end{smallmatrix}\!\big)$, and $B$ may be taken to be any integer. Explicitly, the matrices in $\slz$ with bottom row $(c,d)$ are \[\cbr{\begin{pmatrix}
        1 & k \\ 0 & 1
    \end{pmatrix}\begin{pmatrix}
        a & b \\ c & d
    \end{pmatrix}: k\in \mathbb{Z}},\]
    where $(a,b)$ is one such pair such that $ad-bc = 1$. Thus for $\gamma\in \cbr{\big(\!\begin{smallmatrix}
        1 & k \\ 0 & 1
    \end{smallmatrix}\!\big)\big(\!\begin{smallmatrix}
        a & b \\ c & d
    \end{smallmatrix}\!\big): k\in \mathbb{Z}}$, the intersection $\gamma(V_1)\cap V_2 = \big(\big(\!\begin{smallmatrix}
        a & b \\ c & d
    \end{smallmatrix}\!\big)V_1+ k\big)\cap V_2 =\emptyset$ for all but finitely many $k$; that is, for finitely many $\gamma$. So of the finitely many pairs of integers $(c,d)$ for which \eqref{eqn: sup inf} does not hold, only finitely many matrices $\gamma\in \slz$ with bottom row $(c,d)$ exist with $\gamma(V_1)\cap V_2\neq \emptyset$.

    Let $F = \cbr{\gamma\in \slz: \gamma(V_1)\cap V_2\neq \emptyset, \gamma(z_1)\neq z_2}$, which is finite by the above argument. For each $\gamma \in F$ there exist disjoint neighborhoods $V_{1,\gamma}$ of $\gamma(z_1)$ and $V_{2,\gamma}$ of $z_2$ in $\mathcal H$. Let \[U_1 = V_1\cap \Bigg(\bigcap_{\gamma\in F}\gamma^{-1}(V_{1,\gamma})\Bigg)\quad\text{and}\quad U_2 = V_2\cap \Bigg(\bigcap_{\gamma\in F}V_{2,\gamma}\Bigg),\] and note that $U_1,U_2$ are open, as elements of $\slz$ are open maps on $\mathcal H$ by $\cref{lem: slz acts as open maps}$. Then take any $\gamma\in \slz$ with $\gamma(U_1)\cap U_2\neq \emptyset$. If $\gamma\not\in F$, then we must have $\gamma(z_1) = z_2$. Suppose that $\gamma\in F$. Then $U_1\subset \gamma^{-1}(U_{1,\gamma})$ and $U_2\subset U_{2,\gamma}$ so that $ \gamma(U_1)\cap U_2\subset U_{1,\gamma}\cap U_{2,\gamma}$. But $\gamma(U_1)\cap U_2\neq \emptyset$, which is in contradiction with $U_{1,\gamma}$ and $U_{2,\gamma}$ chosen to be disjoint. Hence $\gamma\not\in F$, so that $\gamma(z_1)= z_2$.
\end{proof}

The equality
\begin{multline}\label{eqn: disjoint nbds in modular curve}
    \emptyset = \pi^{-1}(\pi(U_1)\cap \pi(U_2)) = \Bigg(\bigcup_{\gamma\in \varGamma}\gamma(U_1)\Bigg)\cap \Bigg(\bigcup_{\gamma\in \varGamma}\gamma(U_2)\Bigg) = \bigcup_{\gamma,\eta\in \varGamma}\gamma(U_1)\cap \eta(U_2)\\ = \bigcup_{\eta\in \varGamma}\eta\Bigg[\Bigg(\bigcup_{\gamma\in \varGamma}(\eta^{-1}\gamma)(U_1)\Bigg)\cap U_2\Bigg] = \bigcup_{\eta\in \varGamma}\eta\Bigg[\Bigg(\bigcup_{\gamma^\prime\in \varGamma}\gamma^\prime(U_1)\Bigg)\cap U_2\Bigg]
\end{multline}
implies that $\big(\bigcup_{\gamma^\prime\in \varGamma}\gamma^\prime(U_1)\big)\cap U_2$ must be empty, which proves the following: $\pi(U_1)\cap \pi(U_2)=\emptyset$ in $Y(\varGamma)$ is equivalent to $\big(\bigcup_{\gamma\in \varGamma}\gamma(U_1)\big)\cap U_2 = \emptyset$ in $\mathcal H$.

\begin{proposition}
    The space $Y(\varGamma)$ is second-countable, connected, and Hausdorff.
\end{proposition}

\begin{proof}
    Since $\pi$ is open and $\mathcal H$ is second-countable, $Y(\varGamma)$ is second-countable. As $\mathcal H$ is connected and $\pi$ is continuous, $Y(\varGamma)$ is connected.

    Let $\pi(z_1)$ and $\pi(z_2)$ be distinct points in $Y(\varGamma)$ (so that $\gamma(z_1)\neq z_2$ for any $\gamma\in \varGamma$), and take neighborhoods $U_1$ of $z_1$ and $U_2$ of $z_2$ such that for any $\gamma\in \slz$, if $\gamma(U_1)\cap U_2\neq \emptyset$, then $\gamma(z_1) = z_2$, as per the previous result. Then $\big(\bigcap_{\gamma\in\varGamma}\gamma(U_1)\big)\cap U_2 = \emptyset$ since $\gamma(z_1)\neq z_2$ for every $\gamma\in \varGamma$. From the discussion surrounding \eqref{eqn: disjoint nbds in modular curve}, we have that $\pi(U_1)\cap\pi(U_2) = \emptyset$ as needed, with $\pi(U_1),\pi(U_2)$ open since $\pi$ is an open mapping.
\end{proof}

What remains is to compactify $Y(\varGamma)$ and to put charts on the resulting compact space, which we do without verifying the details. 

The group $\glqp = \{\gamma\in \GL_2(\mathbb Q) : \det \gamma >0\}$ acts on $\mathbb Q \cup \{\infty\}$ by 
\[\abcd \Big(\frac{r}{s}\Big) = \frac{ar+bs}{cr+ds},\] taking $\infty$ to $a/c$ and $-d/c$ to $\infty$ when $c\neq 0$, and taking $\infty$ to $\infty$ when $c = 0$. Since $\smallabcd$ is invertible, the action above never produces indeterminate forms like $0/0$. Furthermore, $\slz$ acts transitively: Any rational number is of the form $a/c$, where $a$ and $c$ are coprime. Choose $b,d$ so that $ad-bc=1$, from which we have $\smallabcd(\infty) = a/c$.

Let $\mathcal H^\ast = \mathcal H\cup \mathbb Q\cup\{\infty\}$ and define the \textib{compactified modular curve} \boldindex{$X(\varGamma)$} as a set by \[X(\varGamma) = \varGamma\backslash\mathcal H^\ast = Y(\varGamma)\cup\varGamma\backslash(\mathbb Q\cup \{\infty\}).\] Call the points $\varGamma s\in \varGamma\backslash(\mathbb Q\cup \{\infty\})$ \boldindex{cusp}\textib{s} of $X(\varGamma)$ (or of $\varGamma$). Denote by $X_0(N)$, $X_1(N)$, and $X(N)$ the modular curves $X(\varGamma_0(N))$, $X(\varGamma_1(N))$, and $X(\varGamma(N))$ respectively.

\begin{lemma}\label{lem: finitely many cusps}
    The modular curve $X(\varGamma)$ has finitely many cusps for any congruence subgroup $\varGamma$.
\end{lemma}

\begin{proof}
    The number of cusps of $X(\varGamma)$ is equal to the index of $\varGamma$ in $\slz$. Let $\slz$ act on $\varGamma\backslash(\mathbb Q \cup \{\infty\})$ by $\gamma(\varGamma s) = \varGamma (\gamma s)$. The action is well defined: If $\varGamma s = \varGamma t$, then $s = \gamma^\prime t$ for some $\gamma^\prime \in \varGamma$. Then $\varGamma \gamma\gamma^\prime s = \varGamma\gamma^{\prime\prime}\gamma s = \varGamma s$ for some $\gamma^{\prime\prime}\in\varGamma$. This group action is transitive since $\slz$ acts transitively on $\mathbb Q \cup \{\infty\}$, and the isotropy group of $\varGamma\backslash(\mathbb Q \cup \{\infty\})$ is $\varGamma$. Thus $[\slz:\varGamma] = \abs{\varGamma\backslash (\mathbb Q\cup \{\infty\})}$.
\end{proof}

It follows that the modular curve $X(1) = \slz\backslash \mathcal H^\ast$ has one cusp, namely, $\infty$.

A base for the topology of $\mathcal H^\ast$ consists of open disks centered at elements of $\mathcal H$ and the neighborhoods 
\[\alpha(\cbr{z\in\mathcal H : \Im(z) > M}\cup\{\infty\})\quad \text{for $M>0,\alpha\in\slz$},\] images of disks centered at $\infty$ under elements of $\slz$. Give $\mathcal H$ the topology these sets generate. As fractional linear transformations are conformal maps, for elements $\gamma\in\slz$ with $\gamma(\infty)\in\mathbb Q$, a disk centered at $\infty$ is mapped to a disk that is tangent to the real axis, containing one rational number.

Choosing this base ensures that $X(\varGamma)$ is Hausdorff (whereas taking open disks at each point in $\mathbb Q \cup \{\infty\}$ would not since $\mathbb Q$ is dense in the real line). Furthermore, every $\gamma\in \slz$ is a homeomorphism of $\mathcal H^\ast$ with itself. Give $X(\varGamma)$ the quotient topology with the quotient map $\pi\colon \mathcal H^\ast\to X(\varGamma)$ defined by extending the natural projection.

\begin{proposition}
    The modular curve $X(\varGamma)$ is Hausdorff, connected, and compact. 
\end{proposition}
\begin{proof}
    See \cite[Proposition 2.4.2]{diamond}.
\end{proof}

Before defining charts on $X(\varGamma)$, we collect some results about elliptic points.

\begin{definition}
    Let $\varGamma$ be a congruence subgroup. For $z\in\mathcal H$ let $\varGamma_z = \{\gamma\in\varGamma : \gamma(z) = z\}$, the isotropy subgroup of $z$. Call $z\in\mathcal H$ (and the corresponding point $\pi(z)\in X(\varGamma)$) an \boldindex{elliptic point} of $\varGamma$ if $\varGamma_z$ is nontrivial; that is, if the containment $\{\pm I\}\subset \{\pm I\}\varGamma_z$ is proper.
\end{definition}

\begin{proposition}\label{prop: cyclic isotropy subgp}
    Let $\varGamma$ be a congruence subgroup. Then $X(\varGamma)$ has finitely many elliptic points, and for each elliptic point $z$ of $\varGamma$, its isotropy subgroup $\varGamma_z$ is finite cyclic.
\end{proposition}
\begin{proof}
    See \cite[Proposition 2.3.3, Corollaries 2.3.4, 2.3.5]{diamond}.
\end{proof}

It follows that every point $z\in\mathcal H$ has an associated positive integer
\[h_z = \abs{\{\pm I\}\varGamma_z/\{\pm I\}} = \begin{cases}
    \abs{\varGamma_z}/2 & \text{if $-I\in\varGamma_z$}\\
    \abs{\varGamma_z} & \text{if $-I\not\in\varGamma_z$}
\end{cases}\] called the period of $z$. Note that $h_z>1$ when $z$ is an elliptic point, and that for $\gamma\in\slz$, the period of $\gamma(z)$ under $\gamma\varGamma\gamma^{-1}$ is the same as the period of $z$ under $\varGamma$. As a result, $h_z = h_{\gamma z}$ for all $\gamma\in\varGamma$, so that the period of $\varGamma z\in Y(\varGamma)\subset X(\varGamma)$ is well defined.

We now define charts for $X(\varGamma)$, starting by defining them on $Y(\varGamma)$. Since $\slz$ acts properly discontinuously on $\mathcal H$, we have the following:

\begin{corollary}[to \cref{prop: properly discontinuous}]
    Let $\varGamma$ be a congruence subgroup. Then each point $z\in\mathcal H$ is contained in a neighborhood $U\subset \mathcal H$ such that for all $\gamma\in\varGamma$, if $\gamma(U)\cap U \neq \emptyset$ then $\gamma\in\varGamma_z$. Furthermore, $U$ has no elliptic points except possibly $z$.
\end{corollary}

Given a point $\pi(z)\in Y(\varGamma)$, take a neighborhood $U$ for $z$ as in the corollary above. Let $\delta_z  = \big(\!\begin{smallmatrix}
    1 & -z \\ 1 & -\overline{z}
\end{smallmatrix}\!\big)\in \GL_2(\mathbb C)$ be the map from $\mathcal H$ (or $U$ by restriction) to $\mathbb C$ which takes $z$ to $0$ and $\overline{z}$ to $\infty$. Note that the isotropy subgroup of $0$ in the transformation group $(\delta_z\cbr{\pm I}\varGamma\delta_z^{-1})_0/\cbr{\pm I}$ is the conjugate $\delta_z(\cbr{\pm I}\varGamma/\cbr{\pm I})\delta_z^{-1}$ of the isotropy subgroup of $z$, hence is cyclic of order $h_z$ as a group of transformations by \cref{prop: cyclic isotropy subgp} (note $h_z$ may be equal to $1$). Since this group of fractional linear transformations fixes $0$ and $\infty$, these maps are given by $w\mapsto aw$, which are rotations about the origin through angular multiples of $2\pi/h_z$. We call $\delta_z$ a ``straightening map'' since it takes neighborhoods of $z$ to neighborhoods of the origin for which equivalent points 
% \sai{suspicious about "equivalent points"}
are evenly spaced apart angularly. (See \cite[Figures 2.1, 2.2]{diamond}.)

Let $\rho\colon \mathbb C\to \mathbb C$ be given by $\rho(w) = w^{h_z}$. Then $\psi = \rho\circ \delta_z \colon U\to V = \psi(U)$ given by $\psi(w) = (\delta(w))^{h_z}$ straightens then ``wraps'' around the neighborhood into a disk. By the open mapping theorem $V$ is an open subset of $\mathbb C$. There exists an bijection $\varphi\colon \pi(U)\to V$ such that $\varphi\circ\pi = \psi$ (see \cite[Section 2.2]{diamond}), which is a homeomorphism as well. The maps $\varphi$ and open sets $U$ for each $z$ are indeed charts for $Y(\varGamma)$.

For neighborhoods containing points in $\mathbb Q\cup\cbr{\infty}$, we specify charts as follows. Let $s\in\mathbb Q\cup\cbr{\infty}$ be a cusp. There exists some $\delta_s\in\slz$ that maps $s$ to $\infty$. We define the width of $s$ to be 
\[h_s = \abs{\slz_\infty/(\delta_s\cbr{\pm I}\varGamma\delta_s^{-1})_\infty}.\] Where the period of an elliptic point is the number of sectors of the disk containing the point that are identified under isotropy, the width of a cusp is the number of sectors (of the infinitely many that come together to $s$) that are distinct under isotropy (see \cite[Figure 2.6]{diamond}). The width of $s$ is finite and independent of the choice of $\delta_s$, and similarly to elliptic points, for $\gamma\in\slz$, the width of $\gamma(s)$ under $\gamma\varGamma\gamma^{-1}$ is the same as the width of $s$ under $\varGamma$. Thus $h_s$ is the same for all the cusps in $\varGamma s$, making the width well defined on $X(\varGamma)$.

Let $U = \delta_s^{-1}(\cbr{z\in\mathcal H : \Im(z) > 2}\cup\{\infty\})$ and let $\psi = \rho\circ \delta_s$ where $\rho\colon \mathbb C\to\mathbb C$ is given by $\rho(z) = e^{2\pi i z/h_s}$. With $V = \psi(U)$ (an open subset of $\mathbb C$), obtain $\psi\colon U\to V$ with $\psi(z) = e^{2\pi i \delta_s(z)/h_s}$. The effect of $\psi$ is to straighten $U$ by making identified points differ by a constant, and the exponential map wraps the upper half plane into a disk centered at $0$ (where $\infty$ maps to). Similarly, there exists a homeomorphism $\varphi\colon \pi(U)\to V$ such that $\varphi\circ\pi = \psi$ that gives the chart as desired. Hence the maps $\phi$ and open sets $U$ for each $s$ give charts on the rest of $X(\varGamma)$.

We briefly define fundamental domains obtained from $X(\varGamma)$, without proofs of claims. We first consider the case $\varGamma = \slz$. Consider the set $\mathcal D = \cbr{z\in\mathcal H : \abs{\Re(z)}\leq 1/2,\abs{z}\geq 1}$. The map $\pi\colon\mathcal D\to Y(1) = Y(\slz)$ given by the natural projection $z\mapsto \slz z$ is a surjection, and is not injective since the boundary half lines where $\Re(z) = \pm 1/2$ are identified by the translation $\big(\!\begin{smallmatrix}
    1 & 1 \\ 0 & 1
\end{smallmatrix}\!\big)$ which takes $z$ to $z+1$ (similarly, the two halves of the boundary arc where $\abs{z} = 1$ are identified by the inversion $\big(\!\begin{smallmatrix}
    0 & -1 \\ 1 & 0
\end{smallmatrix}\!\big)$ taking $z$ to $-1/z$). (See \cite[Section 2.3]{diamond}) In fact, if two distinct points in $\mathcal D$ are $\slz$-equivalent, then the points lie on the boundary of $\mathcal D$ and are either translates or inverses of each other by the aforementioned maps. So by a suitable identification of points on the boundary of $\mathcal D$ the map $\pi$ can be made a bijection. We say that $\mathcal D$ is a \boldindex{fundamental domain} for $\mathcal H$ under the action of $\slz$. A similar procedure may be done for $\mathcal D^\ast = \mathcal D \cup\cbr{\infty}$ and $X(1) = X(\slz)$, by which $\mathcal D^\ast$ becomes a fundamental domain for $\mathcal H^\ast$ under the action of $\slz$. 

For other congruence subgroups $\varGamma$, write $\slz = \bigcup_j \cbr{\pm I}\varGamma \gamma_j$ for some set of representatives $\gamma_j$ of the coset space $\varGamma\backslash \slz$, and consider the surjection $\pi\colon \bigcup_j \gamma_j\mathcal D\to Y(\varGamma)$ given by $z\mapsto \varGamma z$. By identifying the appropriate boundary points of $\bigcup_j \gamma_j\mathcal D$, $\pi$ may be regarded as a bijection, and so $\bigcup_j \gamma_j\mathcal D$ becomes a fundamental domain for $\mathcal H$ under the action of $\varGamma$. A similar procedure for $X(\varGamma)$ produces fundamental domains of $\mathcal H^\ast$ under the action of $\varGamma$.
% \sai{This last statement should be true, but maybe I should find a source}. 

\newpage\section{Modular forms}
In this section, we define modular forms and discuss the growth of their Fourier coefficients. Then we provide dimension formulas for various spaces of modular forms as well as definitions of Eisenstein series as examples of modular forms. We conclude by discussing Dirichlet characters and $L$-functions corresponding to modular forms.

\subsection{Definitions}
\begin{definition}
    For an integer $k$ and a congruence subgroup $\varGamma$, a meromorphic function $f\colon \mathcal H\to \mathbb{C}$ is \boldindex{weakly modular of weight $k$} {\textup{(}with respect to $\varGamma$\textup{)}} if \[f(\gamma(z)) = (cz+d)^kf(z)\quad\text{for any}\quad \gamma = \abcd\in \varGamma, ~z\in \mathcal H.\qedhere\]
\end{definition}
Sometimes we will say that a meromorphic function $f\colon \mathcal H\to \mathcal C$ is just ``weakly modular'' if the weight $k$ and congruence subgroup $\varGamma$ are irrelevant or clear from context.
\begin{definition}\label{def: automorphy and weight k op}
    For any matrix $\gamma=\smallabcd\in \glr$ the \boldindex{factor of automorphy} $j(\gamma,z)$ for $z\in \mathbb{C}$ is defined by \[j(\gamma,z) = cz+d.\] For any integer $k$ and $\gamma\in \glrp$ define the \boldindex{weight-$k$ $\gamma$ operator} $|_k[\gamma]$ on functions $f\colon \mathcal H\to \mathbb{C}$ by \[(f|_k[\gamma])(z) = (\det\gamma)^{k/2} j(\gamma,z)^{-k}f(\gamma(z)),\quad\text{for } z\in \mathcal H.\qedhere\]
\end{definition}
Note that $|_k[\gamma]$ composes left to right with other weight-$k$ operators.
Since the factor of automorphy $j(\gamma,z) = cz+d$ cannot be zero or infinity (as $z\in \mathcal H$), if $f$ is meromorphic on $\mathcal H$, then so is $f|_k[\gamma]$, and the number of poles and zeroes of $f|_k[\gamma]$ and $f$ are the same.

So for an integer $k$ and a congruence subgroup $\varGamma$, a meromorphic function $f\colon \mathcal H \to \mathbb{C}$ is weakly modular of weight $k$ (with respect to $\varGamma$) if \[f|_k[\gamma] = f\quad\text{for all }\gamma\in \varGamma,\] and this is equivalent to the original definition above. Note that if $f$ is weakly modular with respect to $\varGamma$, then the zeroes of $f$ and poles of $f$ are $\varGamma$-invariant sets.

\begin{lemma}\label{lem: props of automorphy}
    For any $\gamma,\eta\in \glrp$ and $z\in \mathcal H$, we have $j(\gamma\eta,z) = j(\gamma,\eta(z))j(\eta,z)$. For any function $f\colon\mathcal H\to\mathbb C$, we have $f|_k[\gamma\eta] = (f|_k[\gamma])|_k[\eta]$.
\end{lemma}
\begin{proof}
    Elements of $\glqp$ act on column vectors in $\mathbb C^2$ by matrix multiplication, and act on points in $\mathbb C$ by fractional linear transformations. Let $\gamma = \smallabcd\in \glqp$ and observe that for $z\in\mathbb C$ we have \[\gamma\begin{pmatrix}
        z \\ 1
    \end{pmatrix} = \begin{pmatrix}
        az+b \\ cz+d
    \end{pmatrix} = \begin{pmatrix}
        \gamma(z) \\ 1
    \end{pmatrix}j(\gamma,z).\]
    Then for $\eta\in\slz$ we have
    \begin{align*}
        \gamma\eta\begin{pmatrix}
            z \\ 1
        \end{pmatrix} &= \begin{pmatrix}
            (\gamma\eta)(z) \\ 1
        \end{pmatrix}j(\gamma\eta,z)\quad\text{and}\\
        \gamma\eta\begin{pmatrix}
            z \\ 1
        \end{pmatrix} &= \gamma\begin{pmatrix}
            \eta(z) \\ 1
        \end{pmatrix}j(\eta,z) = \begin{pmatrix}
            \gamma(\eta(z)) \\ 1
        \end{pmatrix}j(\gamma,\eta(z))j(\eta,z).
    \end{align*} 
    It follows that $j(\gamma\eta,z) = j(\gamma,\eta(z))j(\eta,z)$.
    Then for $f\colon \mathcal H\to \mathbb C$,
    \begin{align*}
        (f|_k[\gamma\eta])(z) &= (\det(\gamma\eta))^{k/2}j(\gamma\eta,z)^{-k}f((\gamma\eta)(z))\\
        &= (\det \eta)^{k/2}j(\eta,z)^{-k}(\det \gamma)^{k/2}j(\gamma,\eta(z))^{-k}f(\gamma(\eta(z)))\\
        &= (\det \eta)^{k/2}j(\eta,z)^{-k}(f|_k[\gamma])(\eta(z))\\
        &= ((f|_k[\gamma])|_k[\eta])(z).\qedhere
    \end{align*}
\end{proof}
From the lemma it follows that if $f\colon \mathcal H\to\mathbb C$ is weakly modular with respect to some set of matrices $A$, then $f$ is weakly modular with respect to the group generated by $A$. So with $\slz = \big\langle\big(\!\begin{smallmatrix}
    1 & 1 \\ 0 & 1
\end{smallmatrix}\!\big), \big(\!\begin{smallmatrix}
    0 & -1 \\ 1 & 0
\end{smallmatrix}\!\big)\big\rangle$, a function $f\colon \mathcal H\to \mathbb C$ is weakly modular of weight $k$ with respect to $\slz$ if
\begin{equation}\label{eqn: periodicity slz}
    f(z+1) = f\big(\big(\!\begin{smallmatrix}
        1 & 1 \\ 0 & 1
    \end{smallmatrix}\!\big)z\big) = (0z+1)^kf(z) = f(z)\quad\text{and}\quad f(-1/z) = f\big(\big(\!\begin{smallmatrix}
        0 & -1 \\ 1 & 0
    \end{smallmatrix}\!\big)z\big) = z^kf(z).
\end{equation} It follows that weakly modular functions of weight $k$ with respect to $\slz$ are $\mathbb Z$-periodic. A similar phenomenon happens for weakly modular functions of weight $k$ with respect to congruence subgroups. If $\varGamma$ is a congruence subgroup of level $N$, then $\varGamma(N)\subset \varGamma$, so that $\varGamma$ contains a translation matrix of the form $\big(\!\begin{smallmatrix}
    1 & h \\ 0 & 1
\end{smallmatrix}\!\big)$ for some minimal positive integer $h$ dividing $N$. To see that $h$ necessarily divides $N$, observe that if 
\[\begin{pmatrix}
    1 & b_1 \\ 0 & 1
\end{pmatrix},\begin{pmatrix}
    1 & b_2 \\ 0 & 1
\end{pmatrix}\in\varGamma, \quad \text{then}\quad \begin{pmatrix}
    1 & \gcd(b_1,b_2) \\ 0 & 1
\end{pmatrix}\in \varGamma.\]
By a similar computation to \eqref{eqn: periodicity slz}, it follows that weakly modular functions with respect to $\varGamma$ are $h\mathbb Z$-periodic. Weakly modular functions with respect to $\varGamma(N)$ are $N\mathbb{Z}$-periodic, and weakly modular functions with respect to $\varGamma_1(N)$ are $\mathbb Z$-periodic.

Let $D = \cbr{q\in\mathbb C : \abs{q} <1}$ be the open unit disk in $\mathbb C$ and let $D^\prime = D \setminus \!\cbr{0}$ denote the punctured open unit disk in $\mathbb C$. The exponential map $z\mapsto e^{2\pi i z/h} = q$ is a $h\mathbb{Z}$-periodic holomorphic map which maps $\mathcal H$ to $D^\prime$. Since $f$ is $h\mathbb{Z}$-periodic, it follows that the function $\tilde f\colon D^\prime\to\mathbb C$ corresponding to $f$ defined by $\tilde f(q) = f(h\log (q)/2\pi i)$ (so $f(z) = \tilde f(e^{2\pi i z/h})$) is well defined because a branch of the logarithm is determined up to integral multiples of $2\pi i$.

The logarithm can be defined holomorphically about each $q\in D^\prime$. It follows that $\tilde f$ is meromorphic on $D^\prime$ since $f$ is meromorphic on $\mathcal H$, and so $\tilde f$ has a Laurent expansion $\tilde f(q) = \sum_{n\in\mathbb Z}a_nq^n$ at each $q$ in a punctured neighborhood of $q=0$ (and $a_n = 0$ for all $n$ sufficiently small). Moreover, if $f$ is holomorphic on $\mathcal H$, it follows that $\tilde f$ is also holomorphic on $D^\prime$.

From $\abs{q} = e^{-2\pi \Im z/h}$, it follows that $q$ tends to zero as $\Im z$ tends to infinity. We define $f$ to be meromorphic (respectively holomorphic) at $\infty$ if $\tilde f$ has a meromorphic (respectively holomorphic) extension to $q = 0$. The Laurent series of $\tilde f$ about $q=0$ is used to obtain a Fourier series expansion of $f$ about $\infty$, given by 
\[f(z) = \sum_{n\in\mathbb Z} a_nq^n,\quad q = e^{2\pi i z/h},\] which converges absolutely and uniformly on compact subsets of the half plane $\cbr{z\in\mathbb Z: \Im z>\tau}$ for some large enough $\tau$ (so that $q$ lies in a punctured neighborhood of zero). When referring to a Fourier series of a $h\mathbb Z$-periodic meromorphic function on $\mathcal H$, we mean the expansion obtained in the above manner. If $f$ is holomorphic on $\mathcal H$ and is holomorphic at $\infty$, the Fourier series expansion becomes $f(z) = \sum_{n=0}^\infty a_nq^n$, $q = e^{2\pi i z/h}$, valid for $z\in \mathcal H$. To see that a weakly modular holomorphic function $f\colon \mathcal H\to \mathbb C$ is holomorphic at $\infty$, it suffices to show that $\lim_{\Im z\to\infty} f(z)$ exists or that $f(z)$ is bounded as $\Im z$ grows unboundedly.
\begin{definition}\label{def: modular form}
    Let $\varGamma$ be a congruence subgroup of $\slz$ and let $k$ be an integer. A function $f\colon \mathcal H\to\mathbb C$ is a \boldindex{modular form of weight $k$} \textup{(}with respect to $\varGamma$\textup{)} if it satisfies the following: \begin{enumerate}[label=(\arabic*)]
        \item $f$ is holomorphic on $\mathcal H$,
        \item $f$ is weakly modular of weight $k$ with respect to $\varGamma$, and
        \item $f|_{k}[\alpha]$ is holomorphic at $\infty$ for all $\alpha\in\slz$.
    \end{enumerate}Furthermore, if for every $\alpha\in\slz$, the coefficient $a_0$ vanishes in the Fourier series expansion of $f|_k[\alpha]$, then we call $f$ a \textib{cusp form of weight $k$} \textup{(}with respect to $\varGamma$\textup{)}. The set of modular forms (respectively cusp forms) of weight $k$ with respect to $\varGamma$ is denoted by \boldindex{$\mathcal M_k(\varGamma)$} (respectively \boldindex{$\mathcal S_k(\varGamma)$}).

    Define also the \textib{weight $k$ Eisenstein space} \textup{(}of $\varGamma$\textup{)} by the quotient of the space of modular forms by the space of cusp forms; that is, \boldindex{$\mathcal E_k(\varGamma)$} $= \mathcal M_k(\varGamma)/\mathcal S_k(\varGamma)$. It can be shown that the Eisenstein space may be obtained as a complement of $\mathcal S_k(\varGamma)$ in $\mathcal M_k(\varGamma)$ under the Petersson inner product (we define the Petersson inner product in \cref{section petersson inner product}, but do not prove this result). We briefly discuss the Eisenstein space at the end of \cref{section eisenstein series}.
\end{definition}

Later on, for example in \cref{section primitive forms}, we specify the congruence subgroup that a modular form is with respect to by its level; that is, we say that a modular form with respect to $\varGamma$ is ``at level $N$'' if $\varGamma$ is of level $N$.

Recall that the $\varGamma$-equivalence classes of points in $\mathbb Q\cup \infty$ are the cusps of $\varGamma$. There are finitely many cusps, at most the index of $\varGamma$ in $\slz$, which we showed was finite in the discussion following \cref{def: congruence subgroup}. Write a cusp $s$ as $\alpha(\infty)$ for some $\alpha\in\slz$.
We say that $f$ is holomorphic at $s$ $f|_k[\alpha]$ is holomorphic at $\infty$, where $f|_k[\alpha]$ is viewed as a modular form of weight $k$ with respect to $\alpha^{-1}\varGamma\alpha$. In view of this definition, condition (3) of \cref{def: modular form} says that $f$ must be holomorphic at the cusps of $\varGamma$.
The group $\alpha^{-1}\varGamma\alpha$ is a congruence subgroup since for some $N>0$, the principal congruence subgroup $\varGamma(N)$ is contained in $\varGamma$ and is normal in $\slz$, so that $\varGamma(N) = \alpha^{-1}\varGamma(N)\alpha\subset \alpha^{-1}\varGamma\alpha$. 

As congruence subgroups have finite index in $\slz$, only finitely many coset representatives $\alpha_j$ in a decomposition $\slz= \bigcup_{j}\varGamma\alpha_j$ are needed to verify condition (3) in the definition above, and in verifying that the term $a_0$ vanishes for all Fourier series expansions for cusp forms: we have by condition (2) of the above definition that $f|_k[\gamma\alpha_j] = f|_k[\alpha_j]$ for all $\gamma\in\varGamma$.

\begin{lemma}\label{lem: weight k operator properties}
    \begin{enumerate}[label = \textup{(\alph*)}]
        \item For any $\gamma\in\glqp$, there exists $\alpha\in\slz$, $r\in\mathbb Q^+$, and $a,b,d\in\mathbb Z$ relatively prime such that $\gamma = r\alpha \big(\!\begin{smallmatrix}
            a & b \\ 0 & d
        \end{smallmatrix}\!\big)$. 
    
        It follows that for $f\in\mathcal M_k(\varGamma)$ and $\gamma$ with $\gamma = r\alpha \big(\!\begin{smallmatrix}
            a & b \\ 0 & d
        \end{smallmatrix}\!\big)$, since $f|_k[\alpha]$ has a Fourier expansion, so does $f|_k[\gamma]$. Moreover, if the constant term in the Fourier expansion for $f|_k[\alpha]$ is $0$, the same holds for the Fourier expansion for $f|_k[\gamma]$.
        \item Let $\varGamma_1,\varGamma_2$ be congruence subgroups with $\gamma\varGamma_2\gamma^{-1}\subset \varGamma_1$ for some $\gamma\in\glqp$. If $f\in\mathcal M_k(\varGamma_1)$, then $f|_k[\gamma]\in\mathcal M_k(\varGamma_2)$, and the same result holds for cusp forms.
        % \item Let $N,d$ be positive integers and let $f\in\mathcal M_k(\varGamma_i(N))$. If $g(z) = f(dz)$, then $g\in\mathcal M_k(\varGamma_i(dN))$ for $i = 0,1$, and the same result holds if $f$ is a cusp form \textup{(}i.e., $g$ is also a cusp form\textup{)}.
    \end{enumerate}  
\end{lemma}
\begin{proof}
    Let $\gamma = \big(\!\begin{smallmatrix}
        a & b \\ c & d
    \end{smallmatrix}\!\big)$ be an element of $\glqp$. If $c = 0$, there is nothing to show. So assume that $c\neq 0$. Then write $a/c = a^\prime/c^\prime$ with $a^\prime$ and $c^\prime$ coprime integers; that is, write $a/c$ in lowest terms. Then choose integers $s,t$ such that $sa^\prime  + tc^\prime = 1$, from which it follows that $\big(\!\begin{smallmatrix}
        s & t \\ -c^\prime & a^\prime
    \end{smallmatrix}\!\big)\in\slz$. We have 
    \[\begin{pmatrix}
        s & t \\ -c^\prime & a^\prime
    \end{pmatrix}\gamma = c\begin{pmatrix}
        s & t \\ -c^\prime & a^\prime
    \end{pmatrix}\begin{pmatrix}
        a^\prime/c^\prime & b/c \\ 1 & d/c
    \end{pmatrix} = \begin{pmatrix}
        sa + tc & sb + tc \\ 0 & a^\prime d - c^\prime b
    \end{pmatrix} = A\]
    There exists an integer $r>0$ such that $rA \in \Mat_2(\mathbb Z)$, from which we can divide through by the greatest common divisor $g$ of the (three) nonzero entries of $rA$ to find that the nonzero entries of $\frac{r}{g}A$ are coprime to each other. Rearranging, we obtain $\gamma = \frac{g}{r}\big(\!\begin{smallmatrix}
        s & t \\ -c^\prime & a^\prime
    \end{smallmatrix}\!\big)^{-1}\big(\frac{r}{g}A\big)$, which proves the first part of (a).

    %Before showing the second part of (a), note that $f|_k[\alpha]$ has a Fourier expansion about $\infty$.

    To see the second part of (a), observe that if $\gamma=r\alpha\big(\!\begin{smallmatrix}
        a & b \\ 0 & d
    \end{smallmatrix}\!\big)$, then $f|_k[\gamma] = f|_k[r\alpha\big(\!\begin{smallmatrix}
        a & b \\ 0 & d
    \end{smallmatrix}\!\big)] = (f|_k[\alpha])|_k[\big(\!\begin{smallmatrix}
        a & b \\ 0 & d
    \end{smallmatrix}\!\big)]$. Since $f$ is a modular form, $f|_k[\alpha]$ is holomorphic at $\infty$ and hence admits a Fourier series expansion of the form $(f|_k[\alpha])(z) = \sum_{n=0}^\infty a_ne^{2\pi i n z / h}$ for some positive integer $h$. Thus the Fourier series expansion for $(f|_k[\alpha])|_k[\big(\!\begin{smallmatrix}
        a & b \\ 0 & d
    \end{smallmatrix}\!\big)]$ is of the form $(f|_k[\alpha\big(\!\begin{smallmatrix}
        a & b \\ 0 & d
    \end{smallmatrix}\!\big)])(z) = \sum_{n=0}^\infty b_ne^{2\pi i az/dh}$. Observe that $b_0$ is proportional to $a_0$.

    Part (b) is clear from the following observation. For any $\alpha\in\slz$ we have that $\gamma\alpha\gamma^{-1}\in\slz$ and
    \[(f|_k[\gamma])|_k[\alpha] = (f|_k[\gamma\alpha\gamma^{-1}])|_k[\gamma].\] 
    If $\alpha\in\varGamma_2$, then $\gamma\alpha\gamma^{-1}\in\varGamma_1$ so that $(f|_k[\gamma\alpha\gamma^{-1}])|_k[\gamma] = f|_k[\gamma]$. From part (a) deduce the holomorphicity of $(f|_k[\gamma\alpha\gamma^{-1}])|_k[\gamma]$ for any $\alpha\in\slz$, and that the result in part (b) is true for cusp forms.
\end{proof}
It is possible to determine if a weakly modular holomorphic function on $\mathcal H$ is a modular form by investigating the growth of its Fourier coefficients. First, we prove a few lemmas.
\begin{lemma}
    If $x$ is a cusp of a congruence subgroup $\varGamma$ and $\sigma\in\slz$ sends $x$ to $\infty$, then 
    \[\sigma\varGamma_x\sigma^{-1}\cdot\cbr{\pm I} = \Big\{\pm\begin{pmatrix}
        1 & h \\ 0 & 1
    \end{pmatrix}^m : m\in\mathbb Z\Big\}\] for some integer $h>0$.
\end{lemma}
\begin{proof}
    Without loss of generality take $x = \infty$ by taking $\sigma\varGamma\sigma^{-1}$ in place of $\varGamma$. By a computation, we have
    \[\varGamma_\infty \subset \cbr{\pm\begin{pmatrix}
        1 & b \\ 0 & 1
    \end{pmatrix} : b\in\mathbb Z}\]
%     There exists an element  $\gamma = \big(\!\begin{smallmatrix}
%     1 & \ell \\ 0 & 1
% \end{smallmatrix}\!\big)$ with $\ell\neq 0$ in the stabilizer $\varGamma_\infty$ of $\infty$. 
%  Suppose there exists an element $\alpha$ of $\varGamma_\infty$ of the form $\big(\!\begin{smallmatrix}
%     a & b \\ 0 & a^{-1}
% \end{smallmatrix}\!\big)$ with $a\neq \pm 1$. Assume $\abs{a}<1$ by taking $\alpha^{-1}$ in place of $\alpha$ if needed. Then for any positive integer $n$, the element 
% \[\alpha^n\gamma\alpha^{-n} = \begin{pmatrix}
%     1 & a^{2n}l \\ 0 & 1
% \end{pmatrix}\in \varGamma.\] We take as a fact that congruence subgroups are discrete topological groups, so that the previous statement yields a contradiction. It follows then that $\varGamma_\infty\subset\cbr{\pm\big(\!\begin{smallmatrix}
%     1 & b \\ 0 & 1
% \end{smallmatrix}\!\big) : b\in\mathbb R}$.
We obtain the result by observing that $\cbr{\pm\big(\!\begin{smallmatrix}
    1 & b \\ 0 & 1
\end{smallmatrix}\!\big) : b\in\mathbb Z}/\cbr{\pm I}\cong \mathbb Z$.
\end{proof}
\begin{lemma}\label{lem: bound growth fourier coeffs}
    Let $f$ be a weakly modular holomorphic function on $\mathcal H$ of weight $k$ with respect to some congruence subgroup $\varGamma$. If there exists a real number $v$ such that $f(z) = O(\Im(z)^{-v})$ as $\Im(z)\to\infty$ uniformly with respect to $\Re(z)$, then $f$ is a modular form of weight $k$. Moreover, if $v$ may be chosen so that $v<k$, then $f$ is a cusp form.
\end{lemma}
\begin{proof}
    % We may assume that $k$ is even by squaring $f$. 
    % \sai{(why do we need $k$ even?)}
    Let $x$ be a cusp of $\varGamma$. If $x = \infty$, then we may send it to a rational cusp via an element of $\varGamma$ since $\varGamma\neq \varGamma_\infty$, and repeat the above argument. So suppose $x$ is rational. Then there exists $\sigma\in\slz$ such that $\sigma(x) = \infty$. Let $h>0$ be an integer satisfying the result of the previous lemma. Then the Fourier series expansion of $f|_k[\sigma^{-1}]$ is 
    \[(f|_k[\sigma^{-1}])(z) = \sum_{n=-\infty}^\infty a_ne^{2\pi i n z/h},\] with \begin{equation}\label{eqn: fourier coeffs z_0}
        a_n = \frac{1}{h}\int_{z_0}^{z_0+h}(f|_k[\sigma^{-1}])(z)e^{-2\pi i n z/h}\dd z,\quad \text{for any fixed $z_0\in\mathcal H$.}
    \end{equation} Let $\smallabcd = \sigma^{-1}$, and note that $c\neq 0$, since $\sigma^{-1}$ does not fix $\infty$. We have that $\Im(\sigma^{-1}(z)) = \Im(z)/\abs{cz+d}^2 = O(1/\Im(z))$ as $\Im(z)\to\infty$, uniformly on $\abs{\Re(z)}\leq h/2$. Then by assumption,
    \[(f|_k[\sigma^{-1}])(z) = f(\sigma^{-1}(z))j(\sigma^{-1},z)^{-k} = O(\Im(z)^{v-k})\quad (\text{as }\Im(z)\to\infty),\] uniformly on $\abs{\Re(z)}\leq h/2$. Choosing $z_0 = iy-h/2$ in \eqref{eqn: fourier coeffs z_0}, we have that $|a_n| = O(y^{v-k}e^{2\pi n y /h})$ as $y\to\infty$.

    So if $n<0$, then $a_n = 0$; moreover, if $v<k$, then $a_0 = 0$. It follows that $f$ is holomorphic at all cusps of $\varGamma$, and has a zero at any cusp if $v<k$.
\end{proof}
\begin{lemma}\label{lem: miyake lem 4.3.3}
    Let $\{a_n\}_{n=0}^\infty$ be a sequence of complex numbers, and let $f\colon \mathcal H\to \mathbb C$ be given by $f(z) = \sum_{n=0}^\infty a_ne^{2\pi i nz}$. If $a_n = O(n^v)$ for some $v >0$, then $\sum_{n=0}^\infty a_ne^{2\pi i nz}$ is convergent absolutely and uniformly on compact subsets of $\mathcal H$. Moreover, 
    \begin{align*}
        f(z) &= O(\Im(z)^{-v-1}) \quad\textup{($\Im(z) \to 0$), and }\\
        f(z) - a_0 &= O(e^{-2\pi \Im(z)}) \quad\textup{($\Im(z) \to \infty$)},
    \end{align*}
    uniformly on $\Re(z)$.
\end{lemma}
\begin{proof}
    From $\varGamma(s) = \lim_{n\to \infty}\frac{n!n^s}{s(s+1)\cdots(s+n)}$ for real $s>0$ (Euler-Gauss), we have for $v>0$ that
    \[\lim_{n\to\infty} n^v/(-1)^n\binom{-v-1}{n} = \varGamma(v).\]
    Thus there exists $L>0$ such that 
    \[\abs{a_n}\leq L(-1)^n\binom{-v-1}{n}\] for all $n\geq 0$. Let $z = x+iy$, so that 
    \begin{align}\label{eqn: miyake eqn 4.3.9}
        \sum_{n=0}^\infty\abs{a_n}|e^{2\pi i n z}|&\leq L\bigg(\sum_{n=0}^\infty (-1)^n\binom{-v-1}{n}e^{-2\pi n y}\bigg)\\
        &= L(1-e^{-2\pi y})^{-v-1}.\nonumber
    \end{align}
    It follows that $f$ is convergent absolutely and uniformly on compact subsets of $\mathcal H$. Since $1-e^{-2\pi y} = O(y)$ as $y\to 0$, we have that $\abs{f(z)} = O(y^{-v-1})$. Furthermore, \eqref{eqn: miyake eqn 4.3.9} implies that $f$ is bounded as $y$ tends to $\infty$.

    Let $g\colon\mathcal H\to\mathbb C$ be given by $g(z) = \sum_{n=0}^\infty a_{n+1}e^{2\pi i nz}$, and observe that $g$ also satisfies the hypotheses of the lemma. It follows that $g$ is also bounded on a neighborhood of $\infty$. Therefore 
    \[f(z)-a_0 = e^{2\pi i z}g(z) = O(e^{-2\pi y})\quad (y\to \infty).\qedhere\]
\end{proof}
Combining the previous few lemmas, we obtain the following:
\begin{proposition}\label{prop: growth of fourier coeffs}
    Let $f$ be a weakly modular holomorphic function on $\mathcal H$ with respect to some congruence subgroup $\varGamma$. If in a Fourier expansion of $f$ given by $f(z) = \sum_{n = 0}^\infty a_ne^{2\pi i n z/h}$, the coefficients $a_n$ are of order $O(n^v)$ for some $v>0$, then $f$ is a modular form.
\end{proposition}
The following result describes the growth of cusp forms:
\begin{proposition}\label{prop: product cuspform boundedness}
    Let $f$ be a weakly modular holomorphic function on $\mathcal H$. Then $f$ is a cusp form if and only if $f(z)\Im(z)^{k/2}$ is bounded on $\mathcal H$.
\end{proposition}
\begin{proof}
    We may assume that $k$ is even, and observe that \cref{lem: bound growth fourier coeffs} provides the `if' direction of the result. So assume that $f$ is a cusp form and let $g$ be given by $g(z) = |f(z)|\Im(z)^{k/2}$. Since $g(\gamma(z)) = g(z)$ for any $\gamma\in\varGamma$, view $g$ as a continuous function on $Y(\varGamma)$. Since $\varGamma$ has only finitely many inequivalent cusps, it suffices to show that $g$ is bounded on a neighborhood of a cusp of $\varGamma$. Let $x$ be a cusp of $\varGamma$, and let $\sigma\in\slz$ take $x$ to $\infty$. Choose a positive integer $h$ so that $\sigma\varGamma_x\sigma^{-1}\cdot\cbr{\pm I} = \cbr{\pm\big(\!\begin{smallmatrix}
        1 & h \\ 0 & 1
    \end{smallmatrix}\!\big)^m : m\in\mathbb Z}$, and let $(f|_k[\sigma^{-1}])(z) = \sum_{n=1}^\infty a_ne^{2\pi i nz /h}$ be the Fourier expansion of $f$ at $x$. Then \begin{align*}
        g(\sigma^{-1}(z)) &= \abs{(f|_k[\sigma^{-1}])(z)}\Im(z)^{k/2}\\
        &= \abs{\sum_{n=1}^\infty a_ne^{2\pi i n z/h}}\Im(z)^{k/2}\quad (\text{as }\Im(z)\to\infty).
    \end{align*} Hence $g$ is bounded on a neighborhood of $x$.
\end{proof}
\begin{corollary}\label{cor: cuspform growth corollary}
    Let $f$ be an element of $\mathcal S_k(\varGamma)$, let $x$ be a cusp of $\varGamma$, and let $\sigma$ be an element of $\slz$ that maps $x$ to $\infty$. Let $(f|_k[\sigma^{-1}])(z) = \sum_{n=1}^\infty a_ne^{2\pi i nz /h}$ be the Fourier expansion of $f$ at $x$. Then $a_n = O(n^{k/2})$.
\end{corollary}
\begin{proof}
    Let $g = f|_k[\sigma^{-1}]$, so that $g\in\mathcal S_k(\sigma\varGamma\sigma^{-1})$. By \cref{prop: product cuspform boundedness}, there exists a constant $M>0$ such that $\abs{g(z)}\leq M\Im(z)^{-k/2}$. Thus 
    \begin{align*}
        |a_n| &= \frac{1}{h}\abs{\int_0^h g(x+iy)e^{-2\pi i n(x+iy)/h}\dd x}\\
        &\leq My^{-k/2}e^{2\pi n y/h}.
    \end{align*}
    By taking $y = 2/n$, it follows that $\abs{a_n}\leq (Me^{4\pi /h}2^{-k/2})n^{k/2}$.
\end{proof}
Lastly, we record the dimensions of various spaces of modular forms below, without proof (See \cite[Chapter 3]{diamond}). These dimension formulas are obtained via Riemann surface theory. In particular, the Riemann-Roch theorem is used, in combination with the geometric data of $X(\varGamma)$: its genus, number of cusps, number of elliptic points, and more.

\begin{theorem}
    Let $k$ be an even integer and let $\varGamma$ be a congruence subgroup. Let $g$ be the genus of $X(\varGamma)$, let $\varepsilon_2,\varepsilon_3$ be the number of elliptic points with period $2$, $3$ respectively, and let $\varepsilon_\infty$ be the number of cusps. Then

    \[\dim \mathcal M_k(\varGamma) = \begin{cases}
        (k-1)(g-1) + \lfloor\frac{k}{4}\rfloor\varepsilon_2 + \lfloor\frac{k}{3}\rfloor\varepsilon_3 + \frac{k}{2}\varepsilon_\infty & \text{if $k\geq 2$},\\
        1 & \text{if $k = 0$},\\
        0 & \text{if $k<0$},
    \end{cases} \quad \text{and}\]

    \[\dim \mathcal S_k(\varGamma) = \begin{cases}
        (k-1)(g-1) + \lfloor\frac{k}{4}\rfloor\varepsilon_2 + \lfloor\frac{k}{3}\rfloor\varepsilon_3 + (\frac{k}{2}-1)\varepsilon_\infty & \text{if $k\geq 4$},\\
        g & \text{if $k = 2$},\\
        0 & \text{if $k<0$},
    \end{cases}\]
\end{theorem}

In particular, the space of modular forms of weight $0$, $\mathcal M_0(\slz)$, is isomorphic to $\mathbb C$ since the genus of $X(1)$ is $0$. Furthermore, $\mathcal M_2(\slz) = 0$ and $\mathcal S_k(\slz) = 0$ for $k = 0,2$. For any even integer $k\geq 4$, $\mathcal M_k(\slz) = \mathcal S_k(\slz)\oplus \mathbb CE_k$ where $E_k$ is the normalized weight $k$ Eisenstein series, which we will define in \cref{section eisenstein series}, and \[\dim S_k(\slz) = \begin{cases}
    \lfloor \frac{k}{12}\rfloor - 1 & \text{if $k\equiv 2 \smod{12}$}, \\
    \lfloor \frac{k}{12}\rfloor & \text{otherwise.}
\end{cases}\]

\begin{theorem}
    Let $k$ be an odd integer and let $\varGamma$ be a congruence subgroup. If $-I\in\varGamma$, then both $\mathcal M_k(\varGamma)$ and $\mathcal S_k(\varGamma)$ are zero. Otherwise, let $g$ be the genus of $X(\varGamma)$, $\varepsilon_3$ be the number of elliptic points with period $3$, $\varepsilon_\infty^{\mathrm{reg}}, \varepsilon_\infty^{\mathrm{irr}}$ be the number of regular and irregular cusps, respectively \textup{(}see \cite[Section 3.6]{diamond}\textup{)}. Then 
    \[\dim \mathcal M_k(\varGamma) = \begin{cases}
        (k-1)(g-1) + \lfloor\frac{k}{3}\rfloor\varepsilon_3 + \frac{k}{2}\varepsilon_\infty^{\mathrm{reg}} + \frac{k-1}{2}\varepsilon_\infty^{\mathrm{irr}} & \text{if $k\geq 3$},\\
        0 & \text{if $k<0$},
    \end{cases}\quad \text{and}\]
    \[\dim \mathcal S_k(\varGamma) = \begin{cases}
        (k-1)(g-1) + \lfloor\frac{k}{3}\rfloor\varepsilon_3 + (\frac{k}{2}-1)\varepsilon_\infty^{\mathrm{reg}} + \frac{k-1}{2}\varepsilon_\infty^{\mathrm{irr}} & \text{if $k\geq 3$},\\
        0 & \text{if $k<0$}.
    \end{cases}\] If $\varepsilon_\infty^{\mathrm{reg}}>2g-2$, then $\dim \mathcal M_1(\varGamma) = \varepsilon_\infty^{\mathrm{reg}}/2$ and $\dim \mathcal S_1(\varGamma) = 0$. If $\varepsilon_\infty^{\mathrm{reg}}\leq 2g-2$, then $\dim \mathcal M_1(\varGamma)\geq \varepsilon_\infty^{\mathrm{reg}}$ and $\dim \mathcal S_1(\varGamma) = \dim \mathcal M_1(\varGamma)- \varepsilon_\infty^{\mathrm{reg}}/2$.
\end{theorem}

Dimension formulas for other spaces of modular forms (varying $k, N, \varGamma$) may be found in \cite[Section 3.9]{diamond}.

Let $k$ be an integer and $\varGamma$ a congruence subgroup. Then the dimensions of the Eisenstein spaces are given by 
\[\dim \mathcal E_k(\varGamma) = \begin{cases}
    \varepsilon_\infty & \text{if $k\geq 4$ is even},\\
    \varepsilon_\infty^{\mathrm{reg}} & \text{if $k\geq 3$ is odd and $-I\not\in\varGamma$},\\
    \varepsilon_\infty-1 & \text{if $k=2$},\\
    \varepsilon_\infty^{\mathrm{reg}}/2 & \text{if $k=1$ and $-I\not\in\varGamma$},\\
    1 & \text{if $k=0$},\\
    0 & \text{if $k<0$, or if $k>0$ is odd and $-I\in\varGamma$}.
\end{cases}\]

\subsection{Eisenstein series}\label{section eisenstein series}
We define the Eisenstein series for a number of spaces of modular forms to provide some examples of modular forms. We do not verify any claims in this section.

Elementary Eisenstein series for $\slz$ are defined for even weights $k\geq 4$, by \[G_k(z) = {\sum_{(c,d)}}^\prime \frac{1}{(cz+d)^k},\quad z\in\mathcal H,\] where the notation $\sum_{(c,d)}^\prime$ denotes summation over all nonzero integer pairs $(c,d)\in\mathbb Z^2$. This sum is absolutely convergent, and converges uniformly on compact subsets of $\mathcal H$. It follows that $G_k$ is holomorphic on $\mathcal H$ and we may rearrange the terms of the sum defining $G_k$. A Fourier series for $G_k$ is given by \[G_k(z) = 2\zeta(k) + 2\frac{(2\pi i)^k}{(k-1)!}\sum_{n=1}^\infty \sigma_{k-1}(n)e^{2\pi i z},\] where $\sigma_{k-1}$ is the arithmetic function $\sigma_{k-1}(n) = \sum_{\substack{m\mid n\\ m>0}}m^{k-1}$. (See \cite[Section 1.1]{diamond}) Obtain the normalized Eisenstein series $E_k(z) = G_k(z)/(2\zeta(k))$, which can be shown to be given by 
\[E_k(z) = \frac{1}{2}\sum_{\substack{(c,d)\in\mathbb Z^2\\\gcd(c,d) = 1}}\frac{1}{(cz+d)^k}\quad \text{or} \quad E_k(z) = \frac{1}{2}\sum_{\gamma\in P_+\backslash \slz}\frac{1}{j(\gamma,z)^k},\] where $P_+ = \cbr{\big(\!\begin{smallmatrix}
    1 & n \\ 0 & 1
\end{smallmatrix}\!\big) : n\in\mathbb Z}$ is the positive part of the parabolic subgroup of $\slz$. It is more elegant to verify that the normalized Eisenstein series is weakly modular of weight $k$ using the last equality above. 

Let $g_2(z) = 60 G_4(z)$ and $g_3(z) = 140G_6(z)$, and define the \textib{discriminant function} $\varDelta\colon\mathcal H\to\mathbb C$ by $\varDelta(z) = (g_2(z))^3-27(g_3(z))^2$. The discriminant function is a nonzero cusp form of weight $12$, with Fourier coefficients $a_0 = 0$, $a_1 = (2\pi)^{12}$. In fact, the only zero of $\varDelta$ is at $\infty$. It follows that the \textib{$j$-invariant} $j\colon \mathcal H\to\mathbb C$ given by $j(z) = 1728(g_2(z))^3/\varDelta(z)$ is holomorphic on $\mathcal H$. The $j$-invariant is $\slz$-invariant, is a surjection, and has a simple pole at $\infty$ with residue $1$.

We summarize what the Eisenstein series for $\varGamma(N)$ look like for weights $k\geq 3$. Let $N$ be a positive integer and let $\overline v \in (\mathbb Z/N\mathbb Z)^2$ be an element of order $2$ written as a row vector, where $v$ is some lift of $\overline v$ to $\mathbb Z^2$ (so $\overline{\,\cdot\,}$ denotes reduction modulo $N$). Let $\delta = \big(\!\begin{smallmatrix}
    a & b \\ c_v & d_v
\end{smallmatrix}\!\big)\in\slz$ with $(c_v,d_v)$ a lift of $\overline v$ to $\mathbb Z^2$, and let $k\geq 3$ be an integer. Lastly, let $\epsilon_N$ be $1/2$ if $N=1,2$ and $1$ otherwise. 

Define the Eisenstein series $E_k^{\overline v}(z)\colon\mathcal H\to\mathbb C$ by \[E_k^{\overline v}(z) =\epsilon_N\sum_{\substack{(c,d)\equiv v\smod N\\\gcd(c,d) = 1}}\frac{1}{(cz+d)^k},\] and it can be shown that \[E_k^{\overline v}(z) = \sum_{\gamma \in(P_+\cap \varGamma(N)\backslash \varGamma(N)\delta)}\frac{1}{j(\gamma,z)^k}.\] Note that when $N = 1$, we have from this definition that $E_k^{\overline v} = E_k$ since there is only one choice of $\overline v$.

A computation yields the fact that for $\gamma\in\slz$, $E_k^{\overline v}|_k[\gamma] = E_k^{\overline{v\gamma}}$, from which it follows that these series are elements of $\mathcal M_k(\varGamma(N))$. By symmetrizing, one can define Eisenstein series for any congruence subgroup $\varGamma$ by
\[E_{k,\varGamma}^{\overline v} = \sum_{\gamma_j\in\varGamma(N)\backslash \varGamma}E_k^{\overline v}|_k[\gamma_j],\] where the $\gamma_j$ constitute a set of coset representatives for $\varGamma(N)\backslash \varGamma$.

For odd weights $k$ and $N = 1,2$, the Eisenstein space $\mathcal E_k(\varGamma(N))$ has dimension $0$, since $-I\in\varGamma(N)$. For $k$ even or for $N>2$, $E_k^{\overline v}$ vanishes at $\infty$ for all $\overline v$ except for $\pm\overline{(0,1)}$. One can show that $E_k^{\overline{(c,d)}}$ is nonvanishing on points in the set $\varGamma(N)(-d/c)$ and vanishes on the other cusps of $\varGamma(N)$. A basis of $\mathcal E_k(\varGamma(N))$ may be formed by choosing a set of vectors $\cbr{\overline v} = \{\overline{(c,d)}\}$ for which the quotients $-d/c$ represent the cusps of $\varGamma(N)$; it follows from the previous sentence that $\{E_k^{\overline v}\}$ is linearly independent. The number of elements in this set is $\varepsilon_\infty$ (for $\varGamma(N)$). Thus for all $k\geq 3$, it is possible to obtain bases for the Eisenstein spaces (recall that for odd weights and $N = 1,2$ the spaces are zero). The basis elements are given by the reductions modulo $\mathcal S_k(\varGamma(N))$ of the modular forms $E_k^{\overline v}$, but it is possible to redefine the Eisenstein spaces $\mathcal E_k(\varGamma(N))$ as subspaces of $\mathcal M_k(\varGamma(N))$ so that the basis elements are the Eisenstein series themselves (see \cite[Section 5.11]{diamond}).

In view of the definition of the normalized Eisenstein series for $N = 1$, define for any $\overline v\in (\mathbb Z/N\mathbb Z)^2$ of order $N$
\[G_k^{\overline v}(z) = \sideset{}{^\prime}\sum_{(c,d)\equiv v\smod N}\frac{1}{(cz+d)^k}.\] It can be shown that (see \cite[Section 4.2]{diamond})
\[G_k^{\overline v}(z) = \frac{1}{\epsilon_N}\sum_{n\in(\mathbb Z/N\mathbb Z)^\times}\zeta_+^n(k)E_k^{n^{-1}\overline v}(z),\] where \[\zeta_+^n(k) = \sum_{\substack{m=1\\m\equiv n\smod N}}^\infty \frac{1}{m^k}\quad \text{for $n\in (\mathbb Z/N\mathbb Z)^\times$}.\] 

The Fourier series expansion of $G_k^{\overline v}(z)$ for $k\geq 3$ and $\overline v\in (\mathbb Z/N\mathbb Z)^2$ of order $N$ is given by \[G_k^{\overline v}(z) = \delta(\overline c_v)\zeta^{\overline d_v}(k) + \frac{(-2\pi i)^k}{N^k(k-1)!}\sum_{n=1}^\infty \sigma_{k-1}^{\overline v}(n)e^{2\pi i n t/N},\] where \[\delta(\overline c_v) = \begin{cases}
    1 & \text{if $\overline c_v =\overline 0$},\\
    0 & \text{otherwise},
\end{cases}\quad \zeta^{\overline d_v}(k) = \sideset{}{^\prime}\sum_{d\equiv d_v\smod N}\frac{1}{d^k},\] and 
\[\sigma_{k-1}^{\overline v}(n) = \sum_{\substack{m\mid n\\n/m \equiv c_v\smod n}}\sgn(m)m^{k-1}\mu_N^{d_vm}.\] In the sum for $\zeta^{\overline d_v}(k)$, we sum over positive and negative $d$, and similarly for $m$ in the sum for $\sigma_{k-1}^{\overline v}$. Similarly to the Eisenstein series, one may form a basis $\cbr{G_k^{\overline v}}$ of the Eisenstein space $\mathcal E_k(\varGamma(N))$ (using the same choice of $\cbr{\overline v}$ as before). Since the $n$-th Fourier coefficient is of the order $n^k$, it follows from \cref{prop: growth of fourier coeffs} that $E_k^{\overline v}$ is a modular form.

\subsection{Dirichlet characters and $L$-functions}
For any positive integer $N$, denote by $Z_N$ the group $\mathbb{Z}/N\mathbb{Z}$.
\begin{definition}
    A \boldindex{Dirichlet character} \textib{modulo $N$} is a homomorphism 
    \[\chi\colon Z_N^\times\to \mathbb{C}^\times.\] (We will sometimes suppress ``Dirichlet'' when referring to Dirichlet characters.)
\end{definition}
The product of two Dirichlet characters $\chi,\psi$, defined by pointwise multiplication $(\chi\psi)(n) = \chi(n)\psi(n)$ for $n\in Z_N^\times$, is a Dirichlet character. The trivial map is a Dirichlet character, called the trivial character. Hence the set of Dirichlet characters of $Z_N^\times$ forms a group called the \textib{dual group} of $Z_N^\times$, denoted $\widehat{Z_N^\times}$. Since $Z_N^\times$ is a finite group, the image of any character lies in the roots of unity. Thus the inverse of a Dirichlet character is its complex conjugate character $\overline \chi$, defined by taking the complex conjugate pointwise: $\overline\chi(n) = \overline{\chi(n)}$ for $n\in Z_N^\times$, where $\overline{\,\cdot\,}$ denotes complex conjugation. Note that the only Dirichlet character of $Z_1^\times$ is the trivial character $\mathbf{1}_1$. 

\begin{proposition}
    The dual group $\widehat{Z_N^\times}$ is (noncanonically) isomorphic to $Z_N^\times$; it follows that the number of Dirichlet characters modulo $N$ is $\phi(N)$.
\end{proposition}

\begin{proof}
    This is \cite[Exercise 5.2.14]{df}, which states that finite Abelian groups are (noncanonically) self-dual. 

    Let $G = \abr{x_1}\times\cdots\times\abr{x_r}$ be a finite Abelian group (recall the structure theorem for finitely generated Abelian groups). Define characters $\chi_i\colon \abr{x_i}\to \mathbb C^\times$ by $x_i\mapsto e^{2\pi i/\abs{x_i}}$ for $1\leq i\leq r$, which have order $\abs{x_i}$ in $\widehat{\abr{x_i}}$. It is evident that the group $G^\prime = \abr{\chi_1}\times\cdots\times \abr{\chi_r}$ is isomorphic to $G$.

    Define the homomorphism $\rho\colon G^\prime\to \widehat G$ by $(\chi_1^{e_1},\dots,\chi_r^{e_r})\mapsto \chi_1^{e_1}\pi_1\cdots\chi_r^{e_r}\pi_r$, where $\pi_i$ denotes the projection $G\to \abr{x_i}$. It is evident that $\ker\rho$ is trivial. Let $f\colon G\to\mathbb C^\times$ be a character of $\widehat G$, and let $\iota_i\colon \abr{x_i}\to G$ denote the inclusion $x_i\mapsto (1,\dots,1,x_i,1,\dots,1)$. Then $f\iota_i\in \widehat{\abr{x_i}}$, and a preimage of $f$ under $\rho$ is $(f\iota_1,\dots,f\iota_r)$. The result follows.
\end{proof}

\begin{proposition}\label{prop: orthogonalities of characters}
    The groups $Z_N^\times$ and $\widehat{Z_N^\times}$ satisfy the following orthogonality relations:
\begin{equation}\label{eqn: ortho rels}
    \sum_{n\in Z_N^\times}\chi(n) = \begin{cases}
    \phi(N) & \text{if $\chi = \mathbf 1$}, \\
    0 & \text{if $\chi\neq \mathbf 1$},
\end{cases}\quad \sum_{\chi\in \widehat{Z_N^\times}}\chi(n) = \begin{cases}
    \phi(N) & \text{if $n = 1$}, \\
    0 & \text{if $n\neq 1$}.
\end{cases}
\end{equation}
\end{proposition}

\begin{proof}
    Let $\chi$ be a Dirichlet character. If $\chi= \mathbf 1$, then $\sum_{n\in Z_N^\times}\chi(n) = \phi(N)$. Suppose $\chi\neq \mathbf 1$, so that there exists $m\in Z_N^\times$ for which $\chi(m)\neq 1$. Then $\sum_{n\in Z_N^\times}\chi(n) = \sum_{n\in Z_N^\times}\chi(mn) = \chi(m)\sum_{n\in Z_N^\times}\chi(n)$, so that $\sum_{n\in Z_N^\times}\chi(n) = 0$.

    Similarly, let $n\in Z_N^\times$. If $n = 1$, $\sum_{\chi\in \widehat{Z_N^\times}}\chi(n) = \phi(N)$. If $n\neq 1$, there exists a character $\eta$ that is not $1$ on $n$, and similarly obtain the equality $\sum_{\chi\in \widehat{Z_N^\times}}\chi(n) =\sum_{\chi\in \widehat{Z_N^\times}}(\eta\chi)(n) =\eta(n)\sum_{\chi\in \widehat{Z_N^\times}}\chi(n)$, from which $\sum_{\chi\in \widehat{Z_N^\times}}\chi(n) = 0$ follows.
\end{proof}

Any Dirichlet character $\chi$ modulo $d$ may be lifted to a character $\chi_N$ modulo $N$ when $d\mid N$, by the rule $\chi_N(n\smod N) = \chi(n\smod d)$ for all $n\in\mathbb{Z}$ coprime to $N$. In other words, if $\pi_{N,d}\colon Z_N^\times\to Z_d^\times$ is the natural projection, then $\chi_N = \chi\circ \pi_{N,d}$.

However, given positive $N,d$ with $d\mid N$ and $\chi$ a character modulo $N$, it is not always possible to find a character $\chi_d$ modulo $d$ such that $\chi = \chi_d\circ \pi_{N,d}$. But for every character modulo $N$ there exists a divisor $d$ of $N$ and a character $\chi_d$ modulo $d$ such that $\chi = \chi_d\circ \pi_{N,d}$. 
\begin{definition}
    The \textib{conductor} \boldindex{$m_\chi$} of a Dirichlet character $\chi$ modulo $N$ is the smallest positive divisor $d$ of $N$ such that there exists a Dirichlet character $\chi_d$ modulo $d$ such that $\chi = \chi_d\circ \pi_{N,d}$, equivalently, such that $\chi$ is trivial on the normal subgroup 
    \[\ker(\pi_{N,d}) = \{n\in Z_N^\times : n\equiv 1\smod d\}.\]

    A Dirichlet character modulo $N$ is called a \boldindex{primitive character} if its conductor is $N$.
\end{definition}

Note that the only character modulo $N$ with conductor $1$ is the trivial character $\mathbf 1_N$, so that the trivial character $\mathbf 1_N$ is primitive only for $N = 1$.

Any Dirichlet character $\chi$ modulo $N$ extends to a function (abusing notation) $\chi\colon Z_N\to\mathbb C$ by the rule $\chi(n) = 0$ for noninvertible elements $n$ of the ring $Z_N$. Further extend $\chi$ to a function $\chi\colon\mathbb Z\to\mathbb C$ by the rule $\chi(n) = \chi(n\smod N)$. The resulting map $\chi\colon \mathbb Z\to \mathbb C$ is a totally multiplicative (set) function; that is, $\chi(nm) = \chi(n)\chi(m)$ for all $n,m\in \mathbb Z$.

For example, the extension of the trivial character $\mathbf 1_N$ to a function on $\mathbb Z$ is given by
\[\mathbf 1_N(n) = \begin{cases}
    1 & \text{if $\gcd(n,N)=1$},\\
    0 & \text{if $\gcd(n,N)\neq1$}.
\end{cases}\]
Also note that the extension of any Dirichlet character $\chi$ satisfies 
\[\chi(0) = \begin{cases}
    1 & \text{if $N = 1$},\\
    0 & \text{if $N > 1$}.
\end{cases}\]
Obtain new orthogonality relations from the ones appearing in \eqref{eqn: ortho rels} by summing from $0$ to $N-1$ in the first orthogonality relation and by taking $n\in\mathbb Z$ in the second:
\begin{equation}\label{eqn: new ortho rels}
    \sum_{n=0}^{N-1}\chi(n) = \begin{cases}
        \phi(N) & \text{if $\chi = \mathbf 1$},\\
        0 & \text{if $\chi\neq \mathbf 1$},
    \end{cases}\quad\sum_{\chi\in \widehat{Z_N^\times}}\chi(n) = \begin{cases}
        \phi(N) & \text{if $n\equiv 1\smod N$},\\
        0 & \text{if $n\not\equiv 1\smod N$}.
    \end{cases} 
\end{equation}
% \begin{definition}
%     Let $\chi$ be a Dirichlet character modulo $N$. The \boldindex{Gauss sum} of $\chi$ is the complex number
%     \[g(\chi) = \sum_{n=0}^{N-1}\chi(n)\mu_N^n,\] where $\mu_N = e^{2\pi i/N}$.
% \end{definition}
% \begin{proposition}
%     If $\chi$ is a primitive character modulo $N$, then for any integer $m$ we have 
%     \[g(\chi) = \sum_{n=0}^{N-1}\chi(n)\mu_N^{nm} = \overline\chi(m)g(\chi).\] Furthermore, the Gauss sum of a primitive character is nonzero.
% \end{proposition}
% \begin{proof}
%     \tbd
% \end{proof}
\begin{lemma}
    Let $N$ be a positive integer. If $N = 1,2$, then every Dirichlet character $\chi$ modulo $N$ satisfies $\chi(-1)=1$. If $N>2$, then the number of Dirichlet characters modulo $N$ is even, of which half satisfy $\chi(-1) = 1$ and the other half satisfy $\chi(-1) = -1$.
\end{lemma}
\begin{proof}
    The extensions of a Dirichlet character to functions on $\mathbb Z$ are unique, so it suffices to study Dirichlet characters given by homomorphisms $Z_N^\times\to\mathbb C^\times$. For $N = 1,2$, $\phi(N) = 1$ so that the only Dirichlet character is the trivial one. For $N>2$, $\phi(N)$ is even, so there are an even number of characters modulo $N$.

    For $N>2$, the map $\widehat{Z_N^\times}\to Z_2$ given by evaluation at $N-1$ (which corresponds to evaluation at $-1$) is a nontrivial homomorphism since there exists a character which is not $1$ on $N-1$. By the first isomorphism theorem, the result follows.
\end{proof}
Dirichlet characters are used to decompose the vector spaces $\mathcal M_k(\varGamma_1(N)), \mathcal S_k(\varGamma_1(N)), \mathcal E_k(\varGamma_1(N))$ into direct sums of interesting subspaces.

For a character $\chi$, define the character (using the same symbol) $\chi\colon \varGamma_0(N)\to\mathbb C^\times$ by $\chi(\gamma) = \chi(d_\gamma)$, where $d_\gamma$ denotes the lower right entry of $\gamma\in\varGamma_0(N)$. (We can also define $\chi$ on any other group of matrices with positive determinant and lower left entry a multiple of $N$.)

\begin{proposition}\label{prop: eigenspace decomp}
    For each Dirichlet character $\chi$ modulo $N$, define the \boldindex{$\chi$-eigenspace} of $\mathcal M_k(\varGamma_1(N))$ by 
    \[\text{\boldindex{$\mathcal M_k(N,\chi)$}} = \{f\in \mathcal M_k(\varGamma_1(N)) : f|_k[\gamma] = \chi(\gamma)f \text{ for all } \gamma\in\varGamma_0(N)\}.\] Similarly define $\chi$-eigenspaces of $\mathcal S_k(\varGamma_1(N))$. Then the following decompositions hold:
    \[\mathcal M_k(\varGamma_1(N)) = \bigoplus_\chi\mathcal M_k(N,\chi),\quad \mathcal S_k(\varGamma_1(N)) = \bigoplus_\chi \text{\boldindex{$\mathcal S_k(N,\chi)$}}.\]
\end{proposition}
We prove this result at the beginning of \cref{section hecke operators}.

% \sai{do i need this discussion of dirichlet L functions}
% \begin{definition}
%     To every Dirichlet character $\chi$ modulo $N$ there is an associated Dirichlet series, called a \boldindex{Dirichlet $L$-function}, given by 
%     \[L(s,\chi) = \sum_{n=1}^\infty \frac{\chi(n)}{n^s} = \prod_{p\in\mathbb P}\frac{1}{1-\chi(p)p^{-s}},\quad \Re(s)>1.\] Here $\mathbb P$ denotes the set of prime numbers and the second equality above is obtained by using the fundamental theorem of arithmetic combined with the fact that $(1-\chi(p)p^{-s})^{-1}$ is given by a geometric series.
% \end{definition}
% It can be shown that these $L$-functions extend meromorphically to the complex plane in $s$, with entire extension unless $\chi = \mathbf 1_N$, which has a simple pole at $s = 1$. (Indeed, $L(s,\mathbf 1_N) = \prod_{p\nmid N}(1-p^{-s})^{-1} = \prod_{p\in\mathbb P}(1-p^{-s})^{-1}\prod_{p\mid N}(1-p^{-s}) = \zeta(s)\prod_{p\mid N}(1-p^{-s})$.) Furthermore, when $\chi(-1) = 1$, the functional equation satisfied by $L(s,\chi)$ is 
% \[\pi^{-s/2}\varGamma\Big(\frac{s}{2}\Big)N^sL(s,\chi) = \pi^{-(1-s)/2}\varGamma\Big(\frac{1-s}{2}\Big)g(\chi)L(1-s,\chi),\] and when $\chi(-1)=-1$, the functional equation is 
% \[\pi^{-(s+1)/2}\varGamma\Big(\frac{s+1}{2}\Big)N^sL(s,\chi) = -i\pi^{-(2-s)/2}\varGamma\Big(\frac{2-s}{2}\Big)g(\chi)L(1-s,\chi).\] Here $\varGamma$ denotes Euler's Gamma function (extended to $\mathbb C$).

Before defining $L$-functions for modular forms, we make a few observations. Let $\varGamma$ be a congruence subgroup, so that by definition there exists a positive integer $N$ such that $\varGamma(N)\subset \varGamma$, from which it follows that $\mathcal M_k(\varGamma)\subset \mathcal M_k(\varGamma(N))$. Furthermore, $\varGamma_1(N^2)$ is contained in 
\[\begin{pmatrix}
    N & 0 \\ 0 & 1
\end{pmatrix}^{\!-1}\varGamma(N)\begin{pmatrix}
    N & 0 \\ 0 & 1
\end{pmatrix} = \Big\{\abcd\in \slz : c\equiv 0 \smod {N^2},~ a\equiv d\equiv 1\smod N\Big\},\] from which we have
\[\begin{pmatrix}
    N & 0 \\ 0 & 1
\end{pmatrix}\varGamma_1(N^2)\subset \varGamma(N)\begin{pmatrix}
    N & 0 \\ 0 & 1
\end{pmatrix}.\] So if $f\in \mathcal M_k(\varGamma(N))$,
\[f(Nz) = f\big(\big(\!\begin{smallmatrix}
    N & 0 \\ 0 & 1
\end{smallmatrix}\!\big)z\big) = N^{-k/2}\Big(f|_k[\big(\!\begin{smallmatrix}
    N & 0 \\ 0 & 1
\end{smallmatrix}\!\big)]\Big)(z)\] and $N^{-k/2}f|_k[\big(\!\begin{smallmatrix}
    N & 0 \\ 0 & 1
\end{smallmatrix}\!\big)]\in \mathcal M_k(\varGamma_1(N^2))$. Indeed, $f(Nz)$ is holomorphic on $\mathcal H$ since $f$ is holomorphic on $\mathcal H$. For $\gamma\in\varGamma_1(N^2)$, $\big(\!\begin{smallmatrix}
    N & 0 \\ 0 & 1
\end{smallmatrix}\!\big)\gamma\in \varGamma(N)\big(\!\begin{smallmatrix}
    N & 0 \\ 0 & 1
\end{smallmatrix}\!\big)$ so that
\[N^{-k/2}f|_k[\big(\!\begin{smallmatrix}
    N & 0 \\ 0 & 1
\end{smallmatrix}\!\big)]|_k[\gamma] = N^{-k/2}f|_k[\gamma^\prime\big(\!\begin{smallmatrix}
    N & 0 \\ 0 & 1
\end{smallmatrix}\!\big)] = N^{-k/2}f|_k[\big(\!\begin{smallmatrix}
    N & 0 \\ 0 & 1
\end{smallmatrix}\!\big)],\] where $\gamma^\prime$ is an element of $\varGamma(N)$. Lastly, let $\alpha\in \slz$, and note that $\slz$ is normal in $\GL_2(\mathbb Z)$ so that $\big(\!\begin{smallmatrix}
    N & 0 \\ 0 & 1
\end{smallmatrix}\!\big)\alpha = \alpha^\prime \big(\!\begin{smallmatrix}
    N & 0 \\ 0 & 1
\end{smallmatrix}\!\big)$ for some $\alpha^\prime\in\slz$. Then 
\[N^{-k/2}\Big(f|_k[\big(\!\begin{smallmatrix}
    N & 0 \\ 0 & 1
\end{smallmatrix}\!\big)]|_k[\alpha]\Big)(z) = N^{-k/2}\Big(f|_k[\alpha^\prime\big(\!\begin{smallmatrix}
    N & 0 \\ 0 & 1
\end{smallmatrix}\!\big)]\Big)(z) = j\big(\alpha^\prime\big(\!\begin{smallmatrix}
    N & 0 \\ 0 & 1
\end{smallmatrix}\!\big),z\big)^{-k}f(\alpha^\prime(Nz)) = (f|_k[\alpha^\prime])(Nz),\] which is holomorphic at $\infty$ since $f\in\mathcal M_k(\varGamma(N))$. If the  series expansion of $f$ at $\infty$ is $f(z) = \sum_{n =0}^\infty a_ne^{2\pi i n z/N}$, we have the Fourier expansion
\[f(Nz) = \sum_{n =0}^\infty a_ne^{2\pi i n z},\] and vice versa. So for the purposes of defining $L$-functions for modular forms, it suffices to define them for modular forms in $\mathcal M_k(\varGamma_1(N))$. We require a few preliminary results about Dirichlet series:
\begin{lemma}\label{lem: miyake lem 3.2.1}
    Assume that both $\sum_{n=1}^\infty a_nn^{-s}$ and $\sum_{n=1}^\infty b_nn^{-s}$ are absolutely convergent at $s = \sigma_0$ with $\sigma_0>0$ real. If $\sum_{n=1}^\infty a_nn^{-s} = \sum_{n=1}^\infty b_nn^{-s}$ on $\Re(s)\geq \sigma_0$, then $a_n = b_n$ for all $n$.
\end{lemma}
\begin{proof}
    By taking differences, it suffices to show that if $\sum_{n=1}^\infty a_nn^{-s} = 0$, then $a_n = 0$ for all $n$. Since $\sum_{n=1}^\infty a_nn^{-s}$ is absolutely convergent at $s = \sigma_0$, it is absolutely and uniformly convergent on $\Re(s)\geq \sigma_0$. Suppose there exists a smallest integer $m$ such that $a_m \neq 0$.
    
    By hypothesis, $-a_m = \sum_{n=m+1}^\infty a_n (n/m)^{-\sigma}$. Let $\sigma = \Re(s)\geq \sigma_0$, and note that for $n>m^2$, we have $(n/m)^{-\sigma}<n^{-\sigma/2}$. It follows that
    \begin{align*}
        \abs{a_m} &\leq \sum_{n=m+1}^\infty \abs{a_n} (n/m)^{-\sigma}\\
        &\leq \sum_{n=m+1}^{m^2} \abs{a_n} (n/m)^{-\sigma} + \sum_{n=m^2+1}^\infty \abs{a_n} n^{-\sigma/2}.
    \end{align*}
    Choose $N$ large enough so that
    \[\sum_{n=N+1}^\infty \abs{a_n}n^{-\sigma_0}\leq \abs{a_m}/3,\]
    and choose $\sigma>2\sigma_0$ large enough so that 
    \[\sum_{n=m+1}^{m^2} \abs{a_n} (n/m)^{-\sigma} + \sum_{n=m^2+1}^N \abs{a_n}n^{-\sigma/2}\leq \abs{a_m}/3.\] With these choices, $\abs{a_m}\leq 2\abs{a_m}/3$, which is a contradiction.
\end{proof}
\begin{proposition}\label{prop: miyake prop pre-lemma 4.3.3}
    Let $f$ be a holomorphic function on $\mathcal H$ such that: 
    \begin{enumerate}[label = \textup{(\arabic*)}]
        \item The Fourier expansion for $f$, \[f(z) = \sum_{n=0}^\infty a_ne^{2\pi i nz},\] converges absolutely and uniformly on compact subsets of $\mathcal H$, and
        \item there exists $v>0$ such that \[f(z) = O(\Im(z)^{-v}) \quad \textup{($\Im(z)\to 0$)}\] uniformly on $\Re(z)$. 
    \end{enumerate}
    Then $a_n = O(n^v)$.
\end{proposition}
\begin{proof}
    By hypothesis, there exists a constant $M>0$ such that $\abs{f(z)}\leq M\Im(z)^{-v}$. Then \begin{align*}
        \abs{a_n} &= \abs{\int_0^1 f(x+iy)e^{-2\pi i n (x+iy)}\dd x}\\
        &\leq My^{-v}e^{2\pi n y}.
    \end{align*} By taking $y = 2/n$, it follows that $\abs{a_n} \leq (Me^{4\pi}2^{-v})n^v$.
\end{proof} A converse to the above proposition is \cref{lem: miyake lem 4.3.3}.
\begin{lemma}[Phragmen-Lindel\"of]\label{lem: miyake lem 4.3.4}
    For two real numbers $v_1,v_2$ with $v_1<v_2$, let 
    \[F = \cbr{s\in \mathbb C : v_1\leq \Re(s)\leq v_2}.\] Let $\phi$ be a holomorphic function on a domain containing $F$ satisfying \[\abs{\phi(s)} = O\big(e^{\abs{\tau}^\delta}\big)\quad \textup{($\abs{\tau}\to\infty$, with $s = \sigma + i\tau$),}\]
    uniformly on $F$ with $\delta >0$. For a real number $b$, if 
    \[\abs{\phi(s)} = O(\abs{\tau}^b)\quad (\abs{\tau}\to\infty) \quad \text{on $\Re(s) = v_1$ and $\Re(s) = v_2$,}\]
    then $\abs{\phi(s)} = O(\abs{\tau}^b)$ as $\abs{\tau}\to\infty$ uniformly on $F$.
\end{lemma}
\begin{proof}
    By hypothesis, there exists $L>0$ such that $\abs{\phi(s)}\leq Le^{\abs{\tau}^\delta}$. Suppose first that $b = 0$. Then there exists $M>0$ such that $\abs{\phi(s)}\leq M$ on $\Re(s) = v_1$ and on $\Re(s) = v_2$. Let $m$ be a positive integer with $m\equiv 2\smod 4$ and let $s = \sigma + i \tau$. Since $\Re(s^m) = \Re((\sigma  + i\tau)^m)$ is a polynomial of $\sigma$ and $\tau$ with highest term of $\tau$ given by $-\tau^m$, we have 
    \[\Re(s^m) = -\tau^m + O(\abs{\tau}^{m-1})\quad (\abs{\tau}\to\infty),\] uniformly on $F$. With $m$ even, $\Re(s^m)$ has an upper bound on $F$. Choose $m$ and $N$ so that $m>\delta$ and $\Re(s^m)<N$, so that for any $\varepsilon>0$,
    \[\big|\phi(s)e^{\varepsilon s^m}\big|\leq Me^{\varepsilon N}\quad \text{on $\Re(s) = v_1$ and $\Re(s) = v_2$,}\]
    and
    \[\big|\phi(s)e^{\varepsilon s^m}\big| = O\big(e^{\abs{\tau}^\delta -\varepsilon\tau^m + K\abs{\tau}^{m-1}}\big)\quad (\abs{\tau}\to\infty\text{, so this quantity tends to zero in the limit}),\] uniformly on $F$. By the maximum principle, it follows that $\big|\phi(s)e^{\varepsilon s^m}\big|\leq Me^{\varepsilon N}$ for $s\in F$. Let $\varepsilon$ tend to $0$ so that $\abs{\phi(s)}\leq M$, that is, $\phi(s) = O(\abs{\tau}^0)$.

    Now let $b\neq 0$. Let $\psi(s) = (s-v_1+1)^b = e^{b\log(s-v_1+1)}$ (using the principal branch of the logarithm); note that $\psi$ is holomorphic. Since $\Re(\log(s-v_1+1)) = \log\abs{s-v_1+1}$, we have 
    \[\abs{\psi(s)} = \abs{s-v_1+1}^b\sim \abs{\tau}^b\quad (\abs{\tau}\to\infty),\] uniformly on $F$ (the notation $\abs{s-v_1+1}^b\sim \abs{\tau}^b$ means that $\lim_{\abs{\tau}\to\infty} \abs{s-v_1+1}^b/\abs{\tau}^b = 1$).

    Let $\phi_1(s) = \phi(s)/\psi(s)$. The function $\phi_1$ satisfies the same assumptions as $\phi$ with $b = 0$, so by repeating the above argument with $\phi_1(s)$ in place of $\phi(s)$, we find that $\phi_1(s)$ is bounded on $F$. It follows that $\abs{\phi(s)} = O(\abs{\tau}^b)$ for $\abs{\tau}\to\infty$. 
    %\sai{I may rewrite the last part of this proof.}
\end{proof}
\begin{definition}\label{def: L function}
    To a holomorphic function $f\colon \mathcal H\to \mathbb C$ satisfying \cref{prop: miyake prop pre-lemma 4.3.3}  with Fourier series expansion $f(z) = \sum_{n=0}^\infty a_ne^{2\pi i nz}$, we associate to it an \boldindex{$L$-function}, given by the Dirichlet series
    \[L(s,f) = \sum_{n=1}^\infty a_nn^{-s}.\] Furthermore, for $N>0$, define \boldindex{$\varLambda_N(s,f)$} by 
    \[\varLambda_N(s,f) = (2\pi/\sqrt{N})^{-s}\varGamma(s)L(s,f).\qedhere\]
\end{definition}

In the above definition, since $a_n = O(n^v)$, the Dirichlet series $L(s,f)$ converges absolutely and uniformly on compact subsets of $\cbr{s\in\mathbb C : \Re(s)>1+v}$. We are interested in $L$-functions of modular forms $f\in\mathcal M_k(\varGamma_1(N))$, which do satisfy \cref{prop: miyake prop pre-lemma 4.3.3}. 

\begin{definition}
    For a positive integer $N$, define \boldindex{$\omega_N$} $\in\GL_2(\mathbb R)$ by 
    \[\omega_n = \begin{pmatrix}
        0 & -1 \\ N & 0
    \end{pmatrix}.\]
    Its action on functions on $\mathcal H$ appears in several results forthcoming, so we define it here.
\end{definition}

\begin{theorem}[Hecke]\label{thm: miyake thm 4.3.6}
    Fix positive integers $k,N$. Let $f,g$ be holomorphic functions on $\mathcal H$ satisfing the hypotheses of \cref{prop: miyake prop pre-lemma 4.3.3}, with Fourier expansions $f(z) = \sum_{n=0}^\infty a_ne^{2\pi i nz}$ and $g(z) = \sum_{n=0}^\infty b_ne^{2\pi i nz}$. Then the following conditions are equivalent: 
    \begin{enumerate}[label = \textup{(\alph*)}]
        \item $g(z) = (f|_k[\omega_N])(z) = (\sqrt{N}z)^{-k}f(-1/Nz)$
        \item Both $\varLambda_N(s,f)$ and $\varLambda(s,g)$ can be analytically continued to the whole $s$-plane, satisfy the functional equation \[\varLambda_N(s,f) = i^k\varLambda_N(k-s,g),\] and the function \[\varLambda_N(s,f) + \frac{a_0}{s} + \frac{i^kb_0}{k-s}\] is holomorphic on the whole $s$-plane and is bounded on any vertical strip.
    \end{enumerate}
\end{theorem}
\begin{proof}
    Suppose (a) holds. Since there exists $v>0$ such that $a_n = O(n^v)$ and $b_n = O(n^v)$, 
    \[\sum_{n=1}^\infty \abs{a_n}e^{-2\pi n t/\sqrt{N}}\quad (t>0)\] and 
    \[\sum_{n=1}^\infty \int_0^\infty \abs{a_n}t^\sigma e^{-2\pi n t/\sqrt{N}}t^{-1}\dd t \quad (\sigma > v+1)\] are convergent. Hence for $\Re(s)>v+1$, 
    \begin{align*}
        \varLambda_N(s,f) &= \sum_{n=1}^\infty a_n(2\pi n/\sqrt{N})^{-s}\int_0^\infty t^{s-1}e^{-t}\dd t\\
        &= \sum_{n=1}^\infty \int_0^\infty a_nt^s e^{-2\pi n t/\sqrt{N}}\\
        &= \int_0^\infty t^s\bigg(\sum_{n=1}^\infty a_n e^{-2\pi nt /\sqrt{N}}\bigg)t^{-1}\dd t\\
        &= \int_0^\infty t^s[f(it/\sqrt{N}) - a_0]t^{-1}\dd t\\
        &= -\frac{a_0}{s} + \int_1^\infty t^{-s}f(i/\sqrt{N}t)t^{-1}\dd t + \int_1^\infty t^s[f(it/\sqrt{N}) - a_0]t^{-1}\dd t.
    \end{align*}
    Since $g(z) = (\sqrt{N}z)^{-k}f(-1/Nz)$, the above equality becomes 
    \begin{equation}\label{eqn: new equality 4.3.6}
        \varLambda_N(s,f) = -\frac{a_0}{s} - \frac{i^kb_0}{k-s} + i^k\int_1^\infty t^{k-s}[g(it/\sqrt{N})-b_0]t^{-1}\dd t + \int_1^\infty t^s[f(it/\sqrt{N}) - a_0]t^{-1}\dd t,
    \end{equation} which holds for $\Re(s)>\max\{k,v+1\}$. By \cref{lem: miyake lem 4.3.3}, as $t$ tends to $\infty$, $f(it)-a_0 = O(e^{-2\pi t})$ and $g(it)-a_0 = O(e^{-2\pi t})$ so that 
    \[\int_1^\infty t^s[f(it/\sqrt{N}) - a_0]t^{-1}\dd t\quad\text{and}\quad \int_1^\infty t^{k-s}[g(it/\sqrt{N})-b_0]t^{-1}\dd t\] are convergent absolutely and uniformly on any vertical strip. Therefore the functions these integrals define are holomorphic on the whole $s$-plane. It follows by \eqref{eqn: new equality 4.3.6} that $\varLambda_N(s,f)$ is a meromorphic function on the whole $s$-plane with $\varLambda_N(s,f) + a_0/s+i^kb_0/(k-s)$ an entire bounded function on any vertical strip. Similarly analytically continue $\varLambda_N(s,g)$ to the whole $s$-plane, satisfying 
    \begin{equation}\label{eqn: similar for g 4.3.6}
        i^k\varLambda_N(k-s,g) = -\frac{a_0}{s}-\frac{i^kb_0}{k-s} + i^k\int_1^\infty t^{k-s}[g(it/\sqrt{N}-b_0)]t^{-1}\dd t + \int_1^\infty t^s[f(it/\sqrt{N})-a_0]t^{-1}\dd t.
    \end{equation}
    It follows from \eqref{eqn: new equality 4.3.6} and \eqref{eqn: similar for g 4.3.6} that $\varLambda_N(s,f) = i^k\varLambda_N(k-s,g)$.

    Conversely, suppose that (b) holds. Since $e^{-e^x}e^{\sigma x}$ is a Schwartz function for $\sigma>0$, the inverse Mellin transform 
    \[e^{-t} = \frac{1}{2\pi i}\int_{\Re(s) = \sigma}\varGamma(s)t^{-s}\dd s\]
    holds for $\sigma>0$. It follows that 
    \[f(iy) = a_0 + \frac{1}{2\pi i}\sum_{n=1}^\infty a_n\int_{\Re(s) =\alpha}\varGamma(s)(2\pi n y)^{-s}\dd s\] for any $\alpha>0$. Let $\alpha>v+1$, so that $L(s,f) = \sum_{n=1}^\infty a_nn^{-s}$ is uniformly convergent and bounded on $\Re(s) = \alpha$. In this case, Stirling's estimate $\varGamma(s)\sim \sqrt{2\pi} \tau^{\sigma -1/2}e^{-\pi\abs{\tau}/2}$ (for $s = \sigma + i\tau$ and $\abs{\tau}\to\infty$) shows that $\varLambda_N(s,f) = (2\pi/\sqrt{N})^{-s}\varGamma(s)L(s,f)$ is absolutely integrable, so that the order of integration and summation above may be interchanged to obtain
    \begin{equation}\label{eqn: integral for Hecke lemma}
        f(iy) = a_0 + \frac{1}{2\pi i} \int_{\Re(s) = \alpha} (\sqrt{N}y)^{-s}\varLambda_N(s,f)\dd s.
    \end{equation}
    Since $L(s,f)$ is bounded on $\Re(s) = \alpha$, by Stirling's estimate we have for any $\mu>0$ that 
    \begin{equation}\label{eqn: varLambda big O}
        |\varLambda_N(s,f)| = O(|\Im(s)|^{-\mu})\quad (|\Im(s)|\to\infty)
    \end{equation}
    on $\Re(s) = \alpha$. Choose $\beta$ so that $k-\beta > v+1$. Using a similar argument we deduce that for any $\mu>0$, 
    \[|\varLambda_N(s,f)| = |\varLambda_N(k-s,g)| = O(|\Im(s)|^{-\mu})\quad (|\Im(s)|\to\infty)\] on $\Re(s) = \beta$. By assumption, $\varLambda_N(s,f) + a_0/s+i^kb_0/(k-s)$ is bounded on the vertical strip $\beta\leq \Re(s)\leq \alpha$. Thus for any $\mu>0$, \cref{lem: miyake lem 4.3.4} implies that \eqref{eqn: varLambda big O} holds uniformly on the strip $\beta\leq \Re(s)\leq \alpha$.
    
    We may also assume that $\alpha>k$ and $\beta<0$. Observe that $(\sqrt{N}y)^{-s}\varLambda_N(s,f)$ has simple poles at $s = 0,k$ with corresponding residues $-a_0,(\sqrt{N}y)^{-k}i^kb_0$. Combined with the fact that \eqref{eqn: varLambda big O} holds uniformly on the strip $\beta\leq \Re(s)\leq \alpha$, we may change the path of integration from $\Re(s) = \alpha$ to $\Re(s) = \beta$ in \eqref{eqn: integral for Hecke lemma} to obtain
    \[f(iy) = (\sqrt{N}y)^{-k}i^kb_0 + \frac{1}{2\pi i} \int_{\Re(s) = \beta} (\sqrt{N}y)^{-s}\varLambda_N(s,f) \dd s.\]
    From the functional equation in (b), we have 
    \begin{align*}
        f(iy) &= (\sqrt{N}y)^{-k}i^kb_0 + \frac{1}{2\pi i} \int_{\Re(s) = \beta} (\sqrt{N}y)^{-s}i^k\varLambda_N(k-s,g) \dd s\\
        &= i^k\bigg[(\sqrt{N}y)^{-k}b_0 + \frac{1}{2\pi i} \int_{\Re(s) = k-\beta} (\sqrt{N}y)^{s-k}\varLambda_N(s,g) \dd s\bigg]\\
        &= i^k(\sqrt{N}y)^{-k}g(-1/iNy).
    \end{align*}
    Since $f,g$ are holomorphic on $\mathcal H$, it follows that $f(z) = i^k(\sqrt{N}z/i)^{-k}g(-1/Nz)$, or equivalently $g(z) = (\sqrt{N}z)^{-k}f(-1/Nz)$.
\end{proof}
\begin{corollary}\label{cor: miyake cor 4.3.7}
    If $f(z)\in \mathcal S_k(N,\chi)$, then because any element of $\mathcal S_k(N,\chi)$ satisfies the conditions of \cref{prop: miyake prop pre-lemma 4.3.3} \textup{(}see \cref{cor: cuspform growth corollary}\textup{)}, it follows that $\varLambda_N(s,f)$ is holomorphic in $s$ and satisfies the functional equation
    \[\varLambda_N(s,f) = i^k\varLambda_N(k-s,f|_k[\omega_N]).\]
\end{corollary}

\newpage\section{Hecke operators}
We outline how double cosets act on spaces of modular forms, and using them we define the Hecke operators $\abr{n}$ and $T_n$ on the space $\mathcal M_k(\varGamma_1(N))$. The Hecke operators are akin to ``averaging'' operators which act on each of the eigenspaces from \cref{prop: eigenspace decomp}. We collect several properties of the Hecke operators for later use. For example, if a modular form is a common eigenfunction of all of the Hecke operators $T_n$, then its Fourier coefficients are proportional to its eigenvalues. Following the discussion of the Hecke operators, we restrict our view to the space of cusp forms $\mathcal S_k(\varGamma_1(N))$, on which we define the Petersson inner product and find adjoints to the Hecke operators $\abr{p}$ and $T_p$ for primes $p$ not dividing $N$. By doing so, we deduce that all of the Hecke operators $\abr{n}$ and $T_n$ are normal operators on $\mathcal S_k(\varGamma_1(N))$ whenever $n$ is coprime to $N$.
\subsection{Double coset operators}
Let $\varGamma_1,\varGamma_2$ be congruence subgroups of $\slz$. For $\alpha\in\glqp$, the set
\[\varGamma_1\alpha\varGamma_2 = \{\gamma_1\alpha\gamma_2 : \gamma_1\in\varGamma_1, \gamma_2\in\varGamma_2\}\] is a \boldindex{double coset} in $\glqp$. The group $\varGamma_1$ acts on the double coset $\varGamma_1\alpha\varGamma_2$ on the left by multiplication, partitioning it into orbits of the form $\varGamma_1\beta$, where $\beta = \gamma_1\alpha\gamma_2$ is some representative for this orbit. We show that the orbit space  $\varGamma_1\backslash\varGamma_1\alpha\varGamma_2$ is a finite disjoint union $\bigsqcup_j\varGamma_1\beta_j$ for some choice of representatives $\beta_j$.
\begin{lemma}\label{lem: conjugation congruence subgroup}
    Let $\alpha\in\glqp$ and $\varGamma$ be a congruence subgroup. Then $\alpha^{-1}\varGamma\alpha\cap\slz$ is also a congruence subgroup.
\end{lemma}
\begin{proof}
    There exists $N$ such that $\varGamma(N)$ is contained in $\varGamma$. Let $M$ be the least common multiple of $N$ and the entries of the matrices $\alpha$ and $\alpha^{-1}$ so that $\varGamma(M)\subset \varGamma$ and $M\alpha, M\alpha^{-1}$ are integer-valued matrices.
    
    Observe that $\varGamma(M^3)$ is contained in $I + M^3\Mat_2(\mathbb Z)$ so that 
    \[\alpha\varGamma(M^3)\alpha^{-1}\subset \alpha(I + M^3\Mat_2(\mathbb Z))\alpha^{-1} = I + M\cdot (M\alpha )\Mat_2(\mathbb Z) (M\alpha^{-1})\subset I+M\Mat_2(\mathbb{Z}).\]
    
    Elements of $\alpha\varGamma(M^3)\alpha^{-1}$ have determinant $1$, so $\alpha\varGamma(M^3)\alpha^{-1}\subset \slz$, from which it follows that $\alpha\varGamma(M^3)\alpha^{-1}\subset \varGamma(M)$. Hence $\varGamma(M^3)\subset \alpha^{-1}\varGamma(M)\alpha \subset \alpha^{-1}\varGamma\alpha$, and so $\varGamma(M^3)\subset \alpha^{-1}\varGamma\alpha\cap \slz$.
\end{proof}
\begin{lemma}\label{lem: double coset coset/orbit spaces}
    Let $\alpha\in\glqp$ and $\varGamma_1,\varGamma_2$ be congruence subgroups. Let $\varGamma_3 = \alpha^{-1}\varGamma_1\alpha\cap \varGamma_2\subset \varGamma_2$. Then the map $\varGamma_2\to \varGamma_1\alpha\varGamma_2$ given by left multiplication by $\alpha$, $\gamma_2\mapsto\alpha\gamma_2$, induces a natural bijection of the coset space $\varGamma_3\backslash\varGamma_2$ to the orbit space $\varGamma_1\backslash\varGamma_1\alpha\varGamma_2$.

    In other words, $\{\gamma_{2,j}\}$ is a set of coset representatives for $\varGamma_3\backslash \varGamma_2$ if and only if $\{\alpha\gamma_{2,j}\}$ is a set of orbit representatives for $\varGamma_1\backslash \varGamma_1\alpha\varGamma_2$.
\end{lemma} 
\begin{proof}
    It is evident that the map $\varGamma_2\to \varGamma_1\backslash\varGamma_1\alpha\varGamma_2$ given by $\gamma_2\mapsto \varGamma_1\alpha\gamma_2$ is surjective, since elements of $\varGamma_1\backslash\varGamma_1\alpha\varGamma_2$ are of the form $\varGamma_1\gamma_1\alpha\gamma_2$ for some $\gamma_1\in\varGamma_1$ and $\gamma_2\in\varGamma_2$. Two elements $\gamma_2,\gamma_2^\prime\in\varGamma_2$ map to the same orbit if $\varGamma_1\alpha\gamma_2 = \varGamma_1\alpha\gamma_2^\prime$, that is, if $\gamma_2^\prime\gamma_2^{-1}\in \alpha^{-1}\varGamma_1\alpha$. It follows that by taking $\varGamma_3 = \alpha^{-1}\varGamma_1\alpha\cap \varGamma_2$ as above, that the induced map from the coset space $\varGamma_3\backslash\varGamma_2\to \varGamma_1\backslash \varGamma_1\alpha\varGamma_2$ is a bijection.
\end{proof}
\begin{lemma}\label{lem: congruence subgroups commensurable}
    Let $\varGamma_1,\varGamma_2$ be congruence subgroups. Then they are \textib{commensurable}, that is, the indices $[\varGamma_1 : \varGamma_1\cap \varGamma_2], [\varGamma_2 : \varGamma_1\cap \varGamma_2]$ are finite.
\end{lemma}
\begin{proof}
    There exist positive integers $N_1,N_2$ such that $\varGamma(N_1)\subset \varGamma_1$ and $\varGamma(N_2)\subset \varGamma_2$. Then $\varGamma(\lcm(N_1,N_2))\subset \varGamma(N_1)\cap \varGamma(N_2)\subset \varGamma_1\cap \varGamma_2$ since $N_1,N_2$ divide their least common multiple. Then \begin{align*}
        [\varGamma_1 : \varGamma(\lcm(N_1,N_2))] &= [\varGamma_1 : \varGamma_1\cap \varGamma_2][\varGamma_1\cap \varGamma_2 : \varGamma(\lcm(N_1,N_2))] \text{ and}\\
        [\varGamma_2 : \varGamma(\lcm(N_1,N_2))] &= [\varGamma_2 : \varGamma_1\cap \varGamma_2][\varGamma_1\cap \varGamma_2 : \varGamma(\lcm(N_1,N_2))].
    \end{align*} The indices $[\varGamma_1 : \varGamma(\lcm(N_1,N_2))],[\varGamma_2 : \varGamma(\lcm(N_1,N_2))]$ are finite (see the discussion following \cref{def: congruence subgroup}), so that the indices $[\varGamma_1 : \varGamma_1\cap \varGamma_2], [\varGamma_2 : \varGamma_1\cap \varGamma_2]$ are also finite.
\end{proof}
In \cref{lem: conjugation congruence subgroup}, since $\alpha^{-1}\varGamma\alpha\cap\slz$ is a congruence subgroup, the coset space $\varGamma_3\backslash \varGamma_2$ from \cref{lem: double coset coset/orbit spaces} is finite, so that the orbit space $\varGamma_1\backslash \varGamma_1\alpha\varGamma_2$ is also finite as desired. With this result, we can define an action of double cosets on modular forms.

\begin{definition}\label{def: weight k double coset operator}
    Let $\alpha\in \glqp$ and $\varGamma_1,\varGamma_2$ be congruence subgroups. The \boldindex{weight-$k$ $\varGamma_1\alpha\varGamma_2$ operator} (``double coset operator'') $|_k[\varGamma_1\alpha\varGamma_2]$ is given by 
    \[f|_k[\varGamma_1\alpha\varGamma_2] = (\det \alpha)^{k/2-1}\sum_jf|_k[\beta_j],\] where $\{\beta_j\}$ is a set of orbit representatives of $\varGamma_1\alpha\varGamma_2$ (that is, $\varGamma_1\alpha\varGamma_2 = \bigsqcup_j \varGamma_1\beta_j$). It takes modular forms in $\mathcal M_k(\varGamma_1)$ to modular forms in $\mathcal M_k(\varGamma_2)$, and takes cusp forms in $\mathcal S_k(\varGamma_1)$ to cusp forms in $\mathcal S_k(\varGamma_2)$.
\end{definition}
We check that this definition is well defined, that is, independent of the choice of orbit representatives for $\varGamma_1\alpha\varGamma_2$. Let $\beta = \gamma_1\alpha\gamma_2$ and $\beta^\prime = \gamma_1^\prime\alpha\gamma_2^\prime$ represent the same orbit in $\varGamma_1\backslash\varGamma_1\alpha\varGamma_2$, so that $\varGamma_1\beta = \varGamma_1\beta^\prime$. It follows that $\alpha\gamma_2\in\varGamma_1\alpha\gamma_2^\prime$. Since $f$ is weight-$k$ invariant under $\varGamma_1$, we have $f|_k[\beta] = f|_k[\alpha\gamma_2] = f|_k[\eta_1\alpha\gamma_2^\prime] = f|_k[\alpha\gamma_2^\prime] = f|_k[\beta^\prime]$. Thus the action of the double coset $\varGamma_1\alpha\varGamma_2$ is well defined.

We also check that the double coset operator takes modular forms to modular forms and cusp forms to cusp forms. For $f\in \mathcal M_k(\varGamma_1)$, we verify that $f|_k[\varGamma_1\alpha\varGamma_2]$ is $\varGamma_2$-invariant and holomorphic at the cusps of $\varGamma_2$.

Note that any $\gamma_2\in\varGamma_2$ permutes the orbit space $\varGamma_1\backslash\varGamma_1\alpha\varGamma_2$ by right multiplication. So given a set of orbit representatives $\cbr{\beta_j}$ for $\varGamma_1\backslash\varGamma_1\alpha\varGamma_2$, then $\cbr{\beta_j\gamma_2}$ is also a set of orbit representatives. Therefore 
\[(f|_k[\varGamma_1\alpha\varGamma_2])|_k[\gamma_2] = (\det\alpha)^{k/2-1}\sum_j f|_k[\beta_j\gamma_2] = f|_k[\varGamma_1\alpha\varGamma_2].\] It follows that $f|_k[\varGamma_1\alpha\varGamma_2]$ is weight-$k$ invariant under $\varGamma_2$ as needed.

Observe that for any $f\in\mathcal M_k(\varGamma_1)$ and any $\gamma\in\glrp$, the function $f|_k[\gamma]$ is holomorphic at infinity; that is, it has a Fourier expansion 
\[(f|_k[\gamma])(z) = \sum_{n=0}^\infty a_ne^{2\pi i n z/h}\] for some positive integer $h$ by \cref{lem: weight k operator properties}. Then for any $\delta\in\slz$, 
\[(f|_k[\varGamma_1\alpha\varGamma_2])|_k[\delta] = (\det\alpha)^{k/2-1}\sum_{j}f|_k[\beta_j\delta]\]
for some set of orbit representatives $\cbr{\beta_j}$ of $\varGamma_1\alpha\varGamma_2$, from which holomorphicity of $f|_k[\varGamma_1\alpha\varGamma_2]$ at all cusps of $\varGamma_2$ follows, since each summand above is holomorphic at infinity. In the case that $f$ is a cusp form, $f|_k[\gamma]$ vanishes at infinity for any $\gamma\in\glrp$ as we saw in \cref{lem: weight k operator properties}. Therefore the double coset operator takes modular forms to modular forms and cusp forms to cusp forms.

Particular choices of $\varGamma_1,\varGamma_2$ yield notable double coset operators:
\begin{enumerate}[label=(\arabic*)]
    \item If $\varGamma_2\subset \varGamma_1$ and $\alpha = I$, then $f|_k[\varGamma_1\alpha\varGamma_2] = f$ and the double coset operator $|_k[\varGamma_1\alpha\varGamma_2]$ is the natural inclusion of $\mathcal M_k(\varGamma_1)$ into $\mathcal M_k(\varGamma_2)$.
    \item If $\varGamma_2 = \alpha^{-1}\varGamma_1\alpha$, then $f|_k[\varGamma_1\alpha\varGamma_2] = (\det\alpha)^{k/2-1}f|_k[\alpha]$, and the double coset operator $|_k[\varGamma_1\alpha\varGamma_2]$ is an isomorphism of $\mathcal M_k(\varGamma_1)$ with $\mathcal M_k(\varGamma_2)$.
    \item Finally, if $\varGamma_1\subset\varGamma_2$ and $\alpha = I$, take $\{\gamma_{2,j}\}$ to be a set of coset representatives for $\varGamma_1\backslash \varGamma_2$. Then $f|_k[\varGamma_1\alpha\varGamma_2] = \sum_j f|_k[\gamma_{2,j}]$, and the double coset operator $|_k[\varGamma_1\alpha\varGamma_2]$ is the natural trace map projecting $\mathcal M_k(\varGamma_1)$ onto its subspace $\mathcal M_k(\varGamma_2)$ by symmetrizing over the quotient, a surjection. (The action of $\varGamma_1$ on $\varGamma_2$ is to permute the coset representatives of $\varGamma_1\backslash\varGamma_2$, so the action of $\varGamma_1$ on elements in the image of the natural trace map is the identity.)
\end{enumerate}

Any double coset operator is a composition of these three operators. Indeed, given $\varGamma_1,\varGamma_2$, and $\alpha$, set $\varGamma_3 = \alpha^{-1}\varGamma_1\alpha \cap \varGamma_2$ and $\varGamma_3^\prime = \alpha\varGamma_3\alpha^{-1} = \varGamma_1\cap \alpha\varGamma_2\alpha^{-1}$. Then $\varGamma_3^\prime\subset \varGamma_1$, $\alpha^{-1}\varGamma_3^\prime\alpha = \varGamma_3$, and $\varGamma_3\subset \varGamma_2$. Composing their corresponding double coset operators, we have for  $f\in\mathcal M_k(\varGamma_1)$ that $f\mapsto f\mapsto (\det\alpha)^{k/2-1}f|_k[\alpha]\mapsto (\det\alpha)^{k/2-1}\sum_j f|_k[\alpha\gamma_{2,j}]$. By \cref{lem: double coset coset/orbit spaces} the composition above agrees with the double coset operator $|_k[\varGamma_1\alpha\varGamma_2]$.
\subsection{Hecke operators $\abr{n}$ and $T_n$}\label{section hecke operators}
Recall the congruence subgroups \begin{align*}
    \varGamma_0(N) &= \cbr{\abcd\in\slz: \abcd\equiv\begin{pmatrix}
        \ast & \ast \\ 0 & \ast
    \end{pmatrix}\smod N}\text{ and}\\
    \varGamma_1(N) &= \cbr{\abcd\in \slz: \begin{pmatrix}
        a & b \\ c & d 
    \end{pmatrix}\equiv \begin{pmatrix}
        1 & \ast \\ 0 & 1
    \end{pmatrix}\smod N}.
\end{align*} Since $\varGamma_1(N)\subset 
\varGamma_0(N)$, we have the containment $\mathcal M_k(\varGamma_0(N))\subset \mathcal M_k(\varGamma_1(N))$ of modular forms. We define the two operators $\abr{n}$ and $T_n$ on the larger space $\mathcal M_k(\varGamma_1(N))$.

Let $\alpha\in\varGamma_0(N)$ and consider the weight-$k$ double coset operator $|_k[\varGamma_1(N)\alpha\varGamma_1(N)]$. From \cref{lem: gamma 1 normal in gamma 0} we have that $\varGamma_1(N)$ is normal in $\varGamma_0(N)$ and that $\varGamma_0(N)/\varGamma_1(N)\cong (\mathbb Z/N\mathbb Z)^\times$, so the double coset operator $|_k[\varGamma_1(N)\alpha\varGamma_1(N)]$ is of the form (2) in the list of notable double coset operators following \cref{def: weight k double coset operator}. Hence this double coset operator translates a function $f\in \mathcal M_k(\varGamma_1(N))$ to $f|_k[\varGamma_1(N)\alpha\varGamma_1(N)] = f|_k[\alpha]\in \mathcal M_k(\varGamma_1(N))$. In this way the group $\varGamma_0(N)$ acts on $\mathcal M_k(\varGamma_1(N))$, and its subgroup $\varGamma_1(N)$ acts trivially. Thus $\varGamma_0(N)/\varGamma_1(N)\cong (\mathbb Z/N\mathbb Z)^\times$ acts on $\mathcal M_k(\varGamma_1(N))$. In particular, if $d\in (\mathbb Z/N\mathbb Z)^\times$, there exists $\big(\!\begin{smallmatrix}
    a^\prime & b^\prime \\ c^\prime & d^\prime
\end{smallmatrix}\!\big)\in \varGamma_0(N)/\varGamma_1(N)$ with $d^\prime\equiv d\smod N$. From this fact, we obtain the diamond operator $\abr{d}$ on $\mathcal M_k(\varGamma_1(N))$: 
\begin{proposition}
    The action of $d$ on $\mathcal M_k(\varGamma_1(N))$, called the diamond operator $\abr{d}\colon \mathcal M_k(\varGamma_1(N))\to \mathcal M_k(\varGamma_1(N))$, is given by 
    \[\abr{d}f = f|_k[\alpha] \quad \text{for any $\alpha = \begin{pmatrix}
    a^\prime & b^\prime \\ c^\prime & d^\prime
    \end{pmatrix}\in \varGamma_0(N)$ with $d^\prime\equiv d\smod N$}.\]
\end{proposition}

For any character $\chi$, we have that the space $\mathcal M_k(N,\chi)$ from \cref{prop: eigenspace decomp} is really the ``$\chi$-eigenspace of the diamond operators''; that is,
\[\mathcal M_k(N,\chi) = \{f\in \mathcal M_k(\varGamma_1(N)) : \abr{d}f = \chi(d)f \text{ for all } d\in(\mathbb Z/N\mathbb Z)^\times\},\]
with a similar definition for $\mathcal S_k(N,\chi)$. In other words, the diamond operator $\abr{d}$ acts on $\mathcal M_k(\varGamma_1(N)) = \bigoplus_\chi\mathcal M_k(N,\chi)$ by acting on each $\chi$-eigenspace by multiplication by $\chi(d)$.
We can now give an elementary proof of \cref{prop: eigenspace decomp}:
\begin{proposition}
    The following decompositions hold:
    \[\mathcal M_k(\varGamma_1(N)) = \bigoplus_\chi\mathcal M_k(N,\chi),\quad \mathcal S_k(\varGamma_1(N)) = \bigoplus_\chi \mathcal S_k(N,\chi).\]
\end{proposition}
\begin{proof}
    % Recall that for each Dirichlet character $\chi$ modulo $N$, we defined the $\chi$-eigenspace of $\mathcal M_k(\varGamma_1(N))$ by 
    % \[\text{\boldindex{$\mathcal M_k(N,\chi)$}} = \{f\in \mathcal M_k(\varGamma_1(N)) : f|_k[\gamma] = \chi(\gamma)f \text{ for all } \gamma\in\varGamma_0(N)\}\] with a similar definition of $\chi$-eigenspaces of $\mathcal S_k(\varGamma_1(N))$ given by replacing $\mathcal M_k(\varGamma_1(N)$ by $\mathcal S_k(\varGamma_1(N)$.
    For any character $\chi\colon (\mathbb Z/N\mathbb Z)^\times\to\mathbb C^\times$, define the operator $\pi_\chi\colon\mathcal M_k(\varGamma_1(N))\to\mathcal M_k(\varGamma_1(N))$ by
    \[\pi_\chi = \frac{1}{\phi(N)}\sum_{d\in(\mathbb Z/N\mathbb Z)^\times}\chi(d)^{-1}\abr{d}.\]
    Observe that for any $d,e\in(\mathbb Z/N\mathbb Z)^\times$, we have $\abr{d}\abr{e} = \abr{e}\abr{d} = \abr{de}$: For any $f\in\mathcal M_k(N,\chi)$, we have $\abr{d}\abr{e}f = \abr{d}\chi(e)f = \chi(de)f$ as needed. It follows that 
    \begin{align*}
        \pi_\chi^2 &= \frac{1}{\phi(N)}\sum_{d\in(\mathbb Z/N\mathbb Z)^\times}\chi(d)^{-1}\abr{d}\bigg(\frac{1}{\phi(N)}\sum_{e\in(\mathbb Z/N\mathbb Z)^\times}\chi(e)^{-1}\abr{e}\bigg)\\
        &= \frac{1}{\phi(N)^2}\sum_{d\in(\mathbb Z/N\mathbb Z)^\times}\sum_{e\in(\mathbb Z/N\mathbb Z)^\times}\chi(de)^{-1}\abr{de}\\
        &= \frac{1}{\phi(N)}\sum_{d\in(\mathbb Z/N\mathbb Z)^\times}\chi(d)^{-1}\abr{d}\\
        &= \pi_\chi.
    \end{align*} Since $\pi_\chi$ is idempotent, it follows from linear algebra that it is a projection operator of $\mathcal M_k(\varGamma_1(N))$ onto some subspace. Let $f\in\mathcal M_k(\varGamma_1(N))$. Then
    \[\abr{e}\pi_\chi(f) = \abr{e}\frac{1}{\phi(N)}\sum_{d\in(\mathbb Z/N\mathbb Z)^\times}\chi(d)^{-1}\abr{d}f = \chi(e)\frac{1}{\phi(N)}\sum_{d\in(\mathbb Z/N\mathbb Z)^\times}\chi(de)^{-1}\abr{de}f = \chi(e)\pi_\chi(f)\] for any $e\in(\mathbb Z/N\mathbb Z)^\times$. Furthermore, for any $f\in\mathcal M_k(N,\chi)$, we have that $\pi_\chi(f) = f$. Hence $\pi_\chi$ is the projection of $\mathcal M_k(\varGamma_1(N))$ onto $\mathcal M_k(N,\chi)$.

    Choose the characters so that they satisfy the orthogonality conditions in \cref{prop: orthogonalities of characters}. Then 
    \begin{align*}
        \sum_\chi\pi_\chi &= \frac{1}{\phi(N)}\sum_{\chi}\sum_{d}\chi(d)^{-1}\abr{d}\\
        &= \frac{1}{\phi(N)}\sum_{d}\bigg(\sum_{\chi}\chi(d)\bigg)\abr{d}\\
        &= \abr{1}\\
        &= \id_{\mathcal M_k(\varGamma_1(N))},
    \end{align*} and for two nonequal characters $\chi,\chi^\prime$, we have 
    \[\pi_{\chi}\circ\pi_{\chi^\prime} = \frac{1}{\phi(N)}\sum_{d}\chi(d)^{-1}\abr{d}\pi_{\chi^\prime} = \frac{1}{\phi(N)}\sum_{d}\chi(d)^{-1}\chi^\prime(d)\pi_{\chi^\prime} = \frac{1}{\phi(N)}\bigg(\sum_{d}(\chi^{-1}\chi^\prime)(d)\bigg)\pi_{\chi^\prime} = 0.\]
    It follows that the $\chi$-eigenspaces span $\mathcal M_k(\varGamma_1(N))$ and are pairwise disjoint, which proves the result. The same argument follows for cusp forms.
\end{proof}

We extend the definition of the diamond operator $\abr{d}$ for $d\in (\mathbb Z/N\mathbb Z)^\times$ to $\abr{n}$ for $n\in\mathbb{Z}^+$.
\begin{definition}
    The \textib{Hecke operator} \boldindex{$\abr{n}$}$\colon \mathcal M_k(\varGamma_1(N))\to \mathcal M_k(\varGamma_1(N))$ for $n\in\mathbb{Z}^+$ is given by the zero operator when $\gcd(n,N)>1$ and is given by $\abr{\overline n}$ when $\gcd(n,N) = 1$ (where $\overline{\,\cdot\,}$ denotes reduction modulo $N$).
\end{definition}

Let $p$ be a prime number, and consider the weight-$k$ double coset operator $T_p \colon \mathcal M_k(\varGamma_1(N))\to \mathcal M_k(\varGamma_1(N))$ given by $T_pf = f|_k[\varGamma_1(N)\alpha\varGamma_1(N)]$, where $\alpha = \big(\!\begin{smallmatrix}
    1 & 0 \\ 0 & p
\end{smallmatrix}\!\big)$. 
\begin{lemma}\label{lem: replacement for T_n}
    For a prime $p$, the double coset $\varGamma_1(N)\big(\!\begin{smallmatrix}
        1 & 0 \\ 0 & p
    \end{smallmatrix}\!\big)\varGamma_1(N)$ is given by 
    \[\varGamma_1(N)\big(\!\begin{smallmatrix}
        1 & 0 \\ 0 & p
    \end{smallmatrix}\!\big)\varGamma_1(N) = \cbr{\gamma\in \Mat_2(\mathbb Z) : \gamma \equiv \begin{pmatrix}
        1 & \ast \\ 0 & p
    \end{pmatrix}\smod N, \det\gamma = p}.\]
\end{lemma}
\begin{proof}
    We sketch the proof since it is rather computational. The containment \[\varGamma_1(N)\big(\!\begin{smallmatrix}
        1 & 0 \\ 0 & p
    \end{smallmatrix}\!\big)\varGamma_1(N) \subseteq \cbr{\gamma\in \Mat_2(\mathbb Z) : \gamma \equiv \begin{pmatrix}
        1 & \ast \\ 0 & p
    \end{pmatrix}\smod N, \det\gamma = p}\] is clear. We show that the other containment holds.

    Let $L = \mathbb Z^2$, and let $L_0 = \cbr{\big(\!\begin{smallmatrix}
        x \\ y
    \end{smallmatrix}\!\big)\in L : y\equiv 0 \smod N}$. Let $\Mat_2(\mathbb Z)$ acts on $L$ by left multiplication. Choose $\gamma\in\Mat_2(\mathbb Z)$ such that $\gamma\equiv \big(\!\begin{smallmatrix}
        1 & \ast \\ 0 & p
    \end{smallmatrix}\!\big)\smod N$ and $\det \gamma = p$. Because the determinant of $\gamma$ is positive, $[L: \gamma L_0] = [L:L_0][L_0:\gamma L_0] = Np$.
    
    By the theory of finitely generated Abelian groups there exists a basis $\cbr{u,v}$ of $L$ such that $\det (u~v) = 1$ and $\gamma L_0 = mu\mathbb Z\oplus nv\mathbb Z$, where $0<m,n$, $m\mid n$, and $mn = Np$. Write the first column of $\gamma$ as $\big(\!\begin{smallmatrix}
        a_\gamma \\ c
    \end{smallmatrix}\!\big)$ with $a_\gamma\equiv 1\smod N$ and $c\equiv 0 \smod N$. We can also write the first column of $\gamma$ as $\gamma e_1 \in \gamma L_0$ ($e_i$ is the $i$-th standard basis vector), and hence is the zero vector modulo $m$. Since $\gcd(a_\gamma,N) = 1$, we have $\gcd(m,N) = 1$, from which we deduce that $m = 1$ and $n = Np$. Hence $\gamma L_0 = u\mathbb Z\oplus Npv\mathbb Z$.
    
    There are unique subgroups $L_0 = u\mathbb Z \oplus Nv\mathbb Z$ and $\gamma L = u\mathbb Z\oplus pv\mathbb Z$, of index $p$ and $N$ respectively inside $L = u\mathbb Z\oplus v\mathbb Z$, that contain $\gamma L_0$. Let $\gamma_1 = (u~v)$. Because $u\in L_0$, it follows that $\gamma_1\in\varGamma_0(N)$. Let $\gamma_2 = \big(\gamma_1\big(\!\begin{smallmatrix}
        1 & 0 \\ 0 & p
    \end{smallmatrix}\!\big)\big)^{-1}\gamma = \big(\!\begin{smallmatrix}
        a & b \\ c & d
    \end{smallmatrix}\!\big)$ be an element of $\glqp$ with determinant $1$. The condition for $\gamma e_1 \in\gamma L_0$ to hold is given by $au + cpv\in u\mathbb Z\oplus Npv\mathbb Z$, so that we must have $a\in\mathbb Z$ and $c\in N\mathbb Z$. Similarly, for $\gamma e_2$ to be an element of $\gamma L$ we must have $b,d \in\mathbb Z$. It follows that $\gamma_2\in\varGamma_0(N)$. Since $a_\gamma\equiv 1\smod N$, $\gamma e_1\equiv e_1\smod N$ and so the equation $\gamma = \gamma_1\big(\!\begin{smallmatrix}
        1 & 0 \\ 0 & p
    \end{smallmatrix}\!\big)\gamma_2$ shows that $au_1\equiv 1\smod N$ (where $u_1$ is the first component of $u$), using only that $\gamma_1,\gamma_2\in\varGamma_0(N)$. Thus if $\gamma = \gamma_1\big(\!\begin{smallmatrix}
        1 & 0 \\ 0 & p
    \end{smallmatrix}\!\big)\gamma_2$ for $\gamma_1,\gamma_2\in\varGamma_0(N)$ such that either $\gamma_1$ or $\gamma_2$ lies in $\varGamma_1(N)$, then both do.

    It now suffices to show that $\varGamma_0(N)\big(\!\begin{smallmatrix}
        1 & 0 \\ 0 & p
    \end{smallmatrix}\!\big)\varGamma_0(N) = \varGamma_1(N)\big(\!\begin{smallmatrix}
        1 & 0 \\ 0 & p
    \end{smallmatrix}\!\big)\varGamma_0(N)$.
    The containment $\varGamma_0(N)\big(\!\begin{smallmatrix}
        1 & 0 \\ 0 & p
    \end{smallmatrix}\!\big)\varGamma_0(N) \supseteq \varGamma_1(N)\big(\!\begin{smallmatrix}
        1 & 0 \\ 0 & p
    \end{smallmatrix}\!\big)\varGamma_0(N)$ is clear. For the other containment, note that $\varGamma_1(N)\backslash \varGamma_0(N)$ is represented by matrices of the form $\big(\!\begin{smallmatrix}
        a & b \\ N & d
    \end{smallmatrix}\!\big)\in\slz$. Thus it suffices to show that for each such matrix there exists a matrix $\delta\in\varGamma_1(N)$ such that $\big(\!\begin{smallmatrix}
        a & b \\ N & d
    \end{smallmatrix}\!\big)\big(\!\begin{smallmatrix}
        1 & 0 \\ 0 & p
    \end{smallmatrix}\!\big)\varGamma_0(N) = \delta \big(\!\begin{smallmatrix}
        1 & 0 \\ 0 & p
    \end{smallmatrix}\!\big)\varGamma_0(N)$; equivalently a matrix $\delta^\prime\in\varGamma_1(N)$ such that \[\begin{pmatrix}
        1 & 0 \\ 0 & p
    \end{pmatrix}^{-1}\delta^\prime\begin{pmatrix}
        a & b \\ N & d
    \end{pmatrix}\begin{pmatrix}
        1 & 0 \\ 0 & p
    \end{pmatrix}\in \varGamma_0(N).\]
    If $p\mid N$, then $\delta^\prime = \big(\!\begin{smallmatrix}
        dN+1 & -1 \\ -dN & 1
    \end{smallmatrix}\!\big)$ is a valid choice. Otherwise, if $p\nmid N$, then any $\delta^\prime = \big(\!\begin{smallmatrix}
        \ast & \ast \\ N & d^\prime
    \end{smallmatrix}\!\big)$ with $d^\prime\equiv 1\smod N$ and $d^\prime\equiv -a\smod p$ is a valid choice in this case.
\end{proof}
So in the definition of $T_p$, we may replace $\alpha$ with any matrix in the double coset $\varGamma_1(N)\big(\!\begin{smallmatrix}
    1 & 0 \\ 0 & p
\end{smallmatrix}\!\big)\varGamma_1(N)$.

We record the explicit action of $T_p$ on elements of $\mathcal M_k(\varGamma_1(N))$.

\begin{proposition}\label{prop: action of Tp}
    The operator $T_p\colon \mathcal M_k(\varGamma_1(N))\to \mathcal M_k(\varGamma_1(N))$ is given explicitly by 
    \[T_pf = \begin{cases}
        p^{k/2-1}\sum_{j=0}^{p-1} f|_k[\big(\!\begin{smallmatrix}
            1 & j \\ 0 & p
        \end{smallmatrix}\!\big)] & \text{if $p\mid N$,}\\
        p^{k/2-1}\sum_{j=0}^{p-1}f|_k[\big(\!\begin{smallmatrix}
            1 & j \\ 0 & p
        \end{smallmatrix}\!\big)] + p^{k/2-1}f|_k[\big(\!\begin{smallmatrix}
            m & n \\ N & p
        \end{smallmatrix}\!\big)\big(\!\begin{smallmatrix}
            p & 0 \\ 0 & 1
        \end{smallmatrix}\!\big)] & \text{if $p\nmid N$, where $mp-nN = 1$}.
    \end{cases}\]
\end{proposition}
\begin{proof}
    We find orbit representatives for $\varGamma_1(N)\backslash \varGamma_1(N)\big(\!\begin{smallmatrix}
        1 & 0 \\ 0 & p
    \end{smallmatrix}\!\big)\varGamma_1(N)$ by finding coset representatives for $\varGamma_3\backslash\varGamma_1(N)$, where $\varGamma_3 = \big(\!\begin{smallmatrix}
        1 & 0 \\ 0 & p
    \end{smallmatrix}\!\big)^{-1}\varGamma_1(N)\big(\!\begin{smallmatrix}
        1 & 0 \\ 0 & p
    \end{smallmatrix}\!\big)\cap \varGamma_1(N)$.

    Define the subgroups 
    \[\varGamma^0(p) = \cbr{\smallabcd\in\slz : \smallabcd \equiv \big(\!\begin{smallmatrix}
        \ast & 0 \\ \ast & \ast
    \end{smallmatrix}\!\big)\smod p}\] and $\varGamma_1^0(N,p) = \varGamma_1(N)\cap \varGamma^0(p)$. A short computation reveals that $\varGamma_3$ is equal to $\varGamma_1^0(N,p)$. We show that the coset representatives of $\varGamma_3\backslash \varGamma_1(N)$ are $\gamma_{2,j} = \big(\!\begin{smallmatrix}
        1& j \\ 0 & 1
    \end{smallmatrix}\!\big)$ for $0\leq j < p$.

    Given $\gamma_2 = \smallabcd\in\varGamma_1(N)$, we have $\gamma_2\in\varGamma_3\gamma_{2,j}$ if $\gamma_2\gamma_{2,j}^{-1}\in\varGamma_3$. Observe that $\gamma_2,\gamma_{2,j}\in\varGamma_1(N)$, so it suffices to find $j$ so that the upper right entry $b-ja$ of $\gamma_2\gamma_{2,j}^{-1}$ is some multiple of $p$.
    
    If $p\nmid a$, then $j = ba^{-1}\smod p$ is a valid choice. So in this case we find that $\gamma_{2,j} = \big(\!\begin{smallmatrix}
        1& j \\ 0 & 1
    \end{smallmatrix}\!\big)$ for $0\leq j < p$ form a complete set of coset representatives as needed (for each representative, one may also take $j + lp$ for any integer $l$ in place of $j$).
    
    If $p\mid a$, then $b-ja$ cannot be zero $\smod p$ for any $j$; if there was a choice of $j$ that made $b-ja$ zero, then $p$ would divide $b$ and hence also $\det \smallabcd = 1$, which is impossible. The cases where $\gamma_2\in\varGamma_1(N)$ with $p\mid a$ occur if and only if $p\nmid N$, in which case $\gamma_{2,j} = \big(\!\begin{smallmatrix}
        1& j \\ 0 & 1
    \end{smallmatrix}\!\big)$ for $0\leq j < p$ fail to represent $\varGamma_3\backslash\varGamma_1(N)$. To complete the set of coset representatives, set 
    \[\gamma_{2,\infty} = \begin{pmatrix}
        mp & n \\ N & 1
    \end{pmatrix}\quad \text{where $mp-nN = 1$}.\]
    So given $\gamma_2 = \smallabcd\in\varGamma_1(N)$ with $p\mid a$, we have $\gamma_2\gamma_{2,\infty}^{-1}\equiv \big(\!\begin{smallmatrix}
        \ast & 0 \\ \ast & \ast
    \end{smallmatrix}\!\big) \smod p$ as needed. It is also routine to check that the $\gamma_{2,j}$ represent distinct cosets.

    The corresponding orbit representatives of $\varGamma_1(N)\backslash \varGamma_1(N)\big(\!\begin{smallmatrix}
        1 & 0 \\ 0 & p
    \end{smallmatrix}\!\big)\varGamma_1(N)$ are given by $\alpha\gamma_{2,j}$ for each $j$: \[\alpha\gamma_{2,j} = \begin{pmatrix}
        1 & j \\ 0 & p
    \end{pmatrix}\text{ for $0\leq j < p$},\quad \alpha\gamma_{2,\infty} = \begin{pmatrix}
        m & n \\ N & p
    \end{pmatrix}\begin{pmatrix}
        p & 0 \\ 0 & 1
    \end{pmatrix}\text{ if $p\nmid N$}.\qedhere\]
\end{proof}

A related result, which we do not prove but will be used later, is the following:
\begin{lemma}\label{lem: miyake lem 4.5.11} 
    Let $N$ be a positive integer and $p$ a prime. Then 
    \[\varGamma_1(pN)\begin{pmatrix}
        1 & 0 \\ 0 & p
    \end{pmatrix}\varGamma_1(N) = \begin{cases}
        \coprod_{v = 0}^{p-1}\varGamma_1(pN)\big(\!\begin{smallmatrix}
            1 & 0 \\ 0 & p
        \end{smallmatrix}\!\big)\gamma_v& \text{if $p\mid N$},\\
        \coprod_{v = 0}^p\varGamma_1(pN)\big(\!\begin{smallmatrix}
            1 & 0 \\ 0 & p
        \end{smallmatrix}\!\big)\gamma_v& \text{if $p\nmid N$};
    \end{cases}\]
    here  $\gamma_v$ for $0\leq v<p$ is an element of $\varGamma_0(N)$ such that $\gamma_v \equiv \big(\!\begin{smallmatrix}
        1 & v \\ 0 & 1
    \end{smallmatrix}\!\big)\smod{p}$, and $\gamma_p$ for $p\nmid N$ is an element of $\varGamma_0(N)$ such that 
    \[\gamma_p \equiv\begin{cases}
        \big(\!\begin{smallmatrix}
            0 & -a \\ a^{-1} & 0
        \end{smallmatrix}\!\big)\smod{p}&\text{for $a$ coprime to $p$},\\
        \big(\!\begin{smallmatrix}
            1 & 0 \\ 0 & 1
        \end{smallmatrix}\!\big)\smod{N}.
    \end{cases}\]
\end{lemma}
See \cite[Lemma 4.5.11]{miyake} for details.

\begin{lemma}\label{lem: Fourier coefficients after Tp-ing}
    Let $f\in \mathcal M_k(\varGamma_1(N))$ have Fourier expansion 
    \[f(z) = \sum_{n=0}^\infty a_n(f)q^n,\quad q = e^{2\pi i z}.\] 
    Then $T_pf$ has Fourier expansion
    \begin{align*}
        (T_pf)(z) &= \sum_{n=0}^\infty a_{np}(f)q^n + \mathbf{1}_N(p)p^{k-1}\sum_{n=0}^\infty a_n(\abr{p}f)q^{np}\\
        &= \sum_{n=0}^\infty (a_{np}(f) + \mathbf{1}_N(p)p^{k-1}a_{n/p}(\abr{p}f))q^n,
    \end{align*}
    where $a_{n/p} = 0$ whenever $n/p$ is not an integer.
\end{lemma}
\begin{proof}
    Let $p\mid N$. For $0\leq j<p$, we have
    \begin{align*}
        p^{k/2-1}(f|_k[\big(\!\begin{smallmatrix}
            1 & j \\ 0 & p
        \end{smallmatrix}\!\big)])(z) &= p^{k-1}p^kf\bigg(\frac{z+j}{p}\bigg)\\
        &= \frac{1}{p}\sum_{n=0}^\infty a_n(f)e^{2\pi i n(z+j)/p}\\
        &= \frac{1}{p}\sum_{n=0}^\infty a_n(f)q_p^n\mu_p^{nj},
    \end{align*}
    where $q_p = e^{2\pi i z/p}$ and $\mu_p = e^{2\pi i/p}$. We use \cref{prop: action of Tp} and the observation that the geometric sum $\sum_{j=0}^{p-1} \mu_p^{nj}$ is equal to $p$ when $p\mid n$ and is zero otherwise to obtain
    \[(T_pf)(z) = p^{k/2-1}\sum_{j=0}^{p-1}(f|_k[\big(\!\begin{smallmatrix}
        1 & j \\ 0 & p
    \end{smallmatrix}\!\big)])(z) = \sum_{n\equiv 0\smod p}a_n(f)q_p^n = \sum_{n=0}^\infty a_{np}(f)q^n.\]

    When $p\nmid N$, the series expansion for $T_pf$ is the one given above plus the term 
    \begin{align*}
        p^{k/2-1}(f|_k[\big(\!\begin{smallmatrix}
            m & n \\ N & p
        \end{smallmatrix}\!\big)\big(\!\begin{smallmatrix}
            p & 0 \\ 0 & 1
        \end{smallmatrix}\!\big)])(z) &= p^{k/2-1}((\abr{p}f)|_k[\big(\!\begin{smallmatrix}
            p & 0 \\ 0 & 1
        \end{smallmatrix}\!\big)])(z)\\
        &= p^{k-1}z^{-k}(\abr{p}f)(pz)\\
        &= p^{k-1}\sum_{n=0}^\infty a_n(\abr{p}f)q^{np}.\qedhere
    \end{align*}
\end{proof}

\begin{corollary}\label{cor: fourier coeffs for eigenspace functions}
    If $f\in\mathcal M_k(N,\chi)$ has Fourier expansion $f(z) = \sum_{n=0}^\infty a_n(f)q^n$, with $q = e^{2\pi i z}$, then $T_pf$ has Fourier expansion
    \begin{align*}
        (T_pf)(z) &= \sum_{n=0}^\infty (a_{np}(f) + \mathbf{1}_N(p)\chi(p)p^{k-1}a_{n/p}(f))q^n,
    \end{align*}
    where $a_{n/p} = 0$ whenever $n/p$ is not an integer. This follows directly from the definition of the diamond operator.
\end{corollary}

\begin{lemma}\label{lem: properties of baby heckes}
    Let $d,e\in(\mathbb Z/N\mathbb Z)^\times$ and $p,q$ be prime. Then \begin{enumerate}[label = \textup{(\arabic*)}]
        \item $\abr{d}T_p = T_p\abr{d}$,
        \item $\abr{d}\abr{e} = \abr{e}\abr{d} = \abr{de}$, and 
        \item $T_pT_q = T_qT_p$.
    \end{enumerate}
\end{lemma}
\begin{proof}
    To show (1), let $\alpha = \big(\!\begin{smallmatrix}
        1 & 0 \\ 0 & p
    \end{smallmatrix}\!\big)$ and observe that $\gamma\alpha\gamma^{-1}\equiv \big(\!\begin{smallmatrix}
        1 & \ast \\ 0 & p
    \end{smallmatrix}\!\big)\smod N$ for any $\gamma\in\varGamma_0(N)$. If $\varGamma_1(N)\alpha\varGamma_1(N) = \bigcup_j \varGamma_1(N)\beta_j$, then by \cref{lem: replacement for T_n} and the fact that $\varGamma_1(N)$ is normal in $\varGamma_0(N)$ we have that 
    \[\varGamma_1(N)\alpha\varGamma_1(N) = \varGamma_1(N)\gamma\alpha\gamma^{-1}\varGamma_1(N) = \gamma\varGamma_1(N)\alpha\varGamma_1(N)\gamma^{-1} = \gamma\bigcup_j \varGamma_1(N)\beta_j\gamma^{-1} = \bigcup_j \varGamma_1(N)\gamma\beta_j\gamma^{-1}.\] Comparing the decompositions $\varGamma_1(N)\alpha\varGamma_1(N) = \bigcup_j \varGamma_1(N)\beta_j = \bigcup_j \varGamma_1(N)\gamma\beta_j\gamma^{-1}$ , we find that $\bigcup_j \varGamma_1(N)\beta_j\gamma = \bigcup_j \varGamma_1(N)\gamma\beta_j$. Thus for any $f\in\mathcal M_k(\varGamma_1(N))$ and any $\gamma\in\varGamma_0(N)$ with lower right entry $\delta\equiv d\smod N$, 
    \[\abr{d}T_pf = \sum_j f|_k[\beta_j\gamma] = \sum_j f|_k[\gamma\beta_j] = T_p\abr{d}.\]

    Because $\abr{d}$ and $T_p$ commute, it follows that the action of $T_p$ on $\mathcal M_k(\varGamma_1(N))$ respects the decomposition $\mathcal M_k(\varGamma_1(N)) = \bigoplus_\chi\mathcal M_k(N,\chi)$, and similarly for $\mathcal S_k(\varGamma_1(N)) = \bigoplus_\chi\mathcal S_k(N,\chi)$ and $\mathcal E_k(\varGamma_1(N)) = \bigoplus_\chi\mathcal E_k(N,\chi)$. That is, $T_p$ maps a $\chi$-eigenspace to itself.

    To see (2) and (3), it suffices to show that these equalities of operators hold on elements of $\mathcal M_k(N,\chi)$, for each $\chi$.

    Indeed, let $f\in\mathcal M_k(N,\chi)$. Then \[\abr{d}\abr{e}f = \abr{d}\chi(e)f = \chi(d)\chi(e)f = \chi(de)f = \abr{de}f,\] and similarly obtain that $\abr{d}\abr{e} = \abr{e}\abr{d}$.
    
    To see that $T_p$ and $T_q$ commute, we investigate the Fourier coefficients of $T_p(T_q(f))$ for $f\in\mathcal M_k(N,\chi)$. From \cref{cor: fourier coeffs for eigenspace functions}, we have that the $n$-th Fourier coefficient of $T_p(f)$ is $a_n(T_pf) = a_{np}(f)+ \chi(p)p^{k-1}a_{n/p}(f)$, where $a_{n/p}(f) = 0$ whenever $n/p$ is not an integer. Applying this formula twice yields the $n$-th Fourier coefficient of $T_p(T_q(f))$:
    \begin{align*}
        a_n(T_p(T_qf)) &= a_{np}(T_qf)+ \chi(p)p^{k-1}a_{n/p}(T_qf)\\
        &= a_{npq}(f) + \chi(q)q^{k-1}a_{np/q}(f)\\
        &\hspace{1em}+ \chi(p)p^{k-1}(a_{nq/p}(f) + \chi(q)q^{k-1}a_{n/pq}(f))\\
        &= a_{npq}(f) + \chi(q)q^{k-1}a_{np/q}(f) + \chi(p)p^{k-1}a_{nq/p}(f) + \chi(pq)(pq)^{k-1}a_{n/pq}(f). 
    \end{align*}
    Since the above expression is symmetric in $p$ and $q$, we obtain the result.
\end{proof}
\begin{definition}
    The \boldindex{Hecke operator} \boldindex{$T_n$}$\colon \mathcal M_k(\varGamma_1(N))\to \mathcal M_k(\varGamma_1(N))$ for $n\in\mathbb Z^+$ is defined inductively. Let $T_1 = \id_{\mathcal M_k(\varGamma_1(N))}$ (the identity operator), and above we had defined $T_pf = f|_k[\varGamma_1(N)\big(\!\begin{smallmatrix}
        1 & 0 \\ 0 & p
    \end{smallmatrix}\!\big)\varGamma_1(N)]$ for primes $p$. For prime powers, let
    \[T_{p^r} = T_pT_{p^{r-1}} - p^{k-1}\abr{p}T_{p^{r-2}}\quad \text{for $r\geq 2$}.\] Then for $n = \prod_i p_i^{r_i}$, define 
    \[T_n = \prod_i T_{p_i^{r_i}},\] where the previous lemma justifies the multiplicative notation above.
\end{definition}

By \cref{lem: properties of baby heckes}, the mapping $n\mapsto \abr{n}$ is totally multiplicative; that is, $\abr{nm} = \abr{n}\abr{m}$ for all positive integers $n,m$. The same lemma implies that the mapping $n\mapsto T_n$ is multiplicative; that is, for $n,m$ coprime, $T_{nm} = T_nT_m$.

\begin{definition}
    The \textbf{\textit{Hecke algebra}} \boldindex{$\mathcal R(N)$} is the polynomial ring 
    \[\mathcal R(N) = \mathbb{Z}[T_p, \abr{q} : p,q\in\mathbb P, q\nmid N].\]
    From previous results it follows that $\mathcal R(N)$ is a commutative unital ring.
\end{definition}

Often (and especially so in the next section), we will restrict the action of the Hecke operators to an arbitrary $\chi$-eigenspace. We use the same symbols to denote the restrictions of these operators, since it will be clear from context when we are taking the restriction.

In \cite{miyake}, the Hecke algebra is generated by Hecke operators $T(p), T(p,p)$ that are defined differently. The Hecke operators $T(p), T(p,p)$ are related to the Hecke operators $T_p, \abr{p}$ by the equations $T(p) = T_p$ and $T(p,p) = p^{k-2}\abr{p}$.

\begin{lemma}\label{lem: miyake lem 4.5.12}
    Let $K$ be a commutative unital ring, and assume that two sequences $\cbr{t_n}_{n=1}^\infty,\cbr{d_n}_{n=1}^\infty$ of elements of $K$ satisfy the following conditions:
    \begin{enumerate}[label=\textup{(\roman*)}]
        \item $t_1 =d_1 = 1$, and 
        \item $d_{mn} = d_md_n$ for any positive integers $m,n$.
    \end{enumerate} Then the following are equivalent:
    \begin{enumerate}[label=\textup{(\arabic*)}]
        \item If $\gcd(m,n) = 1$, then $t_{mn} = t_mt_n$ and \[t_pt_{p^e} = t_{p^{e+1}} + pd_pt_{p^{e-1}}\] for all prime numbers $p$ and all positive integers $e$.
        \item The formal Dirichlet series $\sum_{n=1}^\infty t_nn^{-s}$ has formal Euler product \[\sum_{n=1}^\infty t_nn^{-s} = \prod_{p\in\mathbb P}(1-t_pp^{-s} + pd_pp^{-2s})^{-1}.\]
        \item For any positive integers $m,n$, \[t_mt_n = \sum_{\substack{l>0\\l\mid \gcd(m,n)}}ld_lt_{mn/l^2}.\]
    \end{enumerate}
\end{lemma}
\begin{proof}
    We show that (1) implies (2). From the first condition of (1), we have formally that 
    \[\sum_{n=1}^\infty t_nn^{-s} = \prod_{p\in\mathbb P}\bigg(\sum_{e=0}^\infty t_{p^e}p^{-es}\bigg).\]
    From the second condition of (1), we have \[(1-t_pp^{-s} + pd_pp^{-2s})\bigg(\sum_{e=0}^\infty t_{p^e}p^{-es}\bigg) = 1,\] from which (2) follows. 

    We show that (2) implies (3). Formal Dirichlet series and formal Euler products are elements of the ring of formal power series $K[[p^{-s} : p\in\mathbb P]]$. Finding inverses in $K[[p^{-s}]]$ of the elements $1-t_pp^{-s} + pd_pp^{-2s}$ recursively (cf. \cite[Exercise 7.2.3(c)]{df}), and comparing coefficients of \[\sum_{n=1}^\infty t_nn^{-s} = \prod_{p\in\mathbb P}(1-t_pp^{-s} + pd_pp^{-2s})^{-1}\] reveals that the sequence $\cbr{t_n}_{n=1}^\infty$ is multiplicative. That is, if $\gcd(m,n) =1$, then $t_{mn} = t_mt_n$.
    
    Let $m = \prod_{p\in\mathbb P}p^e$ and $n = \prod_{p\in\mathbb P}p^f$ be the prime factorizations of $m$ and $n$. Then \[\sum_{\substack{l>0\\l\mid\gcd(m,n)}} ld_lt_{mn/l^2} = \prod_{p\in\mathbb P} \bigg(\sum_{0\leq g\leq\min\{e,f\}}p^gd_{p^g}t_{p^{e+f-2g}}\bigg).\] Therefore it suffices to prove that (2) implies (3) in the case that $m$ and $n$ are powers of the same prime $p$. By assumption, we have 
    \[(1-t_pp^{-s} + pd_pp^{-2s})^{-1} = \sum_{e=0}^\infty t_{p^e}p^{-es}.\] Let $\tau_p$ and $\delta_p$ be two indeterminates over $\mathbb Q$, and define a (unital) ring homomorphism $\psi\colon \mathbb Z[\tau_p,\delta_p]\to K$ by $\psi(\tau_p) = t_p$ and $\psi(\delta_p) = d_p$. Define the elements $\tau_{p^e}$ and $\delta_{p^e}$ of $\mathbb Z[\tau_p,\delta_p]$ by $\delta_{p^e} = (\delta_p)^e$ and by the formal power series equality 
    \begin{equation}\label{eqn: tau power series expansion}
        \sum_{e=0}^\infty \tau_{p^e}p^{-es} = (1-\tau_pp^{-s} + p\delta_pp^{-2s})^{-s}.
    \end{equation} It follows that $\psi(\tau_{p^e}) = t_{p^e}$ and $\psi(\delta_{p^e}) = d_{p^e}$.

    Let $u,v$ be indeterminates over $\mathbb Q$, and define a ring homomorphism $\phi\colon \mathbb Z[\tau_p,\delta_p]\to \mathbb Q[u,v]$ by $\phi(\tau_p) = u+v$ and $\phi(\delta_p) = uv/p$. Since $u+v$ and $uv/p$ are algebraically independent over $\mathbb Q$, $\phi$ is injective. Viewing $\mathbb Z[\tau_p,\delta_p]$ as a subring of $\mathbb Q[u,v]$, we may factor $1-\tau_pp^{-s}+p\delta_pp^{-2s}$ into $(1-up^{-s})(1-vp^{-s})$. Inverting this element and comparing coefficients in \eqref{eqn: tau power series expansion}, we find that
    \[\tau_{p^e} = \sum_{i+j=e}u^iv^j = (u^{e+1}-v^{e+1})/(u-v).\]

    Assume that $0\leq e\leq f$. Then \begin{align*}
        \tau_{p^e}\tau_{p^f} &= \bigg(\sum_{i+j=e}u^iv^j\bigg)\cdot (u^{f+1}-v^{f+1})/(u-v)\\
        &= \bigg(u^{f+1}\sum_{j=0}^eu^{e-j}v^j - v^{f+1}\sum_{j=0}^eu^{j}v^{e-j}\bigg)\bigg/(u-v)\\
        &=\sum_{g=0}^e u^gv^g(u^{e+f-2g+1}-v^{e+f-2g+1})/(u-v)\\
        &=\sum_{g=0}^e p^g\delta_{p^g}\tau_{p^{e+f-2g}}.
    \end{align*}
    By taking $\psi$ on both sides of the above equation, obtain \[t_{p^e}t_{p^f} = \sum_{g=0}^ep^gd_{p^g}t_{p^{e+f-2g}}\] as desired. So in all cases, (2) implies (3).

    To see that (3) implies (1), observe that (1) is a special case of (3), with $m = p$ and $n = p^e$.
\end{proof}
\begin{theorem}\label{thm: miyake thm 4.5.13}
    We have \begin{enumerate}[label=\textup{(\arabic*)}]
        \item $T_mT_n = \sum_{\substack{l>0 \\ l\mid \gcd(m,n)\\ \gcd(l,N)=1}}l^{k-1}\abr{l}T_{mn/l^2}$, and
        \item the formal Dirichlet series $\sum_{n=1}^\infty T_nn^{-s}$ has the formal Euler product \[\sum_{n=1}^\infty T_nn^{-s} = \prod_{p\nmid N}(1-T_np^{-s} + (p^{k-2}\abr{p})p^{1-2s})^{-1}\cdot \prod_{p\mid N}(1-T_pp^{-s})^{-1}\]
    \end{enumerate}
\end{theorem}
\begin{proof}
    Apply the previous lemma with 
    \[K = \mathcal R(N),\quad t_n = T_n,\quad \text{and}\quad d_n =\begin{cases}
        n^{k-2}\abr{n} & \text{if $\gcd(n,N) = 1$},\\
        0 & \text{if $\gcd(n,N)\neq 1$}
    \end{cases}\] to obtain the result.
\end{proof}

We record one result (without proof) which will be useful in the next result. For $f\in\mathcal M_k(N,\chi)$, we have
\[(T_nf)(z) = n^{k-1}\sum_{\substack{d>0 \\ d\mid n \\ ad= n}}\sum_{b=0}^{d-1}\chi(a)d^{-k}f((az+b)/d).\]

The following result is similar to \cref{lem: Fourier coefficients after Tp-ing}, but for $T_n$ acting on $\mathcal M_k(N,\chi)$.
\begin{lemma}\label{lem: miyake lem 4.5.14}
    Let $f$ be an element of $\mathcal M_k(N,\chi)$, and let \[
        f(z) = \sum_{n=0}^\infty c_ne^{2\pi i n z},\quad
        (T_mf)(z) = \sum_{n=0}^\infty b_ne^{2\pi i nz}\] be Fourier series expansions. Then \[b_n = \sum_{\substack{d>0 \\ d\mid \gcd(m,n)}}\chi(d)d^{k-1}c_{mn/d^2}.\]
\end{lemma}
\begin{proof}
    Observe that $f((az+b)/d) = \sum_{n=0}^\infty c_ne^{2\pi in(az+b)/d}$. Then by the explicit action of $T_n$ on elements of $\mathcal M_k(N,\chi)$ (which appears just before this lemma), we have
    \[(T_mf)(z) = m^{k-1}\sum_{\substack{d>0\\d\mid n}}\sum_{b=0}^{d-1}\chi(m/d)d^{-k}\sum_{n=0}^\infty c_ne^{2\pi in((m/d)z+b)/d}.\] Then use the fact that $\sum_{b=0}^{d-1}e^{2\pi i n b/d}$ is equal to $d$ if $d\mid n$ and is $0$ otherwise to simplify the above expression into 
    \[\sum_{n=0}^\infty \sum_{\substack{d>0\\d\mid \gcd(m,n)}}\chi(m/d)(m/d)^{k-1}c_ne^{2\pi i(nm/d^2)z}.\] Change summation variables to obtain 
    \[\sum_{n=0}^\infty \sum_{a>0\\ a\mid \gcd(m,n)}\chi(a)a^{k-1}c_{dn/a}e^{2\pi inz},\] and with $ad = m$, it follows that $c_{dn/a} = c_{mn/a^2}$, and the result follows.
\end{proof}

\begin{lemma}\label{lem: miyake lem 4.5.15}
    Let $f(z) = \sum_{n=0}^\infty c_ne^{2\pi i n z}$ be an element of $\mathcal M_k(N,\chi)$, and $P$ a set of prime numbers where $T_pf = t_pf$ for $p\in P$, $t_p\in\mathbb C$. Then 
    \begin{enumerate}[label=\textup{(\arabic*)}]
        \item If all prime factors of a positive integer $m$ lie in $P$, then $f$ is an eigenfunction of $T_m$. In this case, let $T_mf = t_mf$ for some $t_m\in\mathbb C$, from which it follows that $c_m = t_mc_1$.
        \item We have $L(s,f) = \prod_{p\in P}(1-t_pp^{-s} + \chi(p)p^{k-1-2s})^{-1}\cdot \sum_n^\prime c_nn^{-s}$, where the summation $\sum_n^\prime$ is taken over the positive integers coprime to every element of $P$.
    \end{enumerate}
\end{lemma}
\begin{proof}
    The first claim in (1) is clear. Let $n$ be a positive integer coprime to $m$. Comparing the $n$-th Fourier coefficients of $T_mf$ and $t_mf$, we obtain $t_mc_n = c(mn)$ by \cref{lem: miyake lem 4.5.14}. By taking $n = 1$, the second claim in (1) follows.

    By the equality $t_mc_n = c(mn)$, we have the formal equality
    \[L(s,f) = \bigg(\sideset{}{^{\prime\prime}}\sum_{m}t_mm^{-s}\bigg)\cdot \bigg(\sideset{}{^\prime}\sum_nc_nn^{-s}\bigg),\] where the summation $\sum_m^{\prime\prime}$ is taken over $1$ and the positive integers whose prime factors are all contained in $P$, and the summation $\sum_n^\prime$ is taken over the positive integers copimre to all primes in $P$. Then by \cref{thm: miyake thm 4.5.13}(1), 
    \[t_nt_m = \sum_{\substack{l > 0\\ l\mid \gcd(m,n)\\\gcd(l,N) = 1}}\chi(l)l^{k-1}t_{mn/l^2}\] for $m,n$ in the indexing sets for the summations $\sum_m^{\prime\prime}$ and $\sum_n^\prime$, respectively. So formally obtain
    \[\sideset{}{^{\prime\prime}}\sum_{m}t_m m^{-s} = \prod_{p\in P}(1-t_pp^{-s} + \chi(p)p^{k-1-2s})^{-1}\] by a similar argument as in \cref{lem: miyake lem 4.5.12}. Furthermore, if $c_n = O(n^\alpha)$ for some $\alpha$, then the equality above holds on $\Re(s)>\alpha + 1$. If $f$ is a cusp form, then we have this growth condition on $c_n$ due to \cref{cor: cuspform growth corollary}. We do not prove the result when $f$ is not a cusp form.
\end{proof}

The next lemma gives a necessary and sufficient condition for a modular form $f\in\mathcal M_k(N,\chi)$ to be a simultaneous eigenfunction of all of the Hecke operators $T_n$, in terms of being able to write down an Euler product for the $L$-function for $f$. More importantly, the lemma shows that the eigenvalues of such an eigenfunction are proportional to the Fourier coefficients of the modular form.
\begin{theorem}\label{thm: miyake thm 4.5.16}
    Let $f(z) = \sum_{n=0}^\infty c_ne^{2\pi i nz}$ be a nonzero element of $\mathcal M_k(N,\chi)$. Then the following are equivalent: \begin{enumerate}[label=\textup{(\arabic*)}]
        \item $f(z)$ is a common eigenfunction of every Hecke operator $T_n$;
        \item $c_1\neq 0$ and 
        \[L(s,f) = c_1\prod_{p\in\mathbb P} (1-t_pp^{-s} + \chi(p)p^{k-1-2s})^{-1},\quad t_n = \frac{c_n}{c_1}.\]
    \end{enumerate}
    Moreover, if $f(z)$ satisfies the above conditions, then $T_nf = t_nf$ for all $n\geq 1$.
\end{theorem}
\begin{proof}
    Assume that (1) holds, and let $T_nf = t_nf$. By \cref{lem: miyake lem 4.5.15}, we have $c_n = t_nc_1$ for $n\geq 1$. Then if $c_1 = 0$, $f(z) = c_0$. Since $k\geq 1$, we have in this case that $f(z) = 0$, which cannot happen by assumption. Therefore $c_1\neq 0$. Obtain (2) from \cref{lem: miyake lem 4.5.15}(2), taking $P$ to be the set of all prime numbers $\mathbb P$.

    Now assume that (2) holds, and let $t_n = c_n/c_1$ for $n\geq 1$. Then 
    \[\sum_{n=1}^\infty t_nn^{-s} = \prod_{p\in\mathbb P}(1-t_pp^{-s} + \chi(p)p^{k-1-2s})^{-1}.\]
    By \cref{lem: miyake lem 4.5.12}, 
    \[t_mt_n = \sum_{\substack{d > 0\\d\mid \gcd(m,n)}}d^{k-1}\chi(d)t_{mn/d^2},\] and multiplying both sides of this equality by $c_1$ we obtain 
    \[t_mc_1 = \sum_{\substack{d > 0\\d\mid \gcd(m,n)}}d^{k-1}\chi(d)c_{mn/d^2}.\] But the expression on the right hand side of this equality is equal to the $n$-th Fourier coefficient of $T_mf$ by \cref{lem: miyake lem 4.5.14}. Let $b_0$ denote the constant term of the Fourier expansion of $T_mf$. Then $(T_mf)(z) - t_mf = b_0-c_0$, and since $T_mf - t_mf$ is an element of $\mathcal M_k(N,\chi)$ with $k\geq 1$, we have that $b_0 = c_0$. Hence $T_mf = t_mf$ as desired.
\end{proof}

\subsection{The Petersson inner product, adjoints of Hecke operators}\label{section petersson inner product}
We study the space $\mathcal S_k(\varGamma_1(N))$ by endowing it with an inner product. Define the \boldindex{hyperbolic measure} $\mu$ on $\mathcal H$ by $\dd\mu(z) = \dd x \dd y / y^2$ for $z = x+iy\in\mathcal H$ (we will sometimes suppress the $(z)$). This measure is invariant under the action of $\glrp$ on $\mathcal H$: For $\alpha \in \glrp$ and $z=x+iy \in \mathcal H$, let $\alpha(z) = \sigma(x,y) + i\tau(x,y)$. A few computations using Wirtinger derivatives reveal that $\big|\frac{\partial (\sigma,\tau)}{\partial(x,y)}\big| = \abs{\dv{\alpha}{z}}^2 = \big(\det\alpha/\abs{j(\alpha,z)}^2\big)^2$ and $\tau(z) = y\det\alpha/\abs{j(\alpha,z)}^2$. Then for a measurable set $A$ in $\mathcal H$, \[\mu(\alpha(A)) = \int_{\alpha(A)} \frac{\dd\sigma\dd\tau}{\tau^2} = \int_{A}\bigg|\frac{\partial (\sigma,\tau)}{\partial(x,y)}\bigg|\frac{\dd x\dd y}{(\tau(x,y))^2} = \int_A \frac{(\det\alpha)^2}{\abs{j(\alpha,z)}^4}\frac{\abs{j(\alpha,z)}^4\dd x \dd y}{(\det\alpha)^2y^2} =\int_A \frac{\dd x \dd y}{y^2} = \mu(A)\]
as desired. In particular $\dd \mu$ is $\slz$-invariant. Since $\mathbb Q\cup\cbr{\infty}$ is countable, its measure is zero, so we may integrate over the extended upper half plane $\mathcal H^\ast$ with respect to $\mu$.

Recall that a fundamental domain for $\mathcal H^\ast$ under the action of $\slz$ is given by 
\[\mathcal D^\ast = \cbr{z\in\mathcal H : \abs{\Re(z)}\leq 1/2,\abs{z}\geq 1}\cup\cbr{\infty};\] that is, any point of $\mathcal H$ is sent to a point in $\mathcal D$ by a suitable element of $\slz$, which is unique for most points of $\mathcal H$ (there are a few cases involving points on the boundary of $\mathcal D$). Every point $s\in\mathbb Q\cup\cbr{\infty}$ may be sent to $\infty$ by suitable elements of $\slz$. We show that the integral $\int_{\mathcal D^\ast} \varphi(\alpha(z))\dd \mu$ converges for any continuous, bounded function $\varphi\colon\mathcal H\to\mathbb C$ and any $\alpha\in\slz$: Choose $\gamma\in\slz$ such that $\gamma^{-1}\mathcal D^\ast$ is a compact set. By the invariance of the measure under the action of $\slz$, we have
\[\int_{\mathcal D^\ast} \varphi(\alpha(z))\dd \mu = \int_{\mathcal \gamma^{-1} D^\ast} \varphi(\alpha(\gamma(z)))\dd \mu, \] from which it follows that the integral converges.

Let $\varGamma$ be a congruence subgroup, and let $\cbr{\alpha_j}$ be a set of representatives for the coset space $\cbr{\pm I}\varGamma\backslash\slz$; that is, we have the disjoint union $\slz = \bigsqcup_j \cbr{\pm I}\varGamma \alpha_j$. If $\varphi$ is $\varGamma$-invariant, then the sum $\sum_j\int_{\mathcal D^\ast}\varphi(\alpha_j(z))\dd\mu$ is independent of the choice of coset representatives $\alpha_j$. Since $\dd\mu$ is $\slz$-invariant, the sum is equal to $\int_{\bigcup_j \alpha_j(\mathcal D^\ast)}\varphi(z)\dd\mu$. Since $\bigcup_j\alpha_j(\mathcal D^\ast)$ represents the modular curve $X(\varGamma)$ up to identification of boundaries, we define 
\[\int_{X(\varGamma)}\varphi(z)\dd\mu = \int_{\bigcup_j\alpha_j(\mathcal D^\ast)}\varphi(z)\dd\mu = \sum_j\int_{\mathcal D^\ast}\varphi(\alpha_j(z))\dd\mu.\] In particular the volume of $X(\varGamma)$ is given by $V_\varGamma = \int_{X(\varGamma)}\dd\mu = [\slz : \cbr{\pm I}\varGamma]V_{\slz}$.

Observe that for any $f,g\in\mathcal S_k(\varGamma)$ for a congruence subgroup $\varGamma$, the function $\varphi(z) = f(z)\overline{g(z)}(\Im(z))^k$ for $z\in\mathcal H$ is continuous, and more importantly, $\varGamma$-invariant. Indeed, for any $\gamma\in\varGamma$, we have
\begin{align*}
    \varphi(\gamma(z)) &= f(\gamma(z))\overline{g(\gamma(z))}(\Im(\gamma(z)))^k\\
    &= (f|_k[\gamma])(z)j(\gamma,z)^k\overline{(g|_k[\gamma])(z)}\overline{j(\gamma,z)}^k(\Im(z))^k\abs{j(\gamma,z)}^{-2k}\\
    &= (f|_k[\gamma])(z)\overline{(g|_k[\gamma])(z)}(\Im(z))^k\\
    &= \varphi(z).
\end{align*} To show that $\varphi$ is bounded on $\mathcal H$, we show that $\varphi$ is bounded on $\bigcup_j\alpha_j(\mathcal D)$, which is a finite union. Therefore it suffices to show that for any $\alpha\in\slz$, the function $\varphi\circ\alpha$ is bounded on $\mathcal D$. On any compact subset of $\mathcal D$, continuity of $\varphi\circ\alpha$ implies boundedness. For the neighborhoods $\cbr{\Im(z)> M}$ of $i\infty$, note that the Fourier expansions
\[(f|_k[\alpha])(z) = \sum_{n =1}a_n(f|_k[\alpha])q_h^n,\quad (g|_k[\alpha])(z) = \sum_{n =1}a_n(g|_k[\alpha])q_h^n\quad\text{for $q_h = e^{2\pi i z /h}$ for some $h\in\mathbb Z^+$}\] are of order $O(q_h)$ as $\Im(z)\to\infty$. It follows that $\varphi(\alpha(z)) = (f|_k[\gamma])(z)\overline{(g|_k[\gamma])(z)}(\Im(z))^k = O(q_h)^2(\Im(z))^k$. Moreover, since $\abs{q_h} = e^{-2\pi \Im(z)/h}$, $\varphi(\alpha(z))\to 0$ as $\Im(z)\to\infty$, from which it follows $\varphi\circ\alpha$ is bounded on the neighborhoods $\cbr{\Im(z)> M}$ of $i\infty$ of $\mathcal D$, hence on all of $\mathcal D$.
\begin{definition}
    The \boldindex{Petersson inner product} $\abr{\cdot,\cdot}_\varGamma\colon \mathcal S_k(\varGamma)\times \mathcal S_k(\varGamma)\to\mathbb C$ for a congruence subgroup $\varGamma$ is given by 
    \[\abr{f,g}_\varGamma = \frac{1}{V_\varGamma}\int_{X(\varGamma)}f(z)\overline{g(z)}(\Im(z))^k\dd\mu.\] We omit the subscript ${}_\varGamma$ in $\abr{\cdot,\cdot}_\varGamma$ when it is clear from context.
\end{definition}
One can check that the Petersson inner product is a Hermitian inner product. If we have the containment $\varGamma^\prime\subset \varGamma$ of congruence subgroups, then $\abr{\cdot,\cdot}_{\varGamma^\prime} = \abr{\cdot,\cdot}_\varGamma$ on $\mathcal S_k(\varGamma)$: If $\cbr{\beta_m}$ is a set of coset representatives for $\cbr{\pm I}\varGamma\backslash\slz$ and $\cbr{\gamma_{j}}$ is a set of coset representatives for $\cbr{\pm I}\varGamma^\prime\backslash\cbr{\pm I}\varGamma$, then $\cbr{\gamma_j\beta_m}$ form a set of coset representatives for $\cbr{\pm I}\varGamma^\prime\backslash \slz$, from which we obtain
\begin{multline*}
    \frac{1}{V_{\varGamma}}\int_{X(\varGamma)}f(z)\overline{g(z)}(\Im(z))^k\dd\mu = \frac{[\cbr{\pm I}\varGamma:\cbr{\pm I}\varGamma^\prime]}{V_{\varGamma^\prime}}\sum_{m}\int_{\mathcal D^\ast}f(\beta(z))\overline{g(\beta(z))}(\Im(\beta(z)))^k\dd\mu \\ = \frac{1}{V_{\varGamma^\prime}}\sum_{j,m}\int_{\mathcal D^\ast} f((\gamma\beta)(z))\overline{g((\gamma\beta)(z))}(\Im((\gamma\beta)(z)))^k\dd\mu = \frac{1}{V_{\varGamma^\prime}}\int_{X(\varGamma^\prime)}f(z)\overline{g(z)}(\Im(z))^k\dd\mu
\end{multline*} as expected.

Recall that if $T$ is a linear operator on the inner product space $V$, then there exists a unique operator $T^\ast$ on $V$ called the adjoint of $T$, that satisfies $\abr{Tv,w} = \abr{v,T^\ast w}$ for all $v,w\in V$. If $T$ commutes with $T^\ast$, we say that $T$ is a normal operator. Give $\mathcal S_k(\varGamma_1(N))$ the Petersson inner product. We show that the Hecke operators $\abr{n}$ and $T_n$ for $n$ coprime to $N$ are normal operators.

Let $\varGamma$ be a congruence subgroup, and write $\slz = \bigcup_j\cbr{\pm I}\varGamma \alpha_j$ for some representatives $\alpha_j$ of the coset space $\cbr{\pm I}\varGamma\backslash \slz$. If $\alpha\in\glqp$, then the map $\mathcal H\to\mathcal H$ given by $z \mapsto \alpha(z)$ induces a bijection $\alpha^{-1}\varGamma\alpha\backslash \mathcal H^\ast\to X(\varGamma)$. It follows that the union $\bigcup_j \alpha^{-1}\alpha_j(\mathcal D^\ast)$ up to some boundary identifications is in bijection with $\alpha^{-1}\varGamma\alpha\backslash \mathcal H^\ast$. For continuous, bounded, $\alpha^{-1}\varGamma\alpha$-invariant functions $\varphi\colon \mathcal H\to \mathbb C$ define
\[\int_{\alpha^{-1}\varGamma\alpha\backslash \mathcal H^\ast}\varphi(z)\dd\mu = \sum_j \int_{\mathcal D^\ast}\varphi(\alpha^{-1}\alpha_j(z))\dd\mu.\]

\begin{lemma}\label{lem: properties of integral over fund.domain and glqp}
    Let $\varGamma$ be a congruence subgroup and let $\alpha\in\glqp$.\begin{enumerate}[label = \textup{(\arabic*)}]
        \item If $\varphi\colon\mathcal H\to\mathbb C$ is continuous, bounded, and $\varGamma$-invariant, then \[\int_{\alpha^{-1}\varGamma\alpha\backslash \mathcal H^\ast}\varphi(\alpha(z))\dd\mu = \int_{X(\varGamma)}\varphi(z)\dd\mu.\]
        \item If $\alpha^{-1}\varGamma\alpha\subset \slz$, then $V_{\alpha^{-1}\varGamma\alpha} = V_\varGamma$ and $[\slz : \alpha^{-1}\varGamma\alpha] = [\slz : \varGamma]$.
        \item There exist $\beta_1,\dots,\beta_n\in \glqp$ with $n=[\varGamma : \alpha^{-1}\varGamma\alpha\cap\varGamma] = [\varGamma : \alpha\varGamma\alpha^{-1}\cap \varGamma]$, such that \[\varGamma\alpha\varGamma = \bigsqcup_j\varGamma \beta_j = \bigsqcup_j\beta_j\varGamma.\]
    \end{enumerate}
    \begin{proof}
        From the discussion above, (1) follows immediately. The first equality in (2) follows from (1), and the second equality is a consequence of the equation $V_\varGamma = [\slz : \cbr{\pm I}\varGamma]V_{\slz}$ and the observation that $-I\in\alpha^{-1}\varGamma\alpha$ if and only if $-I\in\varGamma$.

        To obtain (3), we apply (2) with $\alpha\varGamma\alpha^{-1}\cap \varGamma$ in place of $\varGamma$ to obtain 
        \[[\slz : \alpha^{-1}\varGamma\alpha\cap \varGamma] = [\slz : \alpha\varGamma\alpha^{-1}\cap\varGamma],\] from which we deduce that $[\varGamma : \alpha^{-1}\varGamma\alpha\cap \varGamma] = [\varGamma : \alpha\varGamma\alpha^{-1}\cap\varGamma]$. Thus there exist coset representatives $\gamma_1,\dots,\gamma_n$ and $\eta^{-1}_1,\dots,\eta^{-1}_n$ in $\varGamma$ for the coset spaces $(\alpha^{-1}\varGamma\alpha\cap \varGamma)\backslash\varGamma$ and $(\alpha\varGamma\alpha^{-1}\cap \varGamma)\backslash\varGamma$ respectively. That is, 
        \[\varGamma = \bigsqcup_j(\alpha^{-1}\varGamma\alpha\cap \varGamma)\gamma_j = \bigsqcup_j(\alpha\varGamma\alpha^{-1}\cap \varGamma)\eta_j^{-1}.\]
        In \cref{lem: double coset coset/orbit spaces}, set $\varGamma_1 = \varGamma_2 = \varGamma$ to obtain 
        \[\varGamma\alpha\varGamma = \bigsqcup_j\varGamma\alpha\gamma_j\quad\text{and}\quad \varGamma\alpha^{-1}\varGamma = \bigsqcup_j \varGamma\alpha^{-1}\eta^{-1}_j.\] The second equation implies that $\varGamma\alpha\varGamma = \bigsqcup_j\eta_j\alpha\varGamma$. For each $j$, the intersection $\varGamma\alpha\gamma_j\cap\eta_j\alpha\varGamma$ is nonempty, since otherwise we would have $\varGamma\alpha\gamma_j\subset\bigsqcup_{i\neq j}\eta_i\alpha\varGamma$. Multiplying this equality from the right by $\varGamma$ gives $\varGamma\alpha\varGamma\subset \bigsqcup_{i\neq j}\eta_i\alpha\varGamma$, a contradiction. So for each $j$, choose some $\beta_j\in \varGamma\alpha\gamma_j\cap\eta_j\alpha\varGamma$, from which we obtain $\varGamma\alpha\varGamma = \bigsqcup_j\varGamma \beta_j = \bigsqcup_j\beta_j\varGamma$.
    \end{proof}
\end{lemma}

The following proposition will be used to compute adjoints of the Hecke operators.
\begin{proposition}\label{prop: computing adjoints}
    Let $\varGamma$ be a congruence subgroup and let $\alpha\in\glqp$. \begin{enumerate}[label = \textup{(\arabic*)}]
        \item If $\alpha^{-1}\varGamma\alpha\subset \slz$, then for all $f\in\mathcal S_k(\varGamma)$ and $g\in\mathcal S_k(\alpha^{-1}\varGamma\alpha)$ \[\abr{f|_k[\alpha],g}_{\alpha^{-1}\varGamma\alpha} = \abr{f,g|_k[\det(\alpha)\alpha^{-1}]}_\varGamma.\]
        \item For all $f,g\in\mathcal S_k(\varGamma)$, \[\abr{f|_k[\varGamma\alpha\varGamma],g}  = \abr{f,g|_k[\varGamma\det(\alpha)\alpha^{-1}\varGamma]}.\]
        We have $|_k[\varGamma\alpha\varGamma]^\ast = |_k[\varGamma\det(\alpha)\alpha^{-1}\varGamma]$. If $\alpha^{-1}\varGamma\alpha = \varGamma$, then $|_k[\alpha]^\ast = |_k[\det(\alpha)\alpha^{-1}]$.
    \end{enumerate}
\end{proposition}
\begin{proof}
    To prove (1), note that $\alpha^\prime$ acts on elements of $\mathcal H^\ast$ in the same way as $\alpha^{-1}$, and apply \cref{lem: properties of integral over fund.domain and glqp}(1) to obtain
    \begin{align*}
        \int_{\alpha^{-1}\varGamma\alpha\backslash\mathcal H^\ast}(f|_k[\alpha])(z)\overline{g(z)}(\Im(z))^k\dd \mu &= \int_{\alpha^{-1}\varGamma\alpha\backslash\mathcal H^\ast}(\det\alpha)^{k/2}j(\alpha,z)^{-k}f(\alpha(z))\overline{g(z)}(\Im(z))^k\dd \mu \\
        &= \int_{X(\varGamma)}f(z)(\det\alpha)^{k/2}j(\alpha,\alpha^\prime(z))^{-k}\overline{g(\alpha^\prime(z))}(\Im(\alpha^\prime(z)))^k\dd \mu
        \intertext{Recall the identities $j(\alpha\alpha^\prime,z) = j(\alpha, \alpha^\prime(z))j(\alpha^\prime,z)$ and $\Im(\alpha^\prime(z)) = (\det\alpha^\prime)\Im(z)/\abs{j(\alpha^\prime,t)}^2$, and observe that $\det\alpha^\prime = \det\alpha$. Continuing with the computation gives}
        &= \int_{X(\varGamma)}f(z)\overline{(\det\alpha)^{k/2}j(\alpha^\prime,z)^{-k}g(\alpha^\prime(z))}(\Im(z))^k\dd \mu\\
        &= \int_{X(\varGamma)}f(z)\overline{(g|_k[\alpha^\prime])(z)}(\Im(z))^k\dd\mu.
    \end{align*} 
    Since $V_{\alpha^{-1}\varGamma\alpha} = V_\varGamma$ by \cref{lem: properties of integral over fund.domain and glqp}(2), the result follows.
    
    To show part (2), from \cref{lem: properties of integral over fund.domain and glqp}(3), there exist $\cbr{\beta_j}$ for which $\varGamma\alpha\varGamma = \bigsqcup_j \beta_j\varGamma$, from which we deduce that $\varGamma\alpha^\prime\varGamma = \bigsqcup_j\varGamma(\det\beta_j)\beta_j^{-1}$. Then apply part (1) of this lemma with each $\beta_j\varGamma\beta_j^{-1}\cap\varGamma$ in place of $\varGamma$ in the following:
    \begin{align*}
        \abr{f|_k[\varGamma\alpha\varGamma],g}_\varGamma &= (\det\alpha)^{k/2-1}\sum_j\abr{f|_k[\beta_j],g}_{\varGamma\cap \beta_j^{-1}\varGamma\beta_j}\\
        &= (\det\alpha^\prime)^{k/2-1}\sum_j\abr{f,g|_k[(\det\beta_j)\beta_j^{-1}]}_{\beta_j\varGamma\beta_j^{-1}\cap \varGamma}\\
        &= \abr{f,g|_k[\varGamma\alpha^\prime\varGamma]}_{\varGamma}.\qedhere
    \end{align*} 
\end{proof}
We obtain the adjoints of the Hecke operators without much trouble using the above proposition.
\begin{theorem}\label{thm: the actual adjoints}
    For $p\nmid N$, the Hecke operators $\abr{p},T_p\colon \mathcal S_k(\varGamma_1(N))\to \mathcal S_k(\varGamma_1(N))$ have adjoints \[\abr{p}^\ast = \abr{p}^{-1}\quad\text{and}\quad T_p^\ast = \abr{p}^{-1}T_p.\]  Furthermore, one can show that $T_n^\ast = \abr{n}^{-1}T_n$ for $n$ coprime to $N$. It follows that the Hecke operators $\abr{n},T_n$ for $n$ coprime to $N$ are normal operators.
\end{theorem}
\begin{proof}
    Let $f,g\in\mathcal S_k(\varGamma_1(N))$. Since $\varGamma_1(N)$ is normal in $\varGamma_0(N)$, we have by \cref{prop: computing adjoints}(1) that $\abr{p}^\ast $ is given by $|_k[\alpha]^\ast = |_k[\alpha^{-1}]$, where $\alpha$ is any element of $\varGamma_0(N)$ such that $\alpha\equiv \big(\!\begin{smallmatrix}
        \ast & \ast \\ 0 & p
    \end{smallmatrix}\!\big)\smod N$. But the action of $\abr{p}^{-1}$ is exactly $|_k[\alpha^{-1}]$.
    
    From \cref{prop: computing adjoints}(2), we have $T_p^\ast$ given by \[|_k[\varGamma_1(N)\big(\!\begin{smallmatrix}
        1 & 0 \\ 0 & p
    \end{smallmatrix}\!\big)\varGamma_1(N)]^\ast = |_k[\varGamma_1(N)\big(\!\begin{smallmatrix}
        p & 0 \\ 0 & 1
    \end{smallmatrix}\!\big)\varGamma_1(N)].\]
    Choose $m,n$ with $mp-nN = 1$, and observe that $\big(\!\begin{smallmatrix}
        p & 0 \\ 0 & 1
    \end{smallmatrix}\!\big) = \big(\!\begin{smallmatrix}
        1 & n \\ N & mp
    \end{smallmatrix}\!\big)^{-1}\big(\!\begin{smallmatrix}
        1 & 0 \\ 0 & p
    \end{smallmatrix}\!\big)\big(\!\begin{smallmatrix}
        p & n \\ N & m
    \end{smallmatrix}\!\big)$, with $\big(\!\begin{smallmatrix}
        1 & n \\ N & mp
    \end{smallmatrix}\!\big)^{-1}\in\varGamma_1(N)$ and $\big(\!\begin{smallmatrix}
        p & n \\ N & m
    \end{smallmatrix}\!\big)\in\varGamma_0(N)$. Thus with $\varGamma_1(N)$ normal in $\varGamma_0(N)$, we have 
    \[\varGamma_1(N)\big(\!\begin{smallmatrix}
        p & 0 \\ 0 & 1
    \end{smallmatrix}\!\big)\varGamma_1(N) = \varGamma_1(N)\big(\!\begin{smallmatrix}
        1 & 0 \\ 0 & p
    \end{smallmatrix}\!\big)\varGamma_1(N)\big(\!\begin{smallmatrix}
        p & n \\ N & m
    \end{smallmatrix}\!\big).\] If $\varGamma_1(N)\big(\!\begin{smallmatrix}
        1 & 0 \\ 0 & p
    \end{smallmatrix}\!\big)\varGamma_1(N) = \bigsqcup_j\varGamma_1(N)\beta_j$, then $\varGamma_1(N)\big(\!\begin{smallmatrix}
        p & 0 \\ 0 & 1
    \end{smallmatrix}\!\big)\varGamma_1(N) = \bigsqcup_j\varGamma_1(N)\beta_j\big(\!\begin{smallmatrix}
        p & n \\ N & m
    \end{smallmatrix}\!\big)$. Since $m\equiv p^{-1}\smod N$, we have $T_p^\ast = \abr{p}^{-1}T_p$ as desired.
\end{proof}
\begin{corollary}
    By the spectral theorem, the space $\mathcal S_k(\varGamma_1(N))$ has an orthogonal basis of simultaneous eigenvectors, called \textit{eigenfunctions} (or \textit{eigenforms}), of the Hecke operators $\cbr{\abr{n},T_n : \gcd(n,N) = 1}$.
\end{corollary}

Recall that $\omega_N = \big(\!\begin{smallmatrix}
    0 & -1 \\ N & 0
\end{smallmatrix}\!\big)$.
The next result will have several consequences in \cref{section primitive forms}, where we discuss simultaneous eigenfunctions of the Hecke operators.

\begin{theorem}\label{thm: miyake thm 4.5.5}
    The following diagram is commutative: 
    % https://q.uiver.app/#q=WzAsNCxbMCwwLCJcXG1hdGhjYWwgTV9rKE4sXFxjaGkpIl0sWzAsMSwiXFxtYXRoY2FsIE1fayhOLFxcb3ZlcmxpbmV7XFxjaGl9KSJdLFs0LDAsIlxcbWF0aGNhbCBNX2soTixcXGNoaSkiXSxbNCwxLCJcXG1hdGhjYWwgTV9rKE4sXFxvdmVybGluZXtcXGNoaX0pIl0sWzAsMSwiXFxvbWVnYV9OIl0sWzEsMywiVF5cXGFzdChuKVxcdGV4dHsgXFx0ZXh0dXB7KH1yZXNwLiAkVF5cXGFzdChtLGwpJFxcdGV4dHVweyl9fSJdLFswLDIsIlQobilcXHRleHR7IFxcdGV4dHVweyh9cmVzcC4gJFQobCxtKSRcXHRleHR1cHspfX0iXSxbMiwzLCJcXG9tZWdhX04iXV0=
\[\begin{tikzcd}
	{\mathcal M_k(N,\chi)} && {\mathcal M_k(N,\chi)} \\\\
	{\mathcal M_k(N,\overline{\chi})} && {\mathcal M_k(N,\overline{\chi}).}
	\arrow["{|_k[\omega_N]}", from=1-1, to=3-1]
	\arrow["{T^\ast_n}", from=3-1, to=3-3]
	\arrow["{T_n}", from=1-1, to=1-3]
	\arrow["{|_k[\omega_N]}", from=1-3, to=3-3]
\end{tikzcd}\]
\end{theorem}
\begin{proof}
    It suffices to show the result for $\abr{p}$ and $T_p$ in place of $T_n$.

    Let $f\in\mathcal M_k(N,\overline{\chi})$. Then $\abr{p}^\ast f = \abr{p}^{-1}f = \chi(p)f$, and for any $\gamma\in\varGamma_0(N)$ with $\gamma \equiv \big(\!\begin{smallmatrix}
        \ast & \ast \\ 0 & p
    \end{smallmatrix}\!\big)\smod N$, we have $\omega_N^{-1}\gamma\omega_N\equiv \big(\!\begin{smallmatrix}
        \ast & \ast \\ 0 & p^{-1}
    \end{smallmatrix}\!\big)\smod N$, where $p^{-1}$ is the inverse of $p$ modulo $N$. Therefore $f|_k[\omega_N^{-1}\gamma\omega_N] = \chi(p)f$ as desired.

    Let $\varGamma_1(N)\big(\!\begin{smallmatrix}
        1 & 0 \\ 0 & p
    \end{smallmatrix}\!\big)\varGamma_1(N) = \bigsqcup_v \varGamma_1(N)\alpha_v$. A computation reveals that $\omega_N^{-1}\varGamma_1(N)\omega_N = \varGamma_1(N)$, from which we obtain
    \[\varGamma_1(N)\big(\!\begin{smallmatrix}
        p & 0 \\ 0 & 1
    \end{smallmatrix}\!\big)\varGamma_1(N) = \omega_N^{-1}\varGamma_1(N)\big(\!\begin{smallmatrix}
        1 & 0 \\ 0 & p
    \end{smallmatrix}\!\big)\varGamma_1(N)\omega_N = \bigsqcup_v \varGamma_1(N)(\omega_N^{-1}\alpha_v\omega_N).\] Thus 
    \begin{align*}
        ((f|_k[\omega_N^{-1}])|_k[\varGamma_1(N)\big(\!\begin{smallmatrix}
            1 & 0 \\ 0 & p
        \end{smallmatrix}\!\big)\varGamma_1(N)])|_k[\omega_N] &= p^{k/2-1}\sum_vf|_k[\omega_N^{-1}\alpha_v\omega_N]\\
        &= f|_k[\varGamma_1(N)\big(\!\begin{smallmatrix}
            p & 0 \\ 0 & 1
        \end{smallmatrix}\!\big)\varGamma_1(N)].
    \end{align*}
    The result follows from the fact that $T_p^\ast$ is given by $|_k[\varGamma_1(N)\big(\!\begin{smallmatrix}
        1 & 0 \\ 0 & p
    \end{smallmatrix}\!\big)\varGamma_1(N)]^\ast = |_k[\varGamma_1(N)\big(\!\begin{smallmatrix}
        p & 0 \\ 0 & 1
    \end{smallmatrix}\!\big)\varGamma_1(N)]$.
\end{proof}

We include the following result without proof.
\begin{theorem}\label{thm: miyake thm 4.5.10}
    Let $M$ be a multiple of $N$, and $n$ a positive integer coprime to $M$. Then the following diagram is commutative:
    % https://q.uiver.app/#q=WzAsNCxbMCwwLCJcXG1hdGhjYWwgTV9rKE4sXFxjaGkpIl0sWzMsMCwiXFxtYXRoY2FsIE1fayhOLFxcY2hpKSJdLFswLDIsIlxcbWF0aGNhbCBNX2soTSxcXGNoaSkiXSxbMywyLCJcXG1hdGhjYWwgTV9rKE0sXFxjaGkpIl0sWzAsMl0sWzEsM10sWzAsMSwiVF9uIFxcdGV4dHsgKHJlc3AuIH0gVF5cXGFzdF9uKSJdLFsyLDMsIlRfbiBcXHRleHR7IChyZXNwLiB9IFReXFxhc3RfbikiXV0=
    \[\begin{tikzcd}[ampersand replacement=\&]
	{\mathcal M_k(N,\chi)} \&\&\& {\mathcal M_k(N,\chi)} \\
	\\
	{\mathcal M_k(M,\chi)} \&\&\& {\mathcal M_k(M,\chi)}
	\arrow[from=1-1, to=3-1]
	\arrow[from=1-4, to=3-4]
	\arrow["{T_n \text{ \textup{(}resp. } T^\ast_n)}", from=1-1, to=1-4]
	\arrow["{T_n \text{ \textup{(}resp. } T^\ast_n)}", from=3-1, to=3-4]
    \end{tikzcd}\]
    Here the vertical arrows indicate the natural embeddings. The Hecke operators are taken at the appropriate levels.
\end{theorem}
See \cite[Theorem 4.5.10]{miyake} for details.

Similarly to $\mathcal R(N)$, we collect the adjoints of the Hecke operators and form the following algebra:
\begin{definition}
    Denote by \boldindex{$\mathcal R^\ast(N)$} the polynomial ring 
    \[\mathcal R^\ast(N) = \mathbb{Z}[T_p^\ast, \abr{q}^\ast : p,q\in\mathbb P, q\nmid N].\]
    The ring $\mathcal R^\ast(N)$ is the ring of all adjoints of Hecke operators in the Hecke algebra $\mathcal R(N)$. 
\end{definition}

We would like to eliminate the requirement that $\gcd(n,N)= 1$ and find a space of cusp forms for which all of the Hecke operators are simultaneously diagonalizable.

\newpage\section{Oldforms, newforms, and strong multiplicity one}
We consider an elementary way of moving between spaces of modular forms at different levels, specifically to take forms of lower levels $M$ dividing $N$ to forms of level $N$. Restricting our view to cusp forms, we define the so-called spaces of oldforms and newforms, which distinguish between the cusp forms that may be obtained by lifting cusp forms of lower levels up to level $N$, and those which are not obtained in this way. We show that the space of newforms has a basis consisting of normalized simultaneous eigenfunctions of all of the Hecke operators, and conclude with the strong multiplicity one property. The strong multiplicity one property shows that a normalized cusp form that is a common eigenfunction of all of the Hecke operators is completely determined by almost all of its eigenvalues.
\subsection{Oldforms and newforms for $\chi$-eigenspaces}
For the remainder of this work, we restrict our view to various $\chi$-eigenspaces of $\mathcal M_k(\varGamma_1(N))$ and $\mathcal S_k(\varGamma_1(N))$. Specifically, we fix $N$ and consider $\chi$-eigenspaces at different levels $M$, where $M$ divides $N$.

\begin{definition}
    In what follows, denote by \boldindex{$\alpha_\ell$} the matrix
    \[\begin{pmatrix}
        \ell & 0 \\ 0 & 1
    \end{pmatrix}\] for a positive integer $\ell$. The action of matrices of this form on modular forms, which is essentially to scale the argument by $\ell$, appears many times in what follows.
\end{definition}

\begin{lemma}\label{lem: miyake lem 4.6.1}
    For $f\in\mathcal M_k(N,\chi)$, the function $f\circ\alpha_\ell$ is an element of $\mathcal M_k(\ell N,\chi)$ with
    \[(f\circ\alpha_\ell)(z) = f(\ell z) = \ell^{-k/2}(f|_k[\alpha_\ell])(z).\] If $f$ is a cusp form, then so is $f\circ\alpha_\ell$.
\end{lemma}
\begin{proof}
    Indeed, for any $\gamma = \big(\!\begin{smallmatrix}
        a & b \\ c\ell N & d
    \end{smallmatrix}\!\big)\in\varGamma_0(\ell N)$, 
    \[(f|_k[\alpha_\ell])|_k[\gamma] = (f|_k[\alpha_\ell\gamma\alpha_\ell^{-1}])|_k[\alpha_\ell] = (f|_k[\big(\!\begin{smallmatrix}
        a & b\ell \\ cN & d
    \end{smallmatrix}\!\big)])|_k[\alpha_\ell] = \chi(d)f|_k[\alpha_\ell] = \chi(\gamma)f|_k[\alpha_\ell],\] which implies that $f|_k[\alpha_\ell]\in\mathcal M_k(\ell N,\chi)$. The last statement is clear.
\end{proof}
In this setting, the map $\alpha_\ell$ provides a way to move between $\chi$-eigenspaces at different levels.

We saw at the end of the previous section that the Hecke operators $T_n$, for $n$ coprime to $N$, are simultaneously diagonalizable on $\mathcal S_k(\varGamma_1(N))$ and hence also on $\mathcal S_k(N,\chi)$. Unfortunately, $\mathcal S_k(N,\chi)$ does not necessarily have a basis of eigenfunctions of \textit{all} Hecke operators $T_n$. However, $\mathcal S_k(N,\chi)$ does have such a basis when $\chi$ is a primitive character of conductor $N$. We might expect that if we exclude cusp forms of lower levels in $\mathcal S_k(N,\chi)$ (that is, the cusp forms appearing as the image of some suitable $\alpha_\ell$ in the manner described above), we may be able to simultaneously diagonalize all of the Hecke operators $T_n$. This is indeed the case, as we will see.

\begin{lemma}\label{lem: miyake lemma 4.6.2}
    The Hecke operators $T_n,\abr{n}$ commute with the action of $\alpha_\ell$ whenever $n$ is coprime to $\ell N$; that is, the diagram 
    \[\begin{tikzcd}
	{\mathcal M_k(N,\chi)} &&& {\mathcal M_k(N,\chi)} \\
	\\
	{\mathcal M_k(\ell N,\chi)} &&& {\mathcal M_k(\ell N,\chi)}
	\arrow["{T_n ~(\text{resp. }\abr{n})}", from=1-1, to=1-4]
	\arrow["{|_k[\alpha_\ell]}", from=1-4, to=3-4]
	\arrow["{|_k[\alpha_\ell]}"', from=1-1, to=3-1]
	\arrow["{T_n ~(\text{resp. }\abr{n})}"', from=3-1, to=3-4]
    \end{tikzcd}\] commutes.
\end{lemma}
\begin{proof}
    Let $f$ be an element of $\mathcal M_k(N,\chi)$. Since $\gcd(n,\ell N) = 1$, we have $(\abr{n}f)|_k[\alpha_\ell] = \chi(n)f|_k[\alpha_\ell] = \abr{n}(f|_k[\alpha_\ell])$.

    Therefore it suffices to show that $T_p$ commutes with the action of $\alpha_\ell$ for primes $p$ not dividing $\ell N$. From \cref{prop: action of Tp}, 
    \begin{align*}
        T_p(f|_k[\alpha_\ell]) &= p^{k/2-1}\sum_{j=0}^{p-1}f|_k[\alpha_\ell\big(\!\begin{smallmatrix}
            1 & j \\ 0 & p
        \end{smallmatrix}\!\big)] + p^{k/2-1}\chi(p)f|_k[\alpha_\ell\big(\!\begin{smallmatrix}
            p & 0 \\ 0 & 1
        \end{smallmatrix}\!\big)]\\
        &= p^{k/2-1}\sum_{j=0}^{p-1}f|_k[\big(\!\begin{smallmatrix}
            1 & \ell j \\ 0 & p
        \end{smallmatrix}\!\big)\alpha_\ell] + p^{k/2-1}\chi(p)f|_k[\big(\!\begin{smallmatrix}
            p & 0 \\ 0 & 1
        \end{smallmatrix}\!\big)\alpha_\ell]\\
        &= (T_pf)|_k[\alpha_\ell],
    \end{align*}
    where in the last equality we used the fact that $\ell$ and $p$ are coprime to see that 
    \[\cbr{\begin{pmatrix}
        1 & \ell j \\ 0 & p
    \end{pmatrix},\begin{pmatrix}
        m & n \\ N & p
    \end{pmatrix}\begin{pmatrix}
        p & 0 \\ 0 & 1
    \end{pmatrix} : 0\leq j < p, mp-nN = 1}\] form a set of coset representatives for $\varGamma_1(N)\backslash \varGamma_1(N)\big(\!\begin{smallmatrix}
        1 & 0 \\ 0 & p
    \end{smallmatrix}\!\big)\varGamma_1(N)$. The result follows.
\end{proof}

We state the next lemma without proof.
\begin{lemma}\label{lem: miyake lemma 4.6.3}
    Let $f$ be an element of $\mathcal M_k(N,\chi)$. If there exists an element $\alpha = \smallabcd\in\Mat_2(\mathbb Z)$ such that $c\equiv 0\smod N$, $\gcd(a,N)=1$, and $\det\alpha >0$ that satisfies the following conditions, then $f = 0$.
    \begin{enumerate}[label=\textup{(\arabic*)}]
        \item $\det\alpha >1$, $\gcd(\det\alpha,N) = 1$, and $\gcd(a,b,c,d) = 1$, and 
        \item $f|_k[\alpha]\in\mathcal M_k(N,\chi)$.
    \end{enumerate} 
\end{lemma}
This is a technical result that will be used in a few results to come. See \cite[Lemma 4.6.3]{miyake} for details.

The next result provides a way to determine if a cusp form may be obtained as a cusp form at a lower level.
\begin{theorem}\label{thm: miyake thm 4.6.4}
    Let $\ell$ be a positive integer, and let $f$ be a holomorphic function on $\mathcal H$ such that $f(z+1) = f(z)$ and $f\circ \alpha_\ell\in\mathcal M_k(N,\chi)$. Denote by $m_\chi$ the conductor of $\chi$. Then \begin{enumerate}[label=\textup{(\arabic*)}]
        \item if $\ell m_\chi\mid N$, then $f\in\mathcal M_k(N/\ell,\chi)$, and 
        \item if $\ell m_\chi\nmid N$, then $f = 0$.
    \end{enumerate}
    Here $m_\chi$ is the conductor of $\chi$. If $f\circ \alpha_\ell$ is a cusp form, then so is $f$.
\end{theorem}
\begin{proof}
    We may assume that $\ell$ is a prime number. We first show that $f\in\mathcal M_k(N,\chi)$. Let 
    \[\varGamma^\prime = \Big\{\begin{pmatrix}
        a & b \\ cN & d
    \end{pmatrix} \in\varGamma_0(N) : b\equiv 0\smod \ell\Big\},\] and let $\gamma = \big(\!\begin{smallmatrix}
        a & b\ell \\ cN & d
    \end{smallmatrix}\!\big)$ be an element of $\varGamma^\prime$. Since $\alpha_\ell^{-1}\gamma\alpha_\ell = \big(\!\begin{smallmatrix}
        a & b \\ c\ell N & d
    \end{smallmatrix}\!\big)\in\varGamma_0(N)$, we have
    \[f|_k[\gamma\alpha_\ell] = f|_k[\alpha_\ell\alpha_\ell^{-1}\gamma\alpha_\ell] =(f|_k[\alpha_\ell])|_k[\alpha_\ell^{-1}\gamma\alpha_\ell] =  \chi(d)f|_k[\alpha_\ell],\] from which it follows that $f|_k[\gamma] = \chi(d)f$. Let $\varGamma^{\prime\prime}$ be the group generated by $\varGamma^\prime$ and the matrix $\big(\!\begin{smallmatrix}
        1 & 1 \\ 0 & 1
    \end{smallmatrix}\!\big)$. Since $[\varGamma^{\prime\prime} : \varGamma^\prime]\geq \ell$ and 
    \[[\varGamma_0(N): \varGamma^\prime] = \begin{cases}
        \ell & \ell\mid N,\\
        \ell + 1 & \ell \nmid N,
    \end{cases}\] we have that $\varGamma^{\prime\prime} = \varGamma_0(N)$ or $\varGamma_0(N)$ is generated by $\big(\!\begin{smallmatrix}
        1 & 1 \\ 0 & 1
    \end{smallmatrix}\!\big)$ and $\varGamma^\prime$. Since $f(z+1) = f(z)$, we have that $f\in\mathcal M_k(N,\chi)$.

    Suppose that $\ell$ is coprime to $N$. By taking $\alpha = \alpha_\ell$ in \cref{lem: miyake lemma 4.6.3}, it follows that $f = 0$. So assume that $\ell\mid N$. For any element $\gamma_1 = \big(\!\begin{smallmatrix}
        a & b \\ cN & d
    \end{smallmatrix}\!\big)\in\varGamma_0(N)$, we have 
    \[f|_k[\big(\!\begin{smallmatrix}
        a & b\ell \\ cN/\ell & d
    \end{smallmatrix}\!\big)] = f|_k[\alpha_\ell\gamma_1\alpha_\ell^{-1}] = \chi(d)f,\] and in particular we have $f|_k[\big(\!\begin{smallmatrix}
        1 & 0 \\ N/\ell & 1
    \end{smallmatrix}\!\big)] = f$. Therefore, with $\gamma = \big(\!\begin{smallmatrix}
        1 & m \\ 0 & 1
    \end{smallmatrix}\!\big)\big(\!\begin{smallmatrix}
        1 & 0 \\ N/\ell & 1
    \end{smallmatrix}\!\big)\big(\!\begin{smallmatrix}
        1 & n \\ 0 & 1
    \end{smallmatrix}\!\big)$ for any integers $n,m$, we have $f|_k[\gamma] = f$. Now choose $n$ to be any integer such that $nN/\ell + 1\not\equiv 0\smod \ell$, from which we may choose $m$ to be an integer such that $n(1 + mN/\ell) + m = n + m(nN/\ell + 1) \equiv 0\smod \ell$. Then
    \[\alpha_\ell^{-1}\gamma\alpha_\ell = \begin{pmatrix}
        1 + mN/\ell & [n(1+mN/\ell) + 1]/\ell\\
        N & nN/\ell + 1
    \end{pmatrix}\in\varGamma_0(N),\] and so $f|_k[\gamma] = \chi(1+nN/\ell)f$. 

    Therefore, if $f \neq 0$, then $\chi(1+nN/\ell) = 1$ for any integer $n$ satisfying $\gcd(nN/\ell+1, \ell) = 1$. This implies that $\chi$ is defined modulo $N/\ell$, and that $N$ is divisible by $\ell m_\chi$. This proves (2) in the statement. Now assume that $\ell m_\chi\mid N$. Since $\varGamma_0(N/\ell)$ is generated by $\big(\!\begin{smallmatrix}
        1 & 0 \\ N/\ell & 1
    \end{smallmatrix}\!\big)$ and $\varGamma_0(N)$, we have that $f\in\mathcal M_k(N/\ell,\chi)$ as needed. Lastly, it is clear that if $f\circ\alpha_\ell$ is a cusp form, then so is $f$.
\end{proof}

\begin{lemma}\label{lem: miyake lemma 4.6.5}
    Let $f\in \mathcal M_k(N,\chi)$ have Fourier series expansion $f(z) = \sum_{n=0}^\infty a_ne^{2\pi i n z}$. For a positive integer $L$, let $g$ be given by \[g(z) = \sum_{\substack{n\geq 0\\\gcd(n,L)=1}}a_ne^{2\pi i n z}.\]
    Then $g\in\mathcal M_k(M,\chi)$, with \[M = N\prod_{\substack{p\mid L \\ p\mid N}}p\prod_{\substack{p\mid L\\ p\nmid N}}p^2,\] where the products are to be taken over primes $p$. If $f$ is a cusp form, then so is $g$.
\end{lemma}
\begin{proof}
    It suffices to show the result when $L$ is a prime number $p$. Let $N^\prime = N$ if $p\mid N$ and $N^\prime = pN$ otherwise. Since $p\mid N^\prime$, we have by \cref{prop: action of Tp} that
    \[(T_pf)(z) = \frac{1}{p}\sum_{n=0}^\infty a_n\sum_{m=0}^{p-1}e^{2\pi i n (z+m)/p} = \sum_{n=0}^\infty a_{np}e^{2\pi i n z},\] where $T_p$ is acting at level $N^\prime$. Hence $(T_pf)(pz) = \sum_{n=0}^\infty a_{np}e^{2\pi i n pz}$, which belongs to $\mathcal M_k(N^\prime p,\chi)$ by \cref{lem: miyake lem 4.6.1}. Then the function $g$, given by
    \[g(z) = f(z) - (T_pf)(pz),\] is an element of $\mathcal M_k(N^\prime p, \chi)$. The last assertion is clear.
\end{proof}

\begin{lemma}\label{lem: miyake lemma 4.6.6}
    Let $\chi$ be a Dirichlet character modulo $N$, $\ell$ a positive integer, and $p$ a prime number that does not divide $\ell$. Let $M = \ell N$. Then the following diagrams commute. 
    % https://q.uiver.app/#q=WzAsNCxbMCwwLCJcXG1hdGhjYWwgTV9rKHBOLFxcY2hpKSJdLFswLDIsIlxcbWF0aGNhbCBNX2socE0sXFxjaGkpIl0sWzMsMiwiXFxtYXRoY2FsIE1fayhNLFxcY2hpKSJdLFszLDAsIlxcbWF0aGNhbCBNX2soTixcXGNoaSkiXSxbMCwzLCJcXHZhckdhbW1hXzAocE4pXFxiaWcoXFwhXFxiZWdpbntzbWFsbG1hdHJpeH0xICYgMCBcXFxcIDAgJiBwXFxlbmR7c21hbGxtYXRyaXh9XFwhXFxiaWcpXFx2YXJHYW1tYV8wKE4pIl0sWzAsMV0sWzMsMl0sWzEsMiwiXFx2YXJHYW1tYV8wKHBNKVxcYmlnKFxcIVxcYmVnaW57c21hbGxtYXRyaXh9MSAmIDAgXFxcXCAwICYgcFxcZW5ke3NtYWxsbWF0cml4fVxcIVxcYmlnKVxcdmFyR2FtbWFfMChNKSJdXQ==
    \begin{equation}\label{diag: diagram 1 4.6.6}
        \begin{tikzcd}[ampersand replacement=\&]
            {\mathcal M_k(pN,\chi)} \&\&\& {\mathcal M_k(N,\chi)} \\
            \\
            {\mathcal M_k(pM,\chi)} \&\&\& {\mathcal M_k(M,\chi)}
            \arrow["{\varGamma_1(pN)\big(\!\begin{smallmatrix}1 & 0 \\ 0 & p\end{smallmatrix}\!\big)\varGamma_1(N)}", from=1-1, to=1-4]
            \arrow[from=1-1, to=3-1]
            \arrow[from=1-4, to=3-4]
            \arrow["{\varGamma_1(pM)\big(\!\begin{smallmatrix}1 & 0 \\ 0 & p\end{smallmatrix}\!\big)\varGamma_1(M)}", from=3-1, to=3-4]
            \end{tikzcd}
    \end{equation} \textup{(}In the above diagram, the vertical arrows indicate the natural embeddings.\textup{)}
    % https://q.uiver.app/#q=WzAsNCxbMCwwLCJcXG1hdGhjYWwgTV9rKHBOLFxcY2hpKSJdLFswLDIsIlxcbWF0aGNhbCBNX2socE0sXFxjaGkpIl0sWzMsMiwiXFxtYXRoY2FsIE1fayhNLFxcY2hpKSJdLFszLDAsIlxcbWF0aGNhbCBNX2soTixcXGNoaSkiXSxbMCwzLCJcXHZhckdhbW1hXzAocE4pXFxiaWcoXFwhXFxiZWdpbntzbWFsbG1hdHJpeH0xICYgMCBcXFxcIDAgJiBwXFxlbmR7c21hbGxtYXRyaXh9XFwhXFxiaWcpXFx2YXJHYW1tYV8wKE4pIl0sWzAsMSwiXFxhbHBoYV9cXGVsbCJdLFszLDIsIlxcYWxwaGFfXFxlbGwiXSxbMSwyLCJcXHZhckdhbW1hXzAocE0pXFxiaWcoXFwhXFxiZWdpbntzbWFsbG1hdHJpeH0xICYgMCBcXFxcIDAgJiBwXFxlbmR7c21hbGxtYXRyaXh9XFwhXFxiaWcpXFx2YXJHYW1tYV8wKE0pIl1d
    \begin{equation}\label{diag: diagram 2 4.6.6}
        \begin{tikzcd}[ampersand replacement=\&]
            {\mathcal M_k(pN,\chi)} \&\&\& {\mathcal M_k(N,\chi)} \\
            \\
            {\mathcal M_k(pM,\chi)} \&\&\& {\mathcal M_k(M,\chi)}
            \arrow["{\varGamma_1(pN)\big(\!\begin{smallmatrix}1 & 0 \\ 0 & p\end{smallmatrix}\!\big)\varGamma_1(N)}", from=1-1, to=1-4]
            \arrow["{\alpha_\ell}", from=1-1, to=3-1]
            \arrow["{\alpha_\ell}", from=1-4, to=3-4]
            \arrow["{\varGamma_1(pM)\big(\!\begin{smallmatrix}1 & 0 \\ 0 & p\end{smallmatrix}\!\big)\varGamma_1(M)}", from=3-1, to=3-4]
        \end{tikzcd}
    \end{equation}
    Similar results hold for cusp forms.
\end{lemma}
\begin{proof}
    By assumption, $p\mid N$ if and only if $p \mid M$. Therefore, by \cref{lem: miyake lem 4.5.11}, the commutativity of diagram \labelcref{diag: diagram 1 4.6.6} is clear.

    Now let $f$ be an element of $\mathcal M_k(pN,\chi)$, and let $g = f|_k[\alpha_\ell]$. Let \[d = \begin{cases}
        p-1 & \text{if $p\mid M$}\\
        p & \text{otherwise},
    \end{cases}\]
    and $\gamma_v = \big(\!\begin{smallmatrix}
        1 & v \\ 0 & 1
    \end{smallmatrix}\!\big)$ for $0\leq v < p$. If $d = p$, choose $\gamma_p$ as in \cref{lem: miyake lem 4.5.11}. Then we have 
    \[\varGamma_1(pM)\big(\!\begin{smallmatrix}
        1 & 0 \\ 0 & p
    \end{smallmatrix}\!\big)\varGamma_1(M) = \bigsqcup_{v=0}^d\varGamma_1(pM)\big(\!\begin{smallmatrix}
        1 & 0 \\ 0 & p
    \end{smallmatrix}\!\big)\gamma_v\] and
    \begin{align*}
        p^{1-k/2}g|_k[\varGamma_1(pM)\big(\!\begin{smallmatrix}
            1 & 0 \\ 0 & p
        \end{smallmatrix}\!\big)\varGamma_1(M)] &= \sum_{v = 0}^d g|_k[\big(\!\begin{smallmatrix}
            1 & 0 \\ 0 & p
        \end{smallmatrix}\!\big)\gamma_v]\\
        &= \sum_{v = 0}^d f|_k[\alpha_\ell\big(\!\begin{smallmatrix}
            1 & 0 \\ 0 & p
        \end{smallmatrix}\!\big)\gamma_v]\\
        &= \sum_{v=0}^d f|_k[\big(\!\begin{smallmatrix}
            1 & 0 \\ 0 & p
        \end{smallmatrix}\!\big)(\alpha_\ell\gamma_v\alpha_\ell^{-1})\alpha_\ell].
    \end{align*}
    Observe that $\alpha_\ell\gamma_v\alpha_\ell^{-1} = \big(\!\begin{smallmatrix}
        1 & \ell v \\ 0 & 1
    \end{smallmatrix}\!\big)$ for $0\leq v < p$, and that for $p\nmid N$, we have \[\alpha_\ell\gamma_p\alpha_\ell^{-1} \equiv \begin{cases}
        \big(\!\begin{smallmatrix}
            0 & -a\ell \\ (a\ell)^{-1} & 0
        \end{smallmatrix}\!\big)\smod p, \\
        \big(\!\begin{smallmatrix}
            1 & 0 \\ 0 & 1
        \end{smallmatrix}\!\big)\smod N.
    \end{cases}\] It follows from \cref{lem: miyake lem 4.5.11} that 
    \[\varGamma_1(pN)\big(\!\begin{smallmatrix}
        1 & 0 \\ 0 & p
    \end{smallmatrix}\!\big)\varGamma_1(N) = \bigsqcup_{v=0}^d\varGamma_1(pN)\big(\!\begin{smallmatrix}
        1 & 0 \\ 0 & p
    \end{smallmatrix}\!\big)(\alpha_\ell\gamma_v\alpha_\ell^{-1}).\] Thus 
    \[g|_k[\varGamma_1(pM)\big(\!\begin{smallmatrix}
        1 & 0 \\ 0 & p
    \end{smallmatrix}\!\big)\varGamma_1(N)] = (f|_k[\varGamma_1(pN)\big(\!\begin{smallmatrix}
        1 & 0 \\ 0 & p
    \end{smallmatrix}\!\big)\varGamma_1(N)])|_k[\alpha_\ell].\qedhere\]
\end{proof}

\begin{lemma}\label{lem: miyake lemma 4.6.7}
    Let $\ell$ be a positive square-free integer, and let $f\in\mathcal M_k(N,\chi)$ have Fourier series expansion $f(z) = \sum_{n=0}^\infty a_ne^{2\pi i n z}$. If $a_n =0$ for all $n$ coprime to $\ell$, then 
    \[f(z) = \sum_{p\mid \ell}g_p(pz)\quad \text{where $g_p\in\mathcal M_k(N\ell^2,\chi)$},\] and $p$ runs over all prime factors of $\ell$. Furthermore, if $\ell\mid N$, then we may choose $g_p$ from $\mathcal M_k(N\ell,\chi)$. If $f$ is a cusp form, then each $g_p$ can be chosen to be cusp forms as well.
\end{lemma}
\begin{proof}
    The proof is by induction on the number of prime factors of $\ell$. Suppose first that $\ell$ is prime. Then $g$ given by $g(z) = f(z/\ell)$ satisfies the conditions of \cref{thm: miyake thm 4.6.4}, so that $g\in\mathcal M_k(N/\ell,\chi)$ or $f = g = 0$, depending on whether $\ell m_\chi$ divides $N$ or not. Thus $g\in \mathcal M_k(N\ell, \chi)$ with $f(z) = g(\ell z)$.
    % \sai{what ^^ ???}
    Now let $\ell$ be a composite number, and assume the result holds for any proper divisor of $\ell$. Let $p$ be a prime factor of $\ell$ and let $\ell^\prime = \ell/p$. Let $h$ be given by 
    \[h(z) = \sum_{\substack{n\geq 0 \\ \gcd(n,p) = 1}}a_ne^{2\pi i n z}.\] Then $h\in\mathcal M_k(Np^2,\chi)$ by \cref{lem: miyake lemma 4.6.5}. Let $b_n$ for $n\geq 0$ be the Fourier coefficients of $f-h$; that is, $f(z)-h(z) = \sum_{n = 0}^\infty b_ne^{2\pi i n z}$. Then if $\gcd(n,p) = 1$, $b_n = 0$. Define $g_p$ by 
    \[g_p(z) = f(z/p)- h(z/p),\] and observe that $g_p\in\mathcal M_k(Np,\chi)$ by the inductive hypothesis. Furthermore, note that $h$ satisfies the assumption of the lemma with $Np^2$ and $\ell^\prime$ in place of $N$ and $\ell$, respectively. Therefore we may find $g_q\in\mathcal M_k(Np^2(\ell^\prime)^2,\chi) = \mathcal M_k(N\ell^2,\chi)$ for each prime factor $q$ of $\ell^\prime$ satisfying $h(z) = \sum_{q\mid \ell^\prime}g_q(qz)$. Combining this with the observation that $f(z) = g_p(pz) + h(z)$ yields the first half of the lemma. From \cref{lem: miyake lemma 4.6.5} it is clear that we may take $g_p$ from $\mathcal M_k(N\ell,\chi)$ if $\ell\mid N$. Lastly, it is clear that the $g_p$ may be chosen to be cusp forms if $f$ is a cusp form.
\end{proof}

The following theorem generalizes \cref{thm: miyake thm 4.6.4}.
\begin{thm}\label{thm: miyake thm 4.6.8}
    Let $\ell$ be a positive integer, and let $f\in\mathcal M_k(N,\chi)$ have Fourier series expansion $f(z) = \sum_{n=0}^\infty a_ne^{2\pi i nz}$. Let $m_\chi$ be the conductor of $\chi$, and assume that $a_n = 0$ for all $n$ coprime to $\ell$.
    \begin{enumerate}[label=\textup{(\arabic*)}]
        \item If $\gcd(\ell,N/m_\chi) = 1$, then $f = 0$.
        \item If $\gcd(\ell,N/m_\chi) \neq 1$, then there exist $f_p\in\mathcal M_k(N/p,\chi)$ for all prime factors $p$ of $\gcd(\ell,N/m_\chi)$ such that 
        \[f(z) = \sum_{p\mid \gcd(\ell,N/m_\chi)}f_p(pz).\] If $f$ is a cusp form, then we may choose the $f_p$ to be cusp forms as well.
    \end{enumerate}
\end{thm}
\begin{proof}
    We may assume that $\ell$ is square-free, and we prove the result by induction on the number of factors of $\ell$. When $\ell$ is a prime number, the result is obtained from \cref{thm: miyake thm 4.6.4} for $f(z/\ell)$. Now assume that $\ell$ is composite and that the lemma holds for any proper divisor of $\ell$. Let $p$ be a prime factor of $\ell$, and let $\ell^\prime = \ell/p$.

    Let $h$ and $g$ be given by 
    \[h(z) = \sum_{\substack{n\geq 0\\\gcd(n,\ell^\prime) \neq 1}}a_ne^{2\pi i n z}\quad\text{and}\quad g(z) = f(z) - h(z) = \sum_{\substack{n\geq 0\\\gcd(n,\ell^\prime) = 1}}a_ne^{2\pi i n z}.\]
    Observe that $h$ and $g$ are elements of $\mathcal M_k(N(\ell^\prime)^2,\chi)$ by \cref{thm: miyake thm 4.6.4}. Let $d = p-1$ if $p^2\mid N$, or let $d = p$ if $p^2\nmid N$. Then choose elements $\gamma_v\in \varGamma_1(N(\ell^\prime)^2/p)$ for $0\leq v < d$ (via \cref{lem: miyake lem 4.5.11}) so that 
    \[\varGamma_1(N(\ell^\prime)^2)\big(\!\begin{smallmatrix}
        1 & 0 \\ 0 & p
    \end{smallmatrix}\!\big)\varGamma_0(N(\ell^\prime)^2/p) = \bigsqcup_{v=0}^d\varGamma_1(N(\ell^\prime)^2)\big(\!\begin{smallmatrix}
        1 & 0 \\ 0 & p
    \end{smallmatrix}\!\big)\gamma_v.\] 
    Then \begin{align*}
        g|_k[\varGamma_1(N(\ell^\prime)^2)\big(\!\begin{smallmatrix}
            1 & 0 \\ 0 & p
        \end{smallmatrix}\!\big)\varGamma_1(N(\ell^\prime)^2/p)] &= p^{k/2-1}\sum_{v=0}^d \overline{\chi(\gamma_v)}g|_k[\big(\!\begin{smallmatrix}
            1 & 0 \\ 0 & p
        \end{smallmatrix}\!\big)\gamma_v]\\
        &= p^{-1}\sum_{v=0}^d\overline{\chi(\gamma_v)}g_p|_k[\gamma_v]\\
        &= (d+1)p^{-1}g_p.
    \end{align*}
    Hence
    \[g(z) = g_p(pz) = p(d+1)^{-1}\big(g|_k[\varGamma_1(N(\ell^\prime)^2)\big(\!\begin{smallmatrix}
        1 & 0 \\ 0 & p
    \end{smallmatrix}\!\big)\varGamma_1(N(\ell^\prime)^2/p)]\big)(pz).\] Let 
    \[f_p(z) = p(d+1)^{-1}\big(f|_k[\varGamma_1(N)\big(\!\begin{smallmatrix}
        1 & 0 \\ 0 & p
    \end{smallmatrix}\!\big)\varGamma_1(N/p)]\big)(z).\]
    Then $f_p\in\mathcal M_k(N/p,\chi)$, and by \cref{lem: miyake lemma 4.6.6}(1),
    \begin{equation}\label{eqn: 4.6.8 equation}
        f_p(z) = p(d+1)^{-1}\big(f|_k[\varGamma_1(N(\ell^\prime)^2)\big(\!\begin{smallmatrix}
            1 & 0 \\ 0 & p
        \end{smallmatrix}\!\big)\varGamma_1(N(\ell^\prime)^2/p)]\big)(z).
    \end{equation}
    We show that $f(z)-f_p(pz)$ satisfies the assumption of the theorem for $\ell^\prime$. It is clear that $f(z)-f_p(pz)$ defines an element of $\mathcal M_k(N,\chi)$. We have 
    \begin{align}\label{eqn: 4.6.8 equation2}
        f(z) - f_p(pz) &= f(z) - f_p(pz) - g(z) + g_p(pz)\\
        &= (f(z) - g(z)) - p(d+1)^{-1}\big((f-g)|_k[\varGamma_1(N(\ell^\prime)^2)\big(\!\begin{smallmatrix}
            1 & 0 \\ 0 & p
        \end{smallmatrix}\!\big)\varGamma_1(N(\ell^\prime)^2/p)]\big)(pz)\nonumber\\
        &= h(z) - p(d+1)^{-1}\big(h|_k[\varGamma_1(N(\ell^\prime)^2)\big(\!\begin{smallmatrix}
            1 & 0 \\ 0 & p
        \end{smallmatrix}\!\big)\varGamma_1(N(\ell^\prime)^2/p)]\big)(pz).\nonumber
    \end{align}
    Since the $n$-th Fourier coefficient of $h$ vanishes if $n$ is coprime to $\ell^\prime$, we may write 
    \[h(z) = \sum_{q\mid \ell^\prime}h_q(qz)\quad \text{with $h_q\in\mathcal M_k(N(\ell^\prime)^3m,\chi)$}\] by \cref{lem: miyake lemma 4.6.7}. Moreover, for any prime factor $q$ of $\ell^\prime$, 
    \[h|_k[\varGamma_1(N(\ell^\prime)^2)\big(\!\begin{smallmatrix}
        1 & 0 \\ 0 & p
    \end{smallmatrix}\!\big)\varGamma_1(N(\ell^\prime)^2/p)] = h|_k[\varGamma_1(N(\ell^\prime)^3q)\big(\!\begin{smallmatrix}
        1 & 0 \\ 0 & p
    \end{smallmatrix}\!\big)\varGamma_1(N(\ell^\prime)^3q/p)]\] by \cref{lem: miyake lemma 4.6.6}(1). Combined with \cref{lem: miyake lemma 4.6.6}(2), we have that 
    \begin{align*}
        \big(h|_k[\varGamma_1(N(\ell^\prime)^2)\big(\!\begin{smallmatrix}
            1 & 0 \\ 0 & p
        \end{smallmatrix}\!\big)\varGamma_1(N(\ell^\prime)^2/p)]\big)(z) &= \bigg(\bigg(\sum_{q\mid \ell^\prime}q^{-k/2}h_q|_k[\alpha_q]\bigg)\bigg|_k[\varGamma_1(N(\ell^\prime)^3)\big(\!\begin{smallmatrix}
            1 & 0 \\ 0 & p
        \end{smallmatrix}\!\big)\varGamma_1(N(\ell^\prime)^3q/p)]\bigg)(z)\\
        &= \sum_{q\mid \ell^\prime}\big(h_q|_k[\varGamma_1(N(\ell^\prime)^3)\big(\!\begin{smallmatrix}
            1 & 0 \\ 0 & p
        \end{smallmatrix}\!\big)\varGamma_1(N(\ell^\prime)^3/p)]\big)(qz).
    \end{align*}
    In particular, the $n$-th Fourier coefficient of $\big(h|_k[\varGamma_1(N(\ell^\prime)^2)\big(\!\begin{smallmatrix}
        1 & 0 \\ 0 & p
    \end{smallmatrix}\!\big)\varGamma_1(N(\ell^\prime)^2/p)]\big)(z)$ vanishes whenever $\gcd(n,\ell^\prime) = 1$. As a result, the $n$-th Fourier coefficient of $f(z)- f_p(pz)$ vanishes by \eqref{eqn: 4.6.8 equation2}. Therefore, by the inductive hypothesis, we have 
    \[f(z) - f_p(pz) = \sum_qf_q(qz),\quad \text{for $f_q\in\mathcal M_k(N/q,\chi)$},\]
    where $q$ runs over all prime factors of $\ell^\prime$. This proves part (2) of the theorem. If $f$ is a cusp form, then it is clear that each of the modular forms appearing in the proof may be chosen to be cusp forms also. 
\end{proof}

In what follows, we only consider cusp forms.
\begin{definition}\label{def: eigenspace oldforms and newforms}
    Denote by \boldindex{$\mathcal S_k^1(N,\chi)$} the subspace of \textbf{\textit{oldforms at level $N$}} of $\mathcal S_k(N,\chi)$, which is generated by the set 
    \[\bigcup_M\bigcup_\ell\cbr{f\circ\alpha_\ell : f\in\mathcal S_k(M,\chi)}.\] Here $M$ runs over all positive integers such that $m_\chi\mid M$, $M\mid N$, and $M\neq N$; here $\ell$ runs over all positive divisors of $N/M$, including $1$ and $N/M$. (As usual, $m_\chi$ denotes the conductor of $\chi$.)

    In other words, $\mathcal S_k^1(N,\chi)$ is the subspace of $\mathcal S_k(N,\chi)$ generated by cusp forms coming from lower levels.

    Denote by \boldindex{$\mathcal S_k^0(N,\chi)$} the subspace of \textbf{\textit{newforms at level $N$}}, given by the orthogonal complement of $\mathcal S_k^1(N,\chi)$ in $\mathcal S_k(N,\chi)$ with respect to the Petersson inner product.
\end{definition}

In other sources, the ${}^1$ and ${}^0$ denoting the spaces of oldforms and newforms is replaced by ${}^\mathrm{old}$ and ${}^\mathrm{new}$, respectively.

By the above definition, we immediately obtain the following lemma:
\begin{lemma}\label{lem: miyake lemma 4.6.9}
    \begin{enumerate}[label=\textup{(\arabic*)}]
        \item If $\chi$ is a primitive Dirichlet character of conductor $N$, then $\mathcal S_k(N,\chi) = \mathcal S_k^0(N,\chi)$.
        \item Let $\chi$ be a character of conductor $m_\chi$. If $m_\chi\mid M$, $M\mid N$, and $M\neq N$, then $\mathcal S_k(M,\chi)\subset\mathcal S_k^1(N,\chi)$.
        \item The space $\mathcal S_k(N,\chi)$ is generated by the set 
        \[\bigcup_M\bigcup_\ell \cbr{f\circ \alpha_\ell : f\in\mathcal S_k^0(N,\chi)}.\] Here $M$ runs over all positive integers such that $m_\chi\mid M$ and $M\mid N$; here $\ell$ runs over all positive divisors of $N/M$, including $1$ and $N/M$.
    \end{enumerate}
\end{lemma}

\begin{lemma}\label{lem: miyake lemma 4.6.10}
    The subspaces $\mathcal S_k^1(N,\chi)$ and $\mathcal S_k^0(N,\chi)$ are stable under the action of the Hecke operators $T_n$ for $n$ coprime to $N$.
\end{lemma}
\begin{proof}
    Let $f\in\mathcal S_k^1(N,\chi)$. By definition, we may write $f$ as 
    \[f(z) = \sum_vf_v(\ell_v z)\quad \text{for $f_v\in\mathcal S_k(M_v,\chi)$,}\quad \text{with $\ell_v M_v\mid N$, $M_v\neq N$.}\]
    Let $g_v$ be given by $g_v(z) = f_v(\ell_v z)$. Since $\gcd(n,\ell_v N) = 1$, we have 
    \[(T_nf)(z) = \sum_v(T_ng_v)(z) = \sum_v (T_nf_v)(\ell_v z)\] by \cref{thm: miyake thm 4.5.10} and \cref{lem: miyake lemma 4.6.2}. Hence $T_nf\in\mathcal S_k^1(N,\chi)$; that is, $\mathcal S_k^1(N,\chi)$ is stable under $T_n$. Since the adjoint of $T_n$ on $\mathcal S_k(N,\chi)$ is $\abr{n}T_n$ (see \cref{thm: the actual adjoints}), it follows that $\mathcal S_k^0(N,\chi)$ is also stable under $T_n$.
\end{proof}

It follows from the above lemma that the subspaces $\mathcal S_k^0(N,\chi)$ and $\mathcal S_k^1(N,\chi)$ of $\mathcal S_k(N,\chi)$ have a basis of simultaneous eigenfunctions of all of the Hecke operators $T_n$, for $n$ coprime to $N$.

\subsection{Primitive forms}\label{section primitive forms}

In this section we briefly focus on simultaneous eigenfunctions of the Hecke operators in the space of newforms.

\begin{lemma}\label{lem: miyake lemma 4.6.11}
    Let $f\in\mathcal S_k^0(N,\chi)$ have series expansion $\sum_{n=1}^\infty a_ne^{2\pi i n z}$. If $f$ is a common eigenfunction of the Hecke operators $T_n$ for all $n$ coprime to some integer $L$, then $a_1\neq 0$.
\end{lemma}
\begin{proof}
    Assume by way of contradiction that $a_1 = 0$. By \cref{lem: miyake lem 4.5.15}(1), we have that $a_n = 0$ for all $n$ coprime to $L$. Hence $f\in\mathcal S_k^1(N,\chi)$ by \cref{thm: miyake thm 4.6.8}, which is a contradiction.
\end{proof}

\begin{theorem}\label{thm: miyake thm 4.6.12}
    Let $f(z)\in \mathcal S_k^0(N,\chi)$ and $g(z)\in \mathcal S_k(N,\chi)$. If $f(z)$ and $g(z)$ are common eigenfunctions of $T_n$ with the same eigenvalue for each $n$ coprime to some integer $L$, then $g(z)$ is a constant multiple of $f(z)$.
\end{theorem}
\begin{proof}
    Let $f$ have Fourier expansion $\sum_{n=1}^\infty a_ne^{2\pi i n z}$. Since $a_1\neq 0$ by \cref{lem: miyake lemma 4.6.11}, normalize $f$ so that $a_1=1$. Furthermore assume that $N\mid L$. Decompose $g$ into $g^0+g^1$, the sum of its newform and oldform components. By \cref{lem: miyake lemma 4.6.10}, both $g^0$ and $g^1$ are common eigenfunctions of the operators $T_n$ with the same eigenvalue $a_n$ for each $n$ coprime to $L$. 

    Assume that $g^0\neq 0$, and write $g^0(z) = \sum_{n=1}^\infty b_ne^{2\pi i n z}$. By \cref{lem: miyake lemma 4.6.11}, $b_1\neq 0$. We show that $g^0 = b_1f$. Let $\sum_{n=1}^\infty c_ne^{2\pi i n z}$ be the Fourier series expansion for $g^0-b_1f$. Since $b_1a_n = b_n$ for all $n$ coprime to $L$ by \cref{lem: miyake lem 4.5.15}(1), we have that $c_n = 0$ for all $n$ coprime to $L$. Then \cref{thm: miyake thm 4.6.8} implies that $g^0-b_1f\in\mathcal S_k^1(N,\chi)$, and that $g^0 = b_1f$.

    Next we show that $g^1(z) = 0$. First suppose that $N = m_\chi$, the conductor of $\chi$. Then $\mathcal S_k^1(N,\chi) = 0$, so that $g^1 = 0$. If $N\neq m_\chi$.
    
    If $g^1\neq 0$, then there exists a proper divisor $M$ of $N$ satisfying $m_\chi\mid M$, and a nonzero element $h$ of $\mathcal S_k^0(M,\chi)$ such that $T_nh = a_nh$ for all $n$ coprime to $L$. Then by definition we may write 
    \[g^1(z) = \sum_v h_v(\ell_v z),\quad h_v\in\mathcal S_k^0(M_v,\chi), \ell_vM_v\mid N, M_v\neq N.\]
    Since $M_v$ divides $N$, \cref{lem: miyake lemma 4.6.10} implies that $\mathcal S_k^0(M_v,\chi)$ has a basis consisting of eigenfunctions of the operators $T_n$ for all $n$ coprime to $L$. Therefore we may assume that all $h_v$ are common eigenfunctions of $T_n$ for all $n$ coprime to $L$. Since eigenfunctions with distinct eigenvalues are linearly independent, the summation of all $h_v(\ell_v z)$ whose eigenvalues for $T_n$ are different from $a_n$ must vanish. So by removing those functions from the sum, we may assume that the $h_v$ appearing in the sum $g^1(z) = \sum_v h_v(\ell_v z)$ satisfy $T_nh_v = a_nh_v$ for all $n$ coprime to $L$. Therefore we may take any $h_v$ as $h$ and any $M_v$ as $M$.

    Let $h(z)= \sum_{n=1}^\infty c_n^\prime e^{2\pi i nz}$ be the Fourier series of the element $h$ from before. Since $T_nh = a_nh$ for all $n$ coprime to $L$, we have that $c^\prime_1\neq 0$ by \cref{lem: miyake lemma 4.6.11}. Let $\sum_{n=1}^\infty d_ne^{2\pi i n z}$ be the Fourier series expansion for $h-c^\prime_1f$. By \cref{lem: miyake lem 4.5.15}(1), $d_n = 0$ if $n$ is coprime to $L$, and so by \cref{thm: miyake thm 4.6.8}, $h-c^\prime_1f\in\mathcal S_k^1(N,\chi)$.

    Hence $f(z) = -c^\prime_1(h(z)-c^\prime_1f(z)) + c^\prime_1h(z)$, so that $f$ is an element of $\mathcal S_k^1(N,\chi)$, which contradicts the assumption that $f$ was an element of $\mathcal S_k^0(N,\chi)$. Therefore $g^1$ must be zero, so that $g=g^0 = b_1f$.
\end{proof}

\begin{definition}
    An element $f\in\mathcal S_k^0(N,\chi)$ is a \boldindex{primitive form} of conductor $N$ if the following conditions are satisfied:
    \begin{enumerate}[label=\textup{(\arabic*)}]
        \item $f$ is a common eigenfunction of all of the Hecke operators $T_n$ for $n$ coprime to $N$, and
        \item $f$ is normalized; that is, $f$ has Fourier series expansion $f(z) = \sum_{n=1}^\infty a_ne^{2\pi i n z}$ with $a_1 = 1$.
    \end{enumerate}
    We also call $\mathcal S_k^0(N,\chi)$ the space of primitive forms of level $N$ with character $\chi$.
\end{definition}

In other sources, primitive forms are called newforms, which may be confusing.

\begin{theorem}\label{thm: miyake thm 4.6.13}
    Primitive forms are common eigenfunctions of all operators in $\mathcal R(N)\cup\mathcal R^\ast(N)$, and $\mathcal S_k^0(N,\chi)$ has a basis consisting of primitive forms.
\end{theorem}
\begin{proof}
    Let $f$ be a primitive form in $\mathcal S_k^0(N,\chi)$, and let $T_nf = a_nf$ for all $n$ coprime to $N$. Let $T$ and $T^\ast$ be elements of $\mathcal R(N)$ and $\mathcal R^\ast(N)$, respectively. Since $T$ commutes with $T_n$, $T^\ast$ commutes with $T_n^\ast$, and $T_nf = \chi(n)T_n^\ast f$, we have that $T^\ast$ commutes with $T_n$. Therefore $Tf$ and $T^\ast f$ are also common eigenfunctions of the operators $T_n$, with the same eigenvalue $a_n$.

    Then \cref{thm: miyake thm 4.6.12} implies that both $Tf$ and $T^\ast f$ are constant multiples of $f$, as desired.

    Since $\mathcal S_k^0(N,\chi)$ and $\mathcal S_k(N,\chi)$ have bases consisting of common eigenfunctions of the operators $T_n$ for $n$ coprime to $N$, what we just proved along with \cref{lem: miyake lemma 4.6.11} imply that $\mathcal S_k^0(N,\chi)$ has a basis of primitive forms.
\end{proof}

% \sai{below is equation 4.6.18, maybe more elaboration is needed}
Let $f$ be a primitive form of $\mathcal S_k^0(N,\chi)$ with Fourier series expansion $f(z) = \sum_{n=1}^\infty a_ne^{2\pi i n z}$. Then  
\begin{equation}
    T_n(f|_k[\omega_N]) = \overline{a_n}f|_k[\omega_N]\quad \text{and}\quad T_n^\ast(f|_k[\omega_N]) = a_nf|_k[\omega_N]
\end{equation}
for all positive integers $n$.

We omit the proofs of the following three results. 
\begin{theorem}\label{thm: miyake thm 4.6.15}
    \begin{enumerate}[label=\textup{(\arabic*)}]
        \item By the action of $\omega_N$, we obtain the isomorphisms $\mathcal S_k^0(N,\chi)\cong \mathcal S_k^0(N,\overline{\chi})$ and $\mathcal S_k^1(N,\chi)\cong \mathcal S_k^1(N,\overline{\chi})$, and 
        \item if $f(z)$ is a primitive form of $\mathcal S_k^0(N,\chi)$, then $f_p(z)$ is a primitive form of $\mathcal S_k^0(N,\overline{\chi})$ and 
        \[f|_k[\omega_N] = cf_p(z)\] for some $c\in\mathbb C$.
    \end{enumerate}
\end{theorem}
See \cite[Theorem 4.6.15]{miyake} for details.

Let $\gamma_q$ and $\gamma_q^\prime$ be two elements of $\slz$ such that
\[\gamma_q \equiv \begin{cases}
    \big(\!\begin{smallmatrix}
        0 & -1 \\ 1 & 0
    \end{smallmatrix}\!\big)\smod{N_q^2}\\
    \big(\!\begin{smallmatrix}
        1 & 0 \\ 0 & 1
    \end{smallmatrix}\!\big) \smod{(N/N_q)^2}
\end{cases}\quad\text{and}\quad \gamma_q^\prime\equiv \begin{cases}
    \big(\!\begin{smallmatrix}
        1 & 0 \\ 0 & 1
    \end{smallmatrix}\!\big)\smod{N_q^2}\\
    \big(\!\begin{smallmatrix}
        0 & -1 \\ 1 & 0
    \end{smallmatrix}\!\big)\smod{(N/N_q)^2}.
\end{cases}\]
Then let $\eta_q = \gamma_q\big(\!\begin{smallmatrix}
    N_q & 0 \\ 0 & 1
\end{smallmatrix}\!\big)$ and $\eta_q^\prime = \gamma_q^\prime \big(\!\begin{smallmatrix}
    N/N_q & 0 \\ 0 & 1
\end{smallmatrix}\!\big)$.
% \sai{need to establish quite a bit of notation here}
\begin{theorem}\label{thm: miyake thm 4.6.16}
    We have
    \begin{enumerate}[label=\textup{(\arabic*)}]
        \item By the action of $\eta_q$, we have the isomorphisms \[\mathcal S_k^0(N,\chi)\cong \mathcal S_k^0(N,\chi^\prime_q\overline{\chi_q}),\quad \mathcal S_k^1(N,\chi)\cong \mathcal S_k^1(N,\chi^\prime_q\overline{\chi_q}).\]
        \item By the action of $\eta^\prime_q$, we have the isomorphisms \[\mathcal S_k^0(N,\chi)\cong \mathcal S_k^0(N,\overline{\chi^\prime_q}\chi_q),\quad \mathcal S_k^1(N,\chi)\cong \mathcal S_k^1(N,\overline{\chi^\prime_q}\chi_q).\]
        \item For $f\in\mathcal S_k(N,\chi)$, we have $f|_k[\eta^2_q] = \chi_q(-1)\overline{\chi^\prime_q}(N_q)f$, $f|_k[{\eta^\prime_q}^2] = \chi^\prime_q(-1)\overline{\chi_q}(N/N_q)f$, and $f|_k[\eta_q\eta^\prime_q] = \overline{\chi^\prime_q}(N_q)f|_k[\omega_N]$.
        \item Let $f\in\mathcal S_k^0(N,\chi)$ be a primitive form with $f(z) = \sum_{n=1}^\infty a_ne^{2\pi i n z}$, and write \[(f|_k[\eta_q])(z) = c\sum_{n=1}^\infty b_ne^{2\pi i n z}\quad \text{\textup{(}with $b_1 = 1$\textup{)}}.\] Let $g_q(z) = \sum_{n=1}^\infty b_ne^{2\pi i n z}$.
        
        Then $g_q$ is a primitive form of $\mathcal S_k^0(N,\chi^\prime_q\overline{\chi_q})$, with \[b_p = \begin{cases}
            \overline{\chi_q}(p)a_p &\text{if $p\neq q$},\\
            \chi^\prime_q(p)\overline{a_p} & \text{if $p = q$}
        \end{cases}\] for any prime $p$.
    \end{enumerate}
\end{theorem}
See \cite[Theorem 4.6.16]{miyake} for details.

\begin{theorem}\label{thm: miyake thm 4.6.17}
    Let $f(z) = \sum_{n=1}^\infty a_ne^{2\pi i nz}$ be a primitive form of $\mathcal S_k^0(N,\chi)$, and $m$ the conductor of $\chi$. For a prime factor $q$ of $N$, denote by $N_q$ and $m_q$ the $q$-components of $N$ and $m$, respectively; that is, $N_q$ and $m_q$ are the largest powers of $p$ dividing $N$ and $m$, respectively. \begin{enumerate}[label=\textup{(\arabic*)}]
        \item If $N_q = m_q$, then $\abs{a_q} = q^{(k-1)/2}$.
        \item If $N_q = q$ and $m_q = 1$, then $a_q^2 = \chi^\prime_q(q)q^{k-2}$.
        \item Otherwise, that is, if $q^2\mid N_q$ and $N_q \neq m_q$, then $a_q = 0$.
    \end{enumerate}
\end{theorem}
See \cite[Theorem 4.6.17]{miyake} for details.

% \begin{lemma}\label{lem: miyake lem 4.5.2}
%     (TO REMOVE?)For any $\alpha\in\varDelta_0(N)$ \textup{(}respectively $\varDelta_0^\ast(N)$\textup{)}, there exists a unique pair of positive integers $(l,m)$ such that $l\mid m$, $\gcd(l,N) = 1$, and 
%     \[\varGamma_0(N)\alpha\varGamma_0(N) = \varGamma_0(N)\begin{pmatrix}
%         l & 0 \\ 0 & m
%     \end{pmatrix}\varGamma_0(N)\] \[(\text{respectively }\varGamma_0(N)\alpha\varGamma_0(N) = \varGamma_0(N)\begin{pmatrix}
%         m & 0 \\ 0 & l
%     \end{pmatrix}\varGamma_0(N)).\]
% \end{lemma}
% \begin{proof}
%     \tbd
% \end{proof}

% \sai{the following need to be reordered correctly!}
% \begin{lemma}\label{lem: miyake lem 4.5.6}
%     (TO REMOVE?)Let $p$ be a prime number, and $e$ a nonnegative integer.Then \begin{enumerate}[label=\textup{(\arabic*)}]
%         \item A set of representatives for \[\varGamma_0(N)\backslash\varGamma_0(N)\begin{pmatrix}
%             1 & 0 \\ 0 & p^e
%         \end{pmatrix}\varGamma_0(N)\] is given by the set \[\begin{cases}
%             \cbr{\big(\!\begin{smallmatrix}
%                 p^{e-f} & m \\ 0 & p^f
%             \end{smallmatrix}\!\big) : 0\leq f\leq e, 0\leq m < p^f, \gcd(m,p^f,p^{e-f}) = 1} & \text{if $p\nmid N$},\\
%             \cbr{\big(\!\begin{smallmatrix}
%                 1 & m \\ 0 & p^e
%             \end{smallmatrix}\!\big) : 0\leq m < p^e} & \text{if $p\mid N$, and}
%         \end{cases} \]
%         \item $\deg\Big(\varGamma_0(N)\backslash\varGamma_0(N)\begin{pmatrix}
%             1 & 0 \\ 0 & p^e
%         \end{pmatrix}\varGamma_0(N)\Big) = \begin{cases}
%             p^e + p^{e-1} & \text{if $p\nmid N$,}\\
%             p^e & \text{if $p\mid N$.}
%         \end{cases}$ \sai{(i have to define degree, perhaps it is the one on p.73 of miyake)}
%     \end{enumerate}
% \end{lemma}
% \begin{proof}
%     \sai{DON'T PROVE}
% \end{proof}
% \begin{lemma}\label{lem: miyake lem 4.5.7}
%     (TO REMOVE)Let $p$ be a prime number and $e$ a positive integer. Then \begin{enumerate}[label=\textup{(\arabic*)}]
%         \item $T(p)T(1,p^e) = T(1,p^{e+1}) + \begin{cases}
%             (p+1)T(p,p) & \text{if $p\nmid N$ and $e=1$,}\\
%             pT(p,p)T(1,p^{e-1}) & \text{if $p\nmid N$ and $e > 1$,}\\
%             0 & \text{if $p\mid N$.}
%         \end{cases}$
%         \item $T(p)T(p^e) = \begin{cases}
%             T(p^{e+1})+pT(p,p)T(p^{e-1}) & \text{if $p\nmid N$,}\\
%             T(p^{e+1}) & \text{if $p\mid N$.}
%         \end{cases}$
%     \end{enumerate}
% \end{lemma}
% \begin{proof}
%     \tbd
% \end{proof}
% \begin{lemma}\label{lem: miyake lem 4.5.8}
%     (TO REMOVE)\begin{enumerate}[label=\textup{(\arabic*)}]
%         \item If $\gcd(lm,l^\prime m^\prime) = 1$, then $T(l,m)T(l^\prime,m^\prime) = T(ll^\prime,mm^\prime)$.
%         \item If $\gcd(m,n) = 1$, then $T(m)T(n) = T(mn)$.
%     \end{enumerate}
% \end{lemma}
% \begin{proof}
%     \tbd
% \end{proof}


\newpage

\subsection{Strong multiplicity one}
We finally prove the strong multiplicity one property for classical modular forms. Let $g\in\mathcal S_k(M,\lambda)$ be a cusp form at level $M$ that is normalized; that is, $g$ has first Fourier coefficient equal to $1$, and is a common eigenfunction of all of the Hecke operators. If there exists a primitive form $f\in\mathcal S_k^0(N,\chi)$ at level $N$ (which need not be the same as $M$) whose Fourier coefficients, that is to say its eigenvalues, agree with the Fourier coefficients of $g$ at all indices $n$ coprime to some integer $L$, we show that $g$ is equal to the primitive form $f$ at level $N$.
\begin{theorem}[Strong multiplicity one]\label{thm: sm1 miyake thm 4.6.19}
    Let $f$ be a primitive form in $\mathcal S_k^0(N,\chi)$ and let $g$ be an element of $ \mathcal S_k(M,\lambda)$, with Fourier series expansions $f(z) = \sum_{n=1}^\infty a_ne^{2\pi i nz}$ and $g(z) = \sum_{n=1}^\infty b_ne^{2\pi i n z}$. If $b_1 = 1$, $g$ is a common eigenfunction of $\mathcal R(M)\cup \mathcal R^\ast(M)$, and $a_n = b_n$ for all $n$ coprime to some integer $L$, then $N = M$ and $f = g$.
\end{theorem}
\begin{proof}
    Without loss of generality we may take $L$ to be a common multiple of $N$ and $M$. If $p$ is a prime number coprime to $L$, we have from $T_{p^2} = T_pT_p - p^{k-1}\abr{p}$ that
    \begin{align*}
        \chi(p)p^{k-1} &= a_p^2 - a_{p^2} \\
        &= b_p^2 - b_{p^2} \\
        &= \lambda(p)p^{k-1},
    \end{align*} so that $\chi(p) = \lambda(p)$. It follows that $\chi(n) = \lambda(n)$ for all $n$ coprime to $L$.

    We show that $M\mid N$. By \cref{thm: miyake thm 4.5.16}, $L(s,f)$ and $L(s,g)$ have Euler product expansions, so that
    \begin{equation}\label{eqn: ratio varLambdas}
        \frac{\varLambda_N(s,f)}{\varLambda_M(s,g)} = \bigg(\frac{\sqrt{N}}{\sqrt{M}}\bigg)^s\prod_{p\mid L}\frac{1-b_pp^{-s} + \lambda(p)p^{k-1-2s}}{1-a_pp^{-s} + \chi(p)p^{k-1-2s}}
    \end{equation} for $\Re(s)>k/2+1$. Since the ratio in the right-hand side is a meromorphic function on the whole $s$-plane, \eqref{eqn: ratio varLambdas} holds for all $s$ on the $s$-plane except at the poles.

    On the other hand, we have
    \[\frac{\varLambda_N(s,f)}{\varLambda_M(s,g)} = \frac{\varLambda_N(k-s,f|_k[\omega_N])}{\varLambda_M(k-s,g|_k[\omega_M])}\] by \cref{cor: miyake cor 4.3.7}. Furthermore, $T_ng = b_ng$, and since $g$ is a common eigenfunction of $T^\ast_n$ by assumption, we have that $T^\ast_n g = \overline{b_n}g$ and 
    \[T^\ast_n(g|_k[\omega_M]) = \overline{b_n}(g|_k[\omega_M])\] by \cref{thm: miyake thm 4.5.5}. Then by \cref{thm: miyake thm 4.5.16} it follows that $L(s, g|_k[\omega_M])$ also has an Euler product expansion. Combined with \cref{thm: miyake thm 4.6.15}(2), we have 
    \begin{equation}\label{eqn: ratio shifted varLambdas}
        \frac{\varLambda_N(s,f)}{\varLambda_M(s,g)} = \frac{\varLambda_N(k-s,f|_k[\omega_N])}{\varLambda_M(k-s,g|_k[\omega_M])} = c\bigg(\frac{\sqrt{N}}{\sqrt{M}}\bigg)^{k-s}\prod_{p\mid L}\frac{1-\overline{b_p}p^{s-k} + \overline{\lambda}(p)p^{2s-k-1}}{1-\overline{a_p}p^{s-k} + \overline{\chi}(p)p^{2s-k-1}}
    \end{equation} for some constant $c$. By combining \eqref{eqn: ratio varLambdas} and \eqref{eqn: ratio shifted varLambdas} together, we obtain
    \[\bigg(\frac{N}{M}\bigg)^s\prod_{p\mid L}\frac{1-b_pp^{-s} + \lambda(p)p^{k-1-2s}}{1-a_pp^{-s} + \chi(p)p^{k-1-2s}} = c\bigg(\frac{\sqrt{N}}{\sqrt{M}}\bigg)^k\prod_{p\mid L}\frac{1-\overline{b_p}p^{s-k} + \overline{\lambda}(p)p^{2s-k-1}}{1-\overline{a_p}p^{s-k} + \overline{\chi}(p)p^{2s-k-1}}.\]
    Denote by $M_p$ and $N_p$ the $p$-components of $M$ and $N$ respectively; that is, $M_p$ and $N_p$ are the largest powers of $p$ dividing $M$ and $N$, respectively. Then for any prime factor $p$ of $L$, 
    \begin{equation}\label{eqn: p-comps}
        \bigg(\frac{N_p}{M_p}\bigg)^s\frac{1-b_pp^{-s} + \lambda(p)p^{k-1-2s}}{1-a_pp^{-s} + \chi(p)p^{k-1-2s}} = c_p
        % \bigg(\frac{\sqrt{N}}{\sqrt{M}}\bigg)^k what happened to this factor? absorbed into $c_p$?
        \frac{1-\overline{b_p}p^{s-k} + \overline{\lambda}(p)p^{2s-k-1}}{1-\overline{a_p}p^{s-k} + \overline{\chi}(p)p^{2s-k-1}}
    \end{equation} holds for some constant $c_p$ by \cref{lem: miyake lem 3.2.1}. Let $x=p^{-s}$, and let $u$ and $v$ be the degrees of
    \begin{align*}
        1-a_pp^{-s} + \chi(p)p^{k-1-2s} &= 1-a_px + \chi(p)p^{k-1}x^2 \text{ and}\\
        1-b_pp^{-s} + \lambda(p)p^{k-1-2s} &= 1-b_px + \lambda(p)p^{k-1}x^2
    \end{align*} as polynomials of $x$, respectively. Thus $0\leq u,v\leq 2$. Furthermore, write $M_p/N_p = p^e$, so that \eqref{eqn: p-comps} becomes \begin{equation}\label{eqn: x substitution}
        x^e\frac{1-b_px + \lambda(p)p^{k-1}x^2}{1-a_px + \chi(p)p^{k-1}x^2} = c_p\frac{1-\overline{b_p}p^{-k}x^{-1} + \overline{\lambda}(p)p^{-k-1}x^{-2}}{1-\overline{a_p}p^{-k}x^{-1} + \overline{\chi}(p)p^{-k-1}x^{-2}}.
    \end{equation} We investigate each combination of values that $u,v$ may take, and show that if a particular combination of $u,v$ implies that $e>0$, that case could not occur. Thus in the remaining cases, $e\leq 0$ so that $M_p\mid N_p$.
    \begin{enumerate}[label=(\arabic*)]
        \item If $u = v$, take $x\to \infty$ in \eqref{eqn: nuts equality} (by taking $\Re(s)\to\infty$) to deduce that $e= 0$.
        \item If $u = 1$ and $v = 0$, then $a_p\neq 0$ so that \eqref{eqn: x substitution} may be rewritten as 
        \[x^e(1-\overline{a_p}p^{-k}x^{-1})=c_p(1-a_px).\] Then $e = 1$, and multiplying the above equality by $a^p$ we may rearrange to obtain 
        \[-a_pc_p = \frac{a_px-\abs{a_p}^2p^{-k}}{a_px-1},\] from which it follows that $\abs{a_p}^2 = p^k$. This is in contradiction with \cref{thm: miyake thm 4.6.17}, so this case could not occur.
        \item If $u = 0$ and $v = 1$, then \eqref{eqn: x substitution} becomes $x^e(1-b_px) = c_p(1-\overline{b_p}p^{-k}x^{-1})$, from which it follows that $e = -1$.
        \item If $u = 2$ and $v = 0$, rewrite \eqref{eqn: x substitution} as 
        \begin{equation}\label{eqn: nuts equality}
            x^e(1-\overline{a_p}p^{-k}x^{-1} + \overline{\chi}(p)p^{-k-1}x^{-2})=c_p(1-a_px + \chi(p)p^{k-1}x^2).
        \end{equation} Then $e=2$, so that we have 
        \[\chi(p)p^{k-1}c_p = \frac{\chi(p)p^{k-1}x^2-\chi(p)\overline{a_p}p^{-1}x + \abs{\chi(p)}^2p^{-2}}{\chi(p)p^{k-1}x^2-a_px+1}.\] It follows that $\chi(p)p^{k-1}c_p = 1$ and $\abs{\chi(p)}^2 = p^{2}$; the latter is impossible.
        \item If $u = 0$ and $v = 2$, deduce that $e = -2$.
        \item If $u = 2$ and $v = 1$, rewrite \eqref{eqn: x substitution} as 
        \[x^e\frac{1-\overline{a_p}p^{-k}x^{-1} + \overline{\chi}(p)p^{-k-1}x^{-2}}{1-\overline{b_p}p^{-k}x^{-1}} = c_p\frac{1-a_px + \chi(p)p^{k-1}x^2}{1-b_px},\] so that $e = 1$, and rewrite this equation as 
        \[\frac{x^2-\overline{a_p}p^{-k}x + \overline \chi(p)p^{-k-1}}{x-\overline{b_p}p^{-k}} = c_p\frac{x^2-\overline\chi(p)a_pp^{-k+1}x + \overline\chi(p)p^{-k+1}}{\overline\chi(p)p^{-k+1}(1-b_px)}.\] The roots of the polynomials 
        \[x^2-\overline{a_p}p^{-k}x + \overline \chi(p)p^{-k-1}\quad \text{and}\quad x^2-\overline\chi(p)a_pp^{-k+1}x + \overline\chi(p)p^{-k+1}\] agree. However, since the product of the roots of these polynomials are equal to the constant terms of each polynomial, we must have $\overline \chi(p)p^{-k-1} = \overline\chi(p)p^{-k+1}$, a contradiction. 
        
        % Let $t= x^{-1}$. The absolute values of the roots of 
        % \[t^2-a_pt+\chi(p)p^{k-1}\] are $p^{(k-1)/2}$ by \sai{eqn: miyake eqn 4.5.41, which is equivalent to the Ramanujan-Petersson conjecture} and the absolute values of the roots of 
        % \[\overline{\chi}(p)p^{-k-1}t^2 -\overline{a_p}p^{-k}t + 1 = \overline{\chi}(p)p^{-k-1}(t^2-a_ppt + \chi(p)p^{k+1}) \] are $p^{(k+1)/2}$, which makes \eqref{eqn: nuts equality} and, hence this case, impossible.
        \item If $u = 1$ and $v = 2$, deduce that $e = -1$ by taking limits.
    \end{enumerate}
    In any case, it follows for any $p$ that $M_p\mid N_p$ so that $M\mid N$, and that $\chi$ is induced by $\lambda$. Then \cref{thm: miyake thm 4.6.12} implies $f(z) = g(z)$, so that $N = M$.
\end{proof}

% \subsection{Oldforms and newforms}
% We start with some results which are used to take forms from lower levels $M$ dividing $N$ to $N$.
% \begin{lemma}
%     If $M\mid N$ then $\mathcal S_k(\varGamma_1(M))\subset \mathcal S_k(\varGamma_1(N))$.
% \end{lemma}
% \begin{proof}
%     \tbd
% \end{proof}
% Let $M\mid N$ and let $d$ be a factor of $N/M$. Define $\alpha_d = \big(\!\begin{smallmatrix}
%     d & 0 \\ 0 & 1
% \end{smallmatrix}\!\big)$ and consider the map $|_k[\alpha_d]$, which is an injective map $\mathcal S_k(\varGamma_1(M))\to\mathcal S_k(\varGamma_1(N))$ by \cref{lem: weight k operator properties}. Explicitly, for $f\colon\mathcal H\to\mathbb C$ we have $(f|_k[\alpha_d])(z) = d^{k/2}f(dz)$.

% We consider the subspace of $\mathcal S_k(\varGamma_1(N))$ containing all cusp forms which may be obtained from cusp forms at lower levels via a combination of the two embeddings discussed earlier.

% \begin{definition}
%     For each divisor $d$ of $N$, define 
%     \[i_d\colon (\mathcal S_k(\varGamma_1(Nd^{-1})))^2\to\mathcal S_k(\varGamma_1(N))\] by $(f,g)\mapsto f + g|_k[\alpha_d]$. The subpace of \boldindex{oldforms at level $N$} is 
%     \[\mathcal S_k(\varGamma_1(N))^\old = \sum_{\substack{p\mid N \\ p\text{ prime}}}i_p((\mathcal S_k(\varGamma_1(Np^{-1})))^2),\] and the subspace of \boldindex{newforms at level $N$} is the orthogonal complement of the subspace of oldforms with respect to the Petersson inner product,
%     \[\mathcal S_k(\varGamma_1(N))^\new = (\mathcal S_k(\varGamma_1(N))^\old)^\perp.\qedhere\]
% \end{definition}

% \begin{proposition}
%     The subspace of oldforms and the subspace of newforms in $\mathcal S_k(\varGamma_1(N))$ are stable under the Hecke operators $\abr{n},T_n$ for all $n\in\mathbb Z^+$.
% \end{proposition}
% \begin{proof}
%     \tbd
% \end{proof}
% \begin{corollary}
%     The subspace of oldforms and the subspace of newforms each have orthogonal bases of eigenforms for the Hecke operators away from the level $N$; that is, for $\cbr{\abr{n},T_n : \gcd(n,N) = 1}$.
% \end{corollary} The condition that $\gcd(n,N) = 1$ can be removed for the space of newforms.

% \sai{The ``Main Lemma'' will appear soon.}

% \newpage\section{}

\newpage\pagestyle{frontmatter}
\printindex\newpage
\begin{bibdiv}
\begin{biblist}

\bib{df}{book}{
    title = {Abstract Algebra},
    author = {Dummit, David S.},
    author = {Foote, Richard M.},
    isbn = {978-0-471-43334-7},
    year = {2004},
    publisher = {John Wiley \& Sons, Inc.}
}

\bib{diamond}{book}{
    title = {A First Course in Modular Forms},
    author = {Diamond, Fred},
    author = {Shurman, Jerry},
    isbn = {978-0-387-23229-4},
    series = {Graduate Texts in Mathematics},
    year = {2005},
    publisher = {Springer New York, NY}
}

\bib{im}{article}{
    title = {Modular forms and modular curves},
    author = {Diamond, Fred},
    author = {Im, John},
    isbn = {0-8218-0313-1},
    journal = {CMS Conf. Proc.},
    volume = {17},
    year = {1995},
    pages = {39-133},
    publisher = {American Mathematical Society, Providence, RI},
    review = {\MR{1357209}}
}

\bib{miyake}{book}{
    title = {Modular Forms},
    author = {Miyake, Toshitsune},
    isbn = {978-3-540-29592-1},
    series = {Springer Monographs in Mathematics},
    year = {2005},
    publisher = {Springer Berlin, Heidelberg}
}

\bib{serre}{book}{
    title = {A Course in Arithmetic},
    author = {Serre, Jean-Pierre},
    isbn = {978-0-387-90040-7},
    series = {Graduate Texts in Mathematics},
    year = {1978},
    publisher = {Springer New York, NY}
}

% \bib{jeremy citation needed}{book}{
%     title = {slides},
%     author = {jeremy},
% }

\end{biblist}
\end{bibdiv}
\end{document}